% Generated by scripts/build_readers.py; change the reviewed sources and rebuild.
\def\EditionMode{1}
\def\EditionLabel{Annotated master}
\def\EditionStatus{WORKING EDITION: 39 of 706 source pages reviewed and included}
\documentclass[10pt,oneside]{book}
\usepackage{fontspec}
\setmainfont{Latin Modern Roman}
\setsansfont{Latin Modern Sans}
\usepackage{amsmath,amssymb,mathtools}
\usepackage{unicode-math}
\setmathfont{Latin Modern Math}
\usepackage{xcolor,tikz-cd,lua-ul,marginnote,pdflscape}
\usetikzlibrary{arrows.meta,decorations.pathmorphing,calc,shapes.geometric}
\tikzcdset{ampersand replacement=\&}
\tikzcdset{every label/.append style={text=black}}
\usepackage[a4paper,left=18mm,right=45mm,top=18mm,bottom=22mm,marginparwidth=28mm,marginparsep=7mm]{geometry}
\usepackage[unicode,colorlinks=true,urlcolor=black,linkcolor=black]{hyperref}
\usepackage{bookmark}
\setlength{\parindent}{0pt}
\setlength{\parskip}{4pt plus 1pt minus 1pt}
\setlength{\emergencystretch}{1.5em}
\raggedbottom
\definecolor{PSHand}{RGB}{25,58,132}
\definecolor{PSDelete}{RGB}{146,37,37}
\definecolor{PSTypedOld}{RGB}{95,95,95}
\ifnum\EditionMode=2\colorlet{blue}{black}\else\colorlet{blue}{PSHand}\fi
\RenewDocumentCommand{\textcircled}{m}{\tikz[baseline=(PSCircle.base)]\node[draw,circle,inner sep=.5pt,line width=.4pt](PSCircle){#1};}
\newif\ifPSInApparatus
\newcounter{psnote}
\newcounter{psvariant}
\newcounter{pslinkmark}
\newif\ifPSInLink
\newsavebox{\PSNoteBox}
\newdimen\PSNoteLimit
{\footnotesize\global\PSNoteLimit=3\baselineskip}
\newcommand{\CurrentScanURL}{}
\newcommand{\CurrentKey}{}
\newcommand{\CurrentScan}{}
\newcommand{\pssource}[3]{\psflushmargins\clearpage\gdef\CurrentKey{#1}\gdef\CurrentScan{#2}\gdef\CurrentScanURL{#3}\setcounter{psvariant}{0}\typeout{PS-SOURCE: #1}}
\newcommand{\psdivision}[2]{\pdfbookmark[0]{#1}{#2}\par\noindent{\large\bfseries #1\par}\medskip}
\newcommand{\pslinkmarker}[2]{\stepcounter{pslinkmark}\latelua{texio.write_nl('log','PS-LINK-SHIP|#1|\CurrentKey|\arabic{pslinkmark}|#2')}}
\newcommand{\pslink}[2]{\ifPSInLink\pslinkmarker{#1}{nested}#2\else\begingroup\PSInLinktrue\href{\CurrentScanURL}{\pslinkmarker{#1}{outer}#2}\endgroup\fi}
\newcommand{\scanpage}[1]{\ifnum\EditionMode=2\leavevmode\marginpar{\scriptsize\sffamily\pslink{scanpage}{scan #1}}\else\par\noindent{\small\sffamily\pslink{scanpage}{[scan #1]}}\par\medskip\fi}
\newcommand{\mspage}[1]{\ifnum\EditionMode=2\leavevmode\marginpar{\scriptsize\sffamily p.\,#1}\else{\small\sffamily[p.\,#1]}\fi}
\newcommand{\hand}[1]{{\ifnum\EditionMode=2\color{black}\else\color{PSHand}\fi #1}}
\let\over\relax
\NewDocumentCommand{\over}{O{} m m}{{\textcolor{PSTypedOld}{#2}\ensuremath{^{\!\nearrow}}\if\relax\detokenize{#1}\relax\hand{#3}\else #3\fi}}
\NewDocumentCommand{\inserted}{O{} m}{{\if\relax\detokenize{#1}\relax\hand{\ensuremath{\langle}#2\ensuremath{\rangle}}\else\ensuremath{\langle}#2\ensuremath{\rangle}\fi}}
\newcommand{\deleted}[1]{{\color{PSDelete}\ifmmode\text{[del.: }#1\text{]}\else[\textsf{del.:} #1]\fi}}
\newcommand{\xout}[1]{{\color{PSTypedOld}\ifmmode\text{[x: }#1\text{]}\else[\textsf{x:} #1]\fi}}
\let\PSMathUnderline\underline
\RenewDocumentCommand{\underline}{+m}{\ifmmode\PSMathUnderline{#1}\else\underLine{#1}\fi}
\NewDocumentCommand{\uline}{+m}{\underline{#1}}
\NewDocumentCommand{\uuline}{+m}{\ifmmode\PSMathUnderline{\PSMathUnderline{#1}}\else\underLine{\underLine{#1}}\fi}
\newunderlinetype\PSWaveLine{\leaders\hbox{\lower.7ex\hbox{\tikz[x=1pt,y=1pt,baseline=0pt]{\draw[line width=.3pt] (0,0) .. controls (1,1) and (2,1) .. (3,0) .. controls (4,-1) and (5,-1) .. (6,0);}}}}
\NewDocumentCommand{\uwave}{+m}{{\PSWaveLine #1}}
\newcommand{\unsure}[1]{\pslink{unsure}{[#1\ifmmode\,?\else\,?\fi]}}
\newcommand{\illegible}[1]{\ifmmode\text{\pslink{illegible}{[illeg.\ #1]}}\else\pslink{illegible}{[illeg.\ #1]}\fi}
\NewDocumentCommand{\diagram}{+m}{\ifmmode\vcenter{\hbox{#1}}\else\par\medskip{\centering #1\par}\medskip\psflushmargins\fi}
\newenvironment{diagramblock}{\par\begin{center}}{\end{center}\par}
\newcommand{\PSMathNotes}{}
\makeatletter
\NewDocumentCommand{\psplacednote}{m +m}{%
      \sbox{\PSNoteBox}{\parbox{\marginparwidth}{\footnotesize\raggedright\hand{\textsf{m#1.} #2}}}%
      \ifdim\dimexpr\ht\PSNoteBox+\dp\PSNoteBox\relax<\PSNoteLimit
        \leavevmode\marginpar{\usebox{\PSNoteBox}}%
      \else
        \par\nobreak\smallskip\hrule\nobreak\smallskip
        {\footnotesize\textsf{Marginal note m#1.}\quad\hand{#2}\par}
        \nobreak\smallskip\hrule\smallskip
      \fi}
\newcommand{\psflushmargins}{\ifmmode\else\ifinner\else\ifx\PSMathNotes\@empty\else\let\PSTempNotes\PSMathNotes\global\let\PSMathNotes\@empty\PSTempNotes\fi\fi\fi}
\NewDocumentCommand{\margin}{+m}{%
  \ifPSInApparatus\ifmmode\hand{\text{[marg.: }#1\text{]}}\else\hand{\textsf{[marg.: }#1\textsf{]}}\fi
  \else\stepcounter{psnote}%
    \ifmmode
      \text{\textsuperscript{m\arabic{psnote}}}%
      \begingroup\edef\PSQueue{\endgroup\noexpand\g@addto@macro\noexpand\PSMathNotes{\noexpand\psplacednote{\arabic{psnote}}{\unexpanded{$#1$}}}}\PSQueue
    \else
      \textsuperscript{m\arabic{psnote}}\psplacednote{\arabic{psnote}}{#1}%
    \fi
  \fi}
\AtBeginDocument{\everymath\expandafter{\the\everymath\aftergroup\psflushmargins}\everydisplay\expandafter{\the\everydisplay\aftergroup\psflushmargins}}
\makeatother
\newcommand{\psnumber}[2]{\expandafter\gdef\csname psnum@#1\endcsname{#2}}
\NewDocumentCommand{\var}{+m +m m m}{\ifnum\EditionMode=0 #1\else\ifnum\EditionMode=1\ifmmode\text{\textsuperscript{\tiny\csname psnum@#4\endcsname}}\else\textsuperscript{\tiny\csname psnum@#4\endcsname}\fi\fi #2\fi}
\newcommand{\psapparatusstart}{\par\bigskip\hrule\smallskip\PSInApparatustrue\begingroup\footnotesize\raggedright\textsf{Apparatus for scan \CurrentScan}\par}
\newcommand{\psapparatusend}{\par\endgroup\PSInApparatusfalse}
\NewDocumentCommand{\psentry}{m m m +m +m m}{\par\noindent\textbf{#2. #3.}\quad\textit{Typescript:} #4\quad\textit{Corrected:} \if\relax\detokenize{#5}\relax\textnormal{[omitted]}\else #5\fi\quad #6\par}
\newcommand{\psfrontmatter}{%
\hypersetup{pageanchor=false}\begin{titlepage}\raggedright
{\Large Alexander Grothendieck\par}\vspace{12mm}
{\Huge Pursuing Stacks\par}\medskip{\LARGE 1983\par}\vspace{12mm}
{\Large\EditionLabel\par}\vspace{8mm}
{\large\EditionStatus\par}\vfill
Université de Montpellier, Grothendieck archive, cotes 134-2--134-8.\par
Including the 26 source pages of R. Brown's translation at cote 134-2, sites 115--143.\par
\end{titlepage}\hypersetup{pageanchor=true}
\chapter*{Editorial statement}
AI-transcribed by Claude Opus 5.5 (Anthropic, effort high); checked against the scans, and this edition built, by GPT-6 Astra (OpenAI, ChatGPT Pro mode); prompts and records by Claude (Anthropic).

The witness consulted and not reproduced is the LaTeX edition of Mateo Carmona with the collaboration of Ulrik Buchholtz (arXiv 2111.01000 v2, 2021), extended from the 2015 edition by ``the scrivener'', CC0 1.0.

The translation in the appendix is R. Brown's, prepared with assistance from L. Breen and J-L. Loday, as his note says; the rights in it belong to its translator or his heirs.

The rights in the original work belong to its author or the author's representatives or heirs. Where we hold rights, in the AI-produced transcription, translation, corrections and typesetting, they are dedicated to the public domain under CC0 1.0. This edition is published for educational, non-profit purposes, and we claim no rights we do not hold. If you hold rights in this work and object to its publication here, please open a GitHub issue and we will take it down.

\section*{How to read this form}
\ifnum\EditionMode=0
The diplomatic form preserves typed and handwritten wording, cancellations, insertions, overwritten readings, underlining, doubts, and original page numbers. Blue marks handwriting; red labels handwritten cancellations; gray labels typewriter cancellations and underlying overwritten text. Angle brackets mark insertions; a superscript arrow separates an overwritten reading from its replacement. Typed insertions and typed overwriting remain plain. The disposition of long marginal notes is identified explicitly.
\else
The corrected text applies the author's insertions, replacements, cancellations, and recorded transpositions. It retains his English and notation, including consistent forms such as ``standart'', ``developped'', ``it's'' for ``its'', and ``A la''. An isolated slip is corrected only when specifically recorded and justified. Mathematical anomalies are preserved and discussed in the guide.
\ifnum\EditionMode=1
Blue continues to identify handwriting. A numbered call identifies each maximal editorial operation, with its full diplomatic and corrected readings and reason in the apparatus following the source page. Internal layers of a wholly replaced or deleted passage are preserved within that entry rather than duplicated as separate entries. These entries are generated from the same variants that govern the running text.
\else
This reading form omits the apparatus, cancelled material, superseded readings, insertion brackets, and color distinctions. It retains the resulting text, marginal notes, underlining, and unresolved doubts. Scan and source-page references appear in the margin.
\fi\fi

Every scan mark and remaining doubt links to its source page in the archive viewer. Doubts are editorial reading marks, not emendations. The source-page breaks identify the archival sequence; a source page may occupy more than one printed page. A partial working edition is explicitly labelled and contains only reviewed source pages.
\clearpage}

\begin{document}
\psfrontmatter
\pssource{ps_134-8_001}{134-8/1}{https://grothendieck.umontpellier.fr/web/viewer.html?file=../134-8.pdf\#page=2}
\psdivision{Cote 134-8}{division-134-8}
%% BEGIN SOURCE ps_134-8_001
% Pursuing Stacks (cotes 134-2 to 134-8): archive cote 134-8, site page 1 (PDF page 2), typescript page 555.
% typescript. Draft compared: ps7.tex:1-56.
% AI-transcribed by Claude Opus 5.5 (claude-opus-5-5, adaptive thinking, effort high), Message Batches API, batch ps4 (msgbatch_012Aoc1NuskATPsyBeKSYK7J), 2026-11-12; reviewed against both supplied scan halves; see CHECK_DECISIONS.tsv.
\scanpage{134-8/1}
\mspage{555}

22.10.

% NOTE: a white box covers the rest of the typed text of the page, from just below the date "22.10." to the foot of the page, including the right margin; nothing of the text under it can be read on this scan.
% NOTE: a small dash-like mark stands in the left margin at the left edge of the white box, about a third of the way down the page; whatever it belongs to is hidden by the box.

\margin{1}

% Visual collation: both supplied image halves inspected; see CHECK_DECISIONS_134_7_8.tsv.

%% END SOURCE ps_134-8_001
\pssource{ps_134-8_002}{134-8/2}{https://grothendieck.umontpellier.fr/web/viewer.html?file=../134-8.pdf\#page=3}
\psnumber{ps_134-8_002-v0001}{1}
\psnumber{ps_134-8_002-v0002}{2}
\psnumber{ps_134-8_002-v0003}{3}
\psnumber{ps_134-8_002-v0004}{4}
\psnumber{ps_134-8_002-v0005}{5}
\psnumber{ps_134-8_002-v0006}{6}
\psnumber{ps_134-8_002-v0007}{7}
\psnumber{ps_134-8_002-v0008}{8}
\psnumber{ps_134-8_002-v0009}{9}
\psnumber{ps_134-8_002-v0010}{10}
\psnumber{ps_134-8_002-v0011}{11}
\psnumber{ps_134-8_002-v0012}{12}
\psnumber{ps_134-8_002-v0013}{13}
\psnumber{ps_134-8_002-v0014}{14}
\psnumber{ps_134-8_002-v0015}{15}
%% BEGIN SOURCE ps_134-8_002
% Pursuing Stacks (cotes 134-2 to 134-8): archive cote 134-8, site page 2 (PDF page 3), typescript page 556.
% typescript. Draft compared: ps7.tex:43-98.
% AI-transcribed by Claude Opus 5.5 (claude-opus-5-5, adaptive thinking, effort high), Message Batches API, batch pilot (msgbatch_0189Nhf1xB615YWnztubtLfc), 2026-11-12; reviewed against both supplied scan halves; see CHECK_DECISIONS.tsv.
\scanpage{134-8/2}
\mspage{556}
% NOTE: a short handwritten stroke at the top left, below the page number, runs into the edge of a white box.
% NOTE: a white box hides most of the upper half of the page, from just below the page number down to the paragraph "This brings me to..."; the text under it is not visible.

This brings me to somewhat more mundane matters - such as the
beginning of school and teaching duties. This year I am \var{\xout{going to be}}{}{typewriter deletion}{ps_134-8_002-v0001} in
charge of preparing the ``concours d'agrégation'', something which,\var{\xout{I}}{}{typewriter deletion}{ps_134-8_002-v0002} I
was affraid, would be rather dull\var{\over{,}{.}}{\hand{.}}{writing over}{ps_134-8_002-v0003} \var{\deleted{whereas}}{}{deletion}{ps_134-8_002-v0004} \var{\inserted{Rather}}{\hand{Rather}}{insertion}{ps_134-8_002-v0005} surprisingly the first
session of common work on a ``problème d'agrègue'' was'nt dull at all -
the problem looked interesting\hand{,}and \var{\xout{the th}}{}{typewriter deletion}{ps_134-8_002-v0006} two of the three students who
turned up so far were interested indeed\hand{,} \var{\inserted{and the atmosphere relaxed and friendly.}}{\hand{and the atmosphere relaxed and friendly.}}{insertion}{ps_134-8_002-v0007} It looks as though I was going
again to learn some mathematics through teaching, or at any rate \var{\xout{get}}{}{typewriter deletion}{ps_134-8_002-v0008}
\var{\xout{come to}}{}{typewriter deletion}{ps_134-8_002-v0009} to apply things the way I have known and understood them (in a
somewhat``highbrow'' way maybe) to the more down-to-earth vision going
with a particular curriculum (here the ``programme d'agrègue'' ) and cor-%
responding virtuosity tests - something rather far\hand{,} of course\hand{,}from my
own relation to mathematical substance ! Besides these arpegios, \var{\over[typed]{w}{W}}{W}{typewriter overwriting}{ps_134-8_002-v0010}e
started a microseminar\unsure{\hand{,}} with t\var{\over[typed]{wo}{hr}}{hr}{typewriter overwriting}{ps_134-8_002-v0011}ee participants besides me, on the
Teichmüller groupoids; I expect one of the participants is going to ta-%
ke a really active interest in the stuff. I have been feeling somewhat
reluctant to \var{\xout{get involved in}}{}{typewriter deletion}{ps_134-8_002-v0012}\var{\inserted[typed]{start}}{start}{typed insertion}{ps_134-8_002-v0013} this seminar while still involved with the
homotopy story, which is going to keep me busy easily for the next
six months still, maybe even longer. Sure enough, the little I told
us about some of the structures and operations to be investigated,
while progressively gaining view of them again after a long oblivion
- even as discovering them again hesitatingly, while pulling them out
of a mist by bribes and bits - this little was enough to revive the
\var{\xout{great}}{}{typewriter deletion}{ps_134-8_002-v0014}\var{\inserted[typed]{special}}{special}{typed insertion}{ps_134-8_002-v0015} fascination of these structures, and all that goes with them.

\margin{2}

% Visual collation: both supplied image halves inspected; see CHECK_DECISIONS_134_7_8.tsv.

%% END SOURCE ps_134-8_002
\ifnum\EditionMode=1
\psapparatusstart
\psentry{ps_134-8_002-v0001}{1}{typewriter deletion}{\xout{going to be}}{}{The typewriter strikes out this material; it leaves the corrected running text.}
\psentry{ps_134-8_002-v0002}{2}{typewriter deletion}{\xout{I}}{}{The typewriter strikes out this material; it leaves the corrected running text.}
\psentry{ps_134-8_002-v0003}{3}{writing over}{\over{,}{.}}{\hand{.}}{The superposed handwritten reading replaces the underlying reading.}
\psentry{ps_134-8_002-v0004}{4}{deletion}{\deleted{whereas}}{}{The author cancels this material by hand; it leaves the corrected running text.}
\psentry{ps_134-8_002-v0005}{5}{insertion}{\inserted{Rather}}{\hand{Rather}}{The handwritten insertion is incorporated at its indicated position; insertion brackets leave the running text.}
\psentry{ps_134-8_002-v0006}{6}{typewriter deletion}{\xout{the th}}{}{The typewriter strikes out this material; it leaves the corrected running text.}
\psentry{ps_134-8_002-v0007}{7}{insertion}{\inserted{and the atmosphere relaxed and friendly.}}{\hand{and the atmosphere relaxed and friendly.}}{The handwritten insertion is incorporated at its indicated position; insertion brackets leave the running text.}
\psentry{ps_134-8_002-v0008}{8}{typewriter deletion}{\xout{get}}{}{The typewriter strikes out this material; it leaves the corrected running text.}
\psentry{ps_134-8_002-v0009}{9}{typewriter deletion}{\xout{come to}}{}{The typewriter strikes out this material; it leaves the corrected running text.}
\psentry{ps_134-8_002-v0010}{10}{typewriter overwriting}{\over[typed]{w}{W}}{W}{The superposed typed reading replaces the underlying reading.}
\psentry{ps_134-8_002-v0011}{11}{typewriter overwriting}{\over[typed]{wo}{hr}}{hr}{The superposed typed reading replaces the underlying reading.}
\psentry{ps_134-8_002-v0012}{12}{typewriter deletion}{\xout{get involved in}}{}{The typewriter strikes out this material; it leaves the corrected running text.}
\psentry{ps_134-8_002-v0013}{13}{typed insertion}{\inserted[typed]{start}}{start}{The typed insertion is incorporated at its indicated position; insertion brackets leave the running text.}
\psentry{ps_134-8_002-v0014}{14}{typewriter deletion}{\xout{great}}{}{The typewriter strikes out this material; it leaves the corrected running text.}
\psentry{ps_134-8_002-v0015}{15}{typed insertion}{\inserted[typed]{special}}{special}{The typed insertion is incorporated at its indicated position; insertion brackets leave the running text.}
\psapparatusend
\fi
\pssource{ps_134-8_003}{134-8/3}{https://grothendieck.umontpellier.fr/web/viewer.html?file=../134-8.pdf\#page=4}
\psnumber{ps_134-8_003-v0001}{1}
\psnumber{ps_134-8_003-v0002}{2}
\psnumber{ps_134-8_003-v0003}{3}
\psnumber{ps_134-8_003-v0004}{4}
\psnumber{ps_134-8_003-v0005}{5}
\psnumber{ps_134-8_003-v0006}{6}
\psnumber{ps_134-8_003-v0007}{7}
\psnumber{ps_134-8_003-v0008}{8}
\psnumber{ps_134-8_003-v0009}{9}
\psnumber{ps_134-8_003-v0010}{10}
\psnumber{ps_134-8_003-v0011}{11}
\psnumber{ps_134-8_003-v0012}{12}
\psnumber{ps_134-8_003-v0013}{13}
\psnumber{ps_134-8_003-v0014}{14}
\psnumber{ps_134-8_003-v0015}{15}
\psnumber{ps_134-8_003-v0016}{16}
\psnumber{ps_134-8_003-v0017}{17}
%% BEGIN SOURCE ps_134-8_003
% Pursuing Stacks (cotes 134-2 to 134-8): archive cote 134-8, site page 3 (PDF page 4), typescript page 557.
% typescript. Draft compared: ps7.tex:86-153.
% AI-transcribed by Claude Opus 5.5 (claude-opus-5-5, adaptive thinking, effort high), Message Batches API, batch ps4 (msgbatch_012Aoc1NuskATPsyBeKSYK7J), 2026-11-12; reviewed against both supplied scan halves; see CHECK_DECISIONS.tsv.
\scanpage{134-8/3}
\mspage{55\var{\over{\unsure{6}}{7}}{\hand{7}}{writing over}{ps_134-8_003-v0001}}

It seems to me there has'nt been a sin\var{\over[typed]{\unsure{g}}{g}}{g}{typewriter overwriting}{ps_134-8_003-v0002}le thing in mathematics,  \uline{inclu-}%
\uline{ding motives},  which has exerted such a fascination upon me.  It will be \margin{it's very close anyhow}
hard I am affraid to carry on a seminar on such a stuff,  and go on and
carry to their (hopefully happy!) end those ponderings and notes on
homotopical algebra so-called  -  and a somewhat crazy one too at times,
I am affraid!  We'll see what \var{\over{\unsure{c}}{c}}{\hand{c}}{writing over}{ps_134-8_003-v0003}omes out of all this !  T$\overset{\illegible{1}}{\underset{\unsure{n}}{h}}$e night after
the seminar session at any rate,  and already during the two hours dri-%
\var{\xout{At any rate I decide}}{}{typewriter deletion}{ps_134-8_003-v0004}
ve home,  \var{\xout{the head}}{}{typewriter deletion}{ps_134-8_003-v0005} \var{\deleted{the stuff}}{}{deletion}{ps_134-8_003-v0006} \var{\inserted{the Teichmüller stuff}}{\hand{the Teichmüller stuff}}{insertion}{ps_134-8_003-v0007} was brewing anew in my head.  Maybe it
would have gone on for days,  but \var{\deleted{next}}{}{deletion}{ps_134-8_003-v0008} \var{\inserted{the very}}{\hand{the very}}{insertion}{ps_134-8_003-v0009} day already several things came \margin{$\sim$}
up demanding special attention,
% NOTE: a white box hides the rest of this line and the beginning of the next line; the next line is hidden as far as the box reaches

26.8.

134.  It is time to come back to the ``review'' of linearization of
(ordinary) homotopy type\hand{s},  and the homology and cohomology formalism
in the context of the modelizer (Cat),  which has been pushed aside now
for over two months,  for the sake of that endless digression on schema-%
tization of homotopy types.\hand{\unsure{*}}  The strong tie between these two strains
of reflection has been the equal importance in both of the linearization
process.  Technically speaking though,  linearization\hand{,} as finally handled
in the previous chapter (via the so-called ``integ\var{\over[typed]{g}{r}}{r}{typewriter overwriting}{ps_134-8_003-v0010}ators'')\hand{,} looks pretty
much different from the similar operation in the schematic set-up,  due
\hand{(}at least partly\hand{)} to the choice we made of models for expressing schematic
homotopy types,  namely taking semisimplicial scheme-like objects\var{\over{,}{;}}{\hand{;}}{writing over}{ps_134-8_003-v0011}
which means,  notably,  working with \hand{$\Delta$} rather than more general test
categories,  and relying heavily on the Kan-Dold-Puppe relationship.
In the disc\var{\over[typed]{\unsure{er}}{re}}{re}{typewriter overwriting}{ps_134-8_003-v0012}te set-up,  it turn\var{\over[typed]{\unsure{o}}{e}}{e}{typewriter overwriting}{ps_134-8_003-v0013}d out (somewhat inexpectedly) that the
latter can be replaced,  when working with more general categories \var{\inserted[typed]{A}}{A}{typed insertion}{ps_134-8_003-v0014} than
\hand{$\Delta$} (which need not even be test categories or anythin\var{\over[typed]{\unsure{g}}{g}}{g}{typewriter overwriting}{ps_134-8_003-v0015} of the kind),  by
the L\hand{$p$}$^{ab}_{!}$ ope\var{\over[typed]{\unsure{a}}{r}}{r}{typewriter overwriting}{ps_134-8_003-v0016}ation (when \hand{$p$} : A $\hand{\longrightarrow}$ e is the map from A to the final ob-\margin{\textcircled{p}}%
ject e in (Ca\var{\over[typed]{y}{t}}{t}{typewriter overwriting}{ps_134-8_003-v0017})),  computable  in terms of a choice of an ``integrator''
for A.  More specifically,  recalling that
\[ (1) \qquad p^{ab}_{!} : A^{\hand{\wedge}}_{ab} \hand{\longrightarrow} (Ab) \]
is defined as the left adjoint of the  inverse image functor
\[ p^{*}: \ (Ab) \hand{\longrightarrow} A^{\hand{\wedge}}_{ab} \quad , \]
associating to every abelian group the corresponding ``constant'' abelian
\margin{3}

% Visual collation: both supplied image halves inspected; see CHECK_DECISIONS_134_7_8.tsv.

%% END SOURCE ps_134-8_003
\ifnum\EditionMode=1
\psapparatusstart
\psentry{ps_134-8_003-v0001}{1}{writing over}{\over{\unsure{6}}{7}}{\hand{7}}{The superposed handwritten reading replaces the underlying reading.}
\psentry{ps_134-8_003-v0002}{2}{typewriter overwriting}{\over[typed]{\unsure{g}}{g}}{g}{The superposed typed reading replaces the underlying reading.}
\psentry{ps_134-8_003-v0003}{3}{writing over}{\over{\unsure{c}}{c}}{\hand{c}}{The superposed handwritten reading replaces the underlying reading.}
\psentry{ps_134-8_003-v0004}{4}{typewriter deletion}{\xout{At any rate I decide}}{}{The typewriter strikes out this material; it leaves the corrected running text.}
\psentry{ps_134-8_003-v0005}{5}{typewriter deletion}{\xout{the head}}{}{The typewriter strikes out this material; it leaves the corrected running text.}
\psentry{ps_134-8_003-v0006}{6}{deletion}{\deleted{the stuff}}{}{The author cancels this material by hand; it leaves the corrected running text.}
\psentry{ps_134-8_003-v0007}{7}{insertion}{\inserted{the Teichmüller stuff}}{\hand{the Teichmüller stuff}}{The handwritten insertion is incorporated at its indicated position; insertion brackets leave the running text.}
\psentry{ps_134-8_003-v0008}{8}{deletion}{\deleted{next}}{}{The author cancels this material by hand; it leaves the corrected running text.}
\psentry{ps_134-8_003-v0009}{9}{insertion}{\inserted{the very}}{\hand{the very}}{The handwritten insertion is incorporated at its indicated position; insertion brackets leave the running text.}
\psentry{ps_134-8_003-v0010}{10}{typewriter overwriting}{\over[typed]{g}{r}}{r}{The superposed typed reading replaces the underlying reading.}
\psentry{ps_134-8_003-v0011}{11}{writing over}{\over{,}{;}}{\hand{;}}{The superposed handwritten reading replaces the underlying reading.}
\psentry{ps_134-8_003-v0012}{12}{typewriter overwriting}{\over[typed]{\unsure{er}}{re}}{re}{The superposed typed reading replaces the underlying reading.}
\psentry{ps_134-8_003-v0013}{13}{typewriter overwriting}{\over[typed]{\unsure{o}}{e}}{e}{The superposed typed reading replaces the underlying reading.}
\psentry{ps_134-8_003-v0014}{14}{typed insertion}{\inserted[typed]{A}}{A}{The typed insertion is incorporated at its indicated position; insertion brackets leave the running text.}
\psentry{ps_134-8_003-v0015}{15}{typewriter overwriting}{\over[typed]{\unsure{g}}{g}}{g}{The superposed typed reading replaces the underlying reading.}
\psentry{ps_134-8_003-v0016}{16}{typewriter overwriting}{\over[typed]{\unsure{a}}{r}}{r}{The superposed typed reading replaces the underlying reading.}
\psentry{ps_134-8_003-v0017}{17}{typewriter overwriting}{\over[typed]{y}{t}}{t}{The superposed typed reading replaces the underlying reading.}
\psapparatusend
\fi
\pssource{ps_134-8_004}{134-8/4}{https://grothendieck.umontpellier.fr/web/viewer.html?file=../134-8.pdf\#page=5}
\psnumber{ps_134-8_004-v0001}{1}
\psnumber{ps_134-8_004-v0002}{2}
\psnumber{ps_134-8_004-v0003}{3}
\psnumber{ps_134-8_004-v0004}{4}
\psnumber{ps_134-8_004-v0005}{5}
\psnumber{ps_134-8_004-v0006}{6}
\psnumber{ps_134-8_004-v0007}{7}
\psnumber{ps_134-8_004-v0008}{8}
\psnumber{ps_134-8_004-v0009}{9}
\psnumber{ps_134-8_004-v0010}{10}
\psnumber{ps_134-8_004-v0011}{11}
%% BEGIN SOURCE ps_134-8_004
% Pursuing Stacks (cotes 134-2 to 134-8): archive cote 134-8, site page 4 (PDF page 5), typescript page 558.
% typescript. Draft compared: ps7.tex:133-208.
% AI-transcribed by Claude Opus 5.5 (claude-opus-5-5, adaptive thinking, effort high), Message Batches API, batch ps4 (msgbatch_012Aoc1NuskATPsyBeKSYK7J), 2026-11-12; reviewed against both supplied scan halves; see CHECK_DECISIONS.tsv.
\scanpage{134-8/4}
\mspage{558}

presheaf on A,  and
\[ \text{(2)}\qquad \mathrm{LH}. = \mathrm{Lp}_{!}^{ab} : D^{-}(A^{\hand{\wedge}}_{ab}) \hand{\longrightarrow} D^{-}(Ab) \hand{,} \quad \text{inducing}\ \margin{][} D.(A^{\hand{\wedge}}_{ab}) \hand{\longrightarrow} D.(Ab)\hand{,} \]
is it's left derived functor (computed using projective resolutions of
complexes bounded above in $A^{\hand{\wedge}}_{ab}$ ),  composing (2) with the tautological
inclusion functor
\[ A^{\hand{\wedge}}_{ab} \hand{\hookrightarrow} D^{\var{\deleted{\illegible{1}}}{}{deletion}{ps_134-8_004-v0001}}.(A^{\hand{\wedge}}_{ab}) \hand{\hookrightarrow} D^{-}(A^{\hand{\wedge}}_{ab}) \quad , \]
we get a canonical functor
\[ \text{(3)}\qquad \mathrm{Lp}_{!}^{ab}\big| A^{\hand{\wedge}}_{ab} : \ A^{\hand{\wedge}}_{ab} \hand{\longrightarrow} D.(Ab) \ , \]
which may be expressed,  in terms of an integrator  $L.^{B}$  for  $B = A^{o}$\hand{,}
i.e. a projective resolution of the constant presheaf $\underline{Z}_B$  in $B^{\hand{\wedge}}_{ab}$ ,  as
the composition
\[ \text{(4)}\qquad A^{\hand{\wedge}}_{ab} \hand{\longrightarrow} Ch.(Ab) \hand{\longrightarrow} D.(Ab) \quad , \]
where the first arrow is
\[ \text{(5)}\qquad F \hand{\longmapsto} F *_{\underline{Z}} L.^{B} \ : \var{\deleted{\illegible{1}}}{}{deletion}{ps_134-8_004-v0002}\ A^{\hand{\wedge}}_{ab} \hand{\longrightarrow} Ch.(Ab) \qquad , \]
and the second is the canonical localization functor.   The functors
(2) and (3) may be viewed as the (total) homology functors of A,  with
coe\var{\over{\unsure{ffi}}{ffi}}{\hand{ffi}}{writing over}{ps_134-8_004-v0003}cients in \var{\xout{Objects of}}{}{typewriter deletion}{ps_134-8_004-v0004}  complexes of abelian presheaves,  resp.
in abelian presheaves simply.  When we focus attention on the latter,
we may introduce in $A^{\hand{\wedge}}_{ab}$ the set of arrows which \var{\xout{are transformed}}{}{typewriter deletion}{ps_134-8_004-v0005} become
isomorphisms under the \var{\xout{functor}}{}{typewriter deletion}{ps_134-8_004-v0006} total homology functor (3),  let's call
them ``\uline{abelian} \uline{weak} \uline{equivalences}'' in $A^{\hand{\wedge}}_{ab}$ ,  not to be confused with the
notion of quasi-isomorphism for a map \var{\over[typed]{\unsure{p}}{b}}{b}{typewriter overwriting}{ps_134-8_004-v0007}etween complexes in $A^{\hand{\wedge}}_{ab}$ .  Let's
denote by
\[ \text{(6)}\qquad \mathrm{Hotab}_A \ = (\var{\over{\unsure{W}}{W}}{\hand{W}}{writing over}{ps_134-8_004-v0008}_A^{ab})^{-1} A^{\hand{\wedge}}_{ab} \]
the corresponding localized category of $A^{\hand{\wedge}}_{ab}$ ,  where  $W_A^{ab}$  denotes the
set of abelian weak equivalences in $A^{\hand{\wedge}}_{ab}$ .   Thus the choice of an inte-%
grator  $L.^{B\var{\deleted{\illegible{1}}}{}{deletion}{ps_134-8_004-v0009}}$  for A (i.e. a cointegrator for B) gives rise to a commutative
diagram of functors
\diagram{
% Hand-drawn arrows are marked blue; typed objects remain plain.
\text{(7)}\qquad
\begin{tikzcd}
A^{\hand{\wedge}}_{ab} \arrow[color=blue,r] \arrow[color=blue,d] \& Ch.(Ab) \arrow[color=blue,d] \\
\mathrm{Hotab}_A \arrow[color=blue,r] \& D.(Ab) \overset{def}{=} \mathrm{Hotab}
\end{tikzcd}
\qquad ,
}
where the \var{\xout{second}}{}{typewriter deletion}{ps_134-8_004-v0010} lower horizontal arrow is defined via (3) \hand{(}independently
of the choice of $L.^{B}$\hand{)},  the vertical arrows being the localization func-%
tors.   Beware that even when  $A=\hand{\Delta}$  and $L.^{B}$  is the usual, ``standar\var{\over{t}{d}}{\hand{d}}{writing over}{ps_134-8_004-v0011}''
\margin{\textcircled{$L.^B$}}
\margin{4}

% Visual collation: both supplied image halves inspected; see CHECK_DECISIONS_134_7_8.tsv.

%% END SOURCE ps_134-8_004
\ifnum\EditionMode=1
\psapparatusstart
\psentry{ps_134-8_004-v0001}{1}{deletion}{$\deleted{\illegible{1}}$}{}{The author cancels this material by hand; it leaves the corrected running text.}
\psentry{ps_134-8_004-v0002}{2}{deletion}{$\deleted{\illegible{1}}$}{}{The author cancels this material by hand; it leaves the corrected running text.}
\psentry{ps_134-8_004-v0003}{3}{writing over}{\over{\unsure{ffi}}{ffi}}{\hand{ffi}}{The superposed handwritten reading replaces the underlying reading.}
\psentry{ps_134-8_004-v0004}{4}{typewriter deletion}{\xout{Objects of}}{}{The typewriter strikes out this material; it leaves the corrected running text.}
\psentry{ps_134-8_004-v0005}{5}{typewriter deletion}{\xout{are transformed}}{}{The typewriter strikes out this material; it leaves the corrected running text.}
\psentry{ps_134-8_004-v0006}{6}{typewriter deletion}{\xout{functor}}{}{The typewriter strikes out this material; it leaves the corrected running text.}
\psentry{ps_134-8_004-v0007}{7}{typewriter overwriting}{\over[typed]{\unsure{p}}{b}}{b}{The superposed typed reading replaces the underlying reading.}
\psentry{ps_134-8_004-v0008}{8}{writing over}{$\over{\unsure{W}}{W}$}{$\hand{W}$}{The superposed handwritten reading replaces the underlying reading.}
\psentry{ps_134-8_004-v0009}{9}{deletion}{$\deleted{\illegible{1}}$}{}{The author cancels this material by hand; it leaves the corrected running text.}
\psentry{ps_134-8_004-v0010}{10}{typewriter deletion}{\xout{second}}{}{The typewriter strikes out this material; it leaves the corrected running text.}
\psentry{ps_134-8_004-v0011}{11}{writing over}{\over{t}{d}}{\hand{d}}{The superposed handwritten reading replaces the underlying reading.}
\psapparatusend
\fi
\pssource{ps_134-8_005}{134-8/5}{https://grothendieck.umontpellier.fr/web/viewer.html?file=../134-8.pdf\#page=6}
\psnumber{ps_134-8_005-v0001}{1}
\psnumber{ps_134-8_005-v0002}{2}
\psnumber{ps_134-8_005-v0003}{3}
\psnumber{ps_134-8_005-v0004}{4}
\psnumber{ps_134-8_005-v0005}{5}
\psnumber{ps_134-8_005-v0006}{6}
\psnumber{ps_134-8_005-v0007}{7}
\psnumber{ps_134-8_005-v0008}{8}
\psnumber{ps_134-8_005-v0009}{9}
\psnumber{ps_134-8_005-v0010}{10}
\psnumber{ps_134-8_005-v0011}{11}
\psnumber{ps_134-8_005-v0012}{12}
\psnumber{ps_134-8_005-v0013}{13}
%% BEGIN SOURCE ps_134-8_005
% Pursuing Stacks (cotes 134-2 to 134-8): archive cote 134-8, site page 5 (PDF page 6), typescript page 559.
% typescript. Draft compared: ps7.tex:195-261.
% AI-transcribed by Claude Opus 5.5 (claude-opus-5-5, adaptive thinking, effort high), Message Batches API, batch ps4 (msgbatch_012Aoc1NuskATPsyBeKSYK7J), 2026-11-12; reviewed against both supplied scan halves; see CHECK_DECISIONS.tsv.
\scanpage{134-8/5}
\mspage{559}

integrat or for \hand{$\Delta$}, the upper horizontal arrow in (7) is \uline{not} the
Kan-Dold-Puppe equivalence of categories, it has to be followed still
by the ``nor\var{\over[typed]{\illegible{1}}{ma}}{ma}{typewriter overwriting}{ps_134-8_005-v0001}lization'' operation. Thus \uline{we certainly should not expect} \margin{!x}
\uline{in any case the functor} (5) \uline{to be an equivalence} - \uline{however we suspect}
\uline{that when A is a test category} (maybe even a weak test category would
do it), \uline{then the lower horizontal arrow in} (7)
\[ (8) \qquad \mathrm{Hotab}_A \hand{\longrightarrow} \mathrm{D.(Ab)} = \mathrm{H}_{\mathrm{o}}\mathrm{tab} \]
\uline{is an equivalence of categories}. \var{\inserted[typed]{Whenever true, \xout{this}}}{Whenever true, }{typed insertion}{ps_134-8_005-v0002} \var{\xout{Such kind of statement,}}{}{typewriter deletion}{ps_134-8_005-v0003} for a given
category A , \var{\inserted[typed]{this statement}}{this statement}{typed insertion}{ps_134-8_005-v0004} looks like a reasonable substitute (on the level of the
relevant derived categories) \var{\xout{of}}{}{typewriter deletion}{ps_134-8_005-v0005} for the Kan-Dold-Puppe theorem, known
in the two cases \hand{$\Delta$} and \hand{\unsure{$\bigcirc$}}.

There are however still two \var{\inserted[typed]{important}}{important}{typed insertion}{ps_134-8_005-v0006} extra features in the case A=\hand{$\Delta$}, which
deserve to be understood in the case of more general A . One is that
a map $u\colon F\hand{\to} G$ in $A^{\hand{\wedge}}_{ab}$ is in $W^{ab}_A$ (i.e. is an ``abelian weak equivalence''\hand{)}
iff it is in $W_A$ i.e. iff it is a weak equivalence, when forgetting the
abelian structures. In terms of a final object e in A (when such ob-%
ject exists indeed), viewing the categories $A_{/F}$ (for F in $A^{\hand{\wedge}}_{ab}$) pointed \margin{it is enough that A be 1 connected instead of having a final object}
by the zero map $e \hand{\to} F$ , which yields the necessary base-point \var{\inserted{$e_F$}}{\hand{$e_F$}}{insertion}{ps_134-8_005-v0007} for
defining homotopy invariants $\hand{\pi}_i(F)$ for $i \hand{\geqslant} 0$, the relationship just
considered between abelian weak equival\unsure{ë}nce and weak equivalence, will
follow of course whenever we have functorial isomorphisms
\[ (9) \qquad \hand{\pi}_i(F) \overset{\mathrm{def}}{=} \hand{\pi}_i(A_{/F} \hand{;} e_F) \hand{\simeq} H_i(A,F) \var{\xout{\illegible{1}}}{}{typewriter deletion}{ps_134-8_005-v0008} \quad , \]
which are known to exist indeed in the case A =\hand{$\Delta$}. It should be noted
that in section 92, when starting (in a somewhat casual way) with some
reflections on ``abelianization'', we introduced a category (denoted
by $\mathrm{Hotab}_A$) by localizing $A^{\hand{\wedge}}_{ab}$ with respect to \var{\inserted{the}}{\hand{the}}{insertion}{ps_134-8_005-v0009} maps which are weak
equivalences (when forgetting the abelian structures), whereas it has
now become clear that, in case \var{\inserted[typed]{the latter}}{the latter}{typed insertion}{ps_134-8_005-v0010} \var{\xout{this notion}}{}{typewriter deletion}{ps_134-8_005-v0011} should not coincide with
the \var{\xout{notion}}{}{typewriter deletion}{ps_134-8_005-v0012} ``abelian weak equivalences'' defined in terms of linearization
or homology, it is the category (6) definitely which is the right one.
Still the question of defining isomorphisms (9) when A has a final
object, and whether the equality
\[ (10) \qquad W^{ab}_A = \var{\xout{\illegible{1}}}{}{typewriter deletion}{ps_134-8_005-v0013} \ \mathrm{forg}_A{}^{-1}(W_A) \]
holds, where
\[ \mathrm{forg}_A\colon A^{\hand{\wedge}}_{ab} \hand{\longrightarrow} A^{\hand{\wedge}} \]
is the ``forgetful functor'' , should be settled for general A.

\margin{5}

% Visual collation: both supplied image halves inspected; see CHECK_DECISIONS_134_7_8.tsv.

%% END SOURCE ps_134-8_005
\ifnum\EditionMode=1
\psapparatusstart
\psentry{ps_134-8_005-v0001}{1}{typewriter overwriting}{\over[typed]{\illegible{1}}{ma}}{ma}{The superposed typed reading replaces the underlying reading.}
\psentry{ps_134-8_005-v0002}{2}{typed insertion}{\inserted[typed]{Whenever true, \xout{this}}}{Whenever true, }{The typed insertion is incorporated at its indicated position; insertion brackets leave the running text.}
\psentry{ps_134-8_005-v0003}{3}{typewriter deletion}{\xout{Such kind of statement,}}{}{The typewriter strikes out this material; it leaves the corrected running text.}
\psentry{ps_134-8_005-v0004}{4}{typed insertion}{\inserted[typed]{this statement}}{this statement}{The typed insertion is incorporated at its indicated position; insertion brackets leave the running text.}
\psentry{ps_134-8_005-v0005}{5}{typewriter deletion}{\xout{of}}{}{The typewriter strikes out this material; it leaves the corrected running text.}
\psentry{ps_134-8_005-v0006}{6}{typed insertion}{\inserted[typed]{important}}{important}{The typed insertion is incorporated at its indicated position; insertion brackets leave the running text.}
\psentry{ps_134-8_005-v0007}{7}{insertion}{\inserted{$e_F$}}{\hand{$e_F$}}{The handwritten insertion is incorporated at its indicated position; insertion brackets leave the running text.}
\psentry{ps_134-8_005-v0008}{8}{typewriter deletion}{$\xout{\illegible{1}}$}{}{The typewriter strikes out this material; it leaves the corrected running text.}
\psentry{ps_134-8_005-v0009}{9}{insertion}{\inserted{the}}{\hand{the}}{The handwritten insertion is incorporated at its indicated position; insertion brackets leave the running text.}
\psentry{ps_134-8_005-v0010}{10}{typed insertion}{\inserted[typed]{the latter}}{the latter}{The typed insertion is incorporated at its indicated position; insertion brackets leave the running text.}
\psentry{ps_134-8_005-v0011}{11}{typewriter deletion}{\xout{this notion}}{}{The typewriter strikes out this material; it leaves the corrected running text.}
\psentry{ps_134-8_005-v0012}{12}{typewriter deletion}{\xout{notion}}{}{The typewriter strikes out this material; it leaves the corrected running text.}
\psentry{ps_134-8_005-v0013}{13}{typewriter deletion}{$\xout{\illegible{1}}$}{}{The typewriter strikes out this material; it leaves the corrected running text.}
\psapparatusend
\fi
\pssource{ps_134-8_006}{134-8/6}{https://grothendieck.umontpellier.fr/web/viewer.html?file=../134-8.pdf\#page=7}
\psnumber{ps_134-8_006-v0001}{1}
\psnumber{ps_134-8_006-v0002}{2}
\psnumber{ps_134-8_006-v0003}{3}
\psnumber{ps_134-8_006-v0004}{4}
\psnumber{ps_134-8_006-v0005}{5}
\psnumber{ps_134-8_006-v0006}{6}
\psnumber{ps_134-8_006-v0007}{7}
\psnumber{ps_134-8_006-v0008}{8}
\psnumber{ps_134-8_006-v0009}{9}
\psnumber{ps_134-8_006-v0010}{10}
\psnumber{ps_134-8_006-v0011}{11}
\psnumber{ps_134-8_006-v0012}{12}
\psnumber{ps_134-8_006-v0013}{13}
%% BEGIN SOURCE ps_134-8_006
% Pursuing Stacks (cotes 134-2 to 134-8): archive cote 134-8, site page 6 (PDF page 7), typescript page 560.
% typescript. Draft compared: ps7.tex:243-309.
% AI-transcribed by Claude Opus 5.5 (claude-opus-5-5, adaptive thinking, effort high), Message Batches API, batch ps4 (msgbatch_012Aoc1NuskATPsyBeKSYK7J), 2026-11-12; reviewed against both supplied scan halves; see CHECK_DECISIONS.tsv.
\scanpage{134-8/6}
\mspage{560}

T\var{\over[typed]{h}{h}}{h}{typewriter overwriting}{ps_134-8_006-v0001}e other "extra feature"  is about  the relationship of $W^{ab}_{A}$  with
the notion of homotopism.  In case $A = \hand{\Delta}$ ,  a map in $A^{\hand{\wedge}}_{ab}$  is a weak equi-%
valence (or equivalently,  an abelian weak equivalence)  iff it is a
homotopism when forgetting the abelian structures  -  this follows from
the well known fact that semisimplicial groups are Kan complexes.  There
is of course also a notion \var{\xout{of}}{}{typewriter deletion}{ps_134-8_006-v0002} of homotopism in the stronger abelian sen-%
se -  a particular case of the notion of a homotopism between semisimpli-%
cial objects in an arbitrary category (here in (Ab)).  The Kan-Dold-Pupp\unsure{e}
theory implies that if F and G in $\hand{\Delta}^{\hand{\wedge}}_{ab}$  have as values projective $\underline{Z}$-modu-%
les,  then \var{\over{A}{a}}{\hand{a}}{writing over}{ps_134-8_006-v0003} map  $F \hand{\longrightarrow} G$ in $\hand{\Delta}^{\hand{\wedge}}_{ab}$  is a weak equivalence iff it is an
"abelian" homotopism;  and likewise,  two maps  $u,v\colon F \hand{\rightrightarrows} G$  are equal
in $\mathrm{Hotab}_{\hand{\Delta}}$  iff they are homotopic (in the strict,\uline{abelian} sense of the
word).  T\var{\over[typed]{h}{h}}{h}{typewriter overwriting}{ps_134-8_006-v0004}e corresponding statements still make a sense and are true,
when replacing $\underline{Z}$ by any ground ring k ,  working in $\hand{\Delta}^{\hand{\wedge}}_{k}$  rather than
$\hand{\Delta}^{\hand{\wedge}}_{\underline{Z}} = \hand{\Delta}^{\hand{\wedge}}_{ab}$   -  and more generally still,  when working in
\[ \hand{\Delta}^{\hand{\wedge}}_{\hand{\mathcal{M}}} = \underline{\mathrm{Hom}}(\hand{\Delta}^{o}, \hand{\mathcal{M}}) \quad , \]
where $\hand{\mathcal{M}}$ is any abelian category.  Now \var{\deleted{\xout{for}}}{}{deletion}{ps_134-8_006-v0005} \var{\inserted[typed]{replacing \hand{$\Delta$} by}}{replacing \hand{$\Delta$} by}{typed insertion}{ps_134-8_006-v0006} an arbitrary object A in
(Cat),  these statements still make a sense, \var{\inserted[typed]{it would seem,}}{it would seem,}{typed insertion}{ps_134-8_006-v0007} provided we got on $\hand{A}^{\hand{\wedge}}$ a
suitable "homotopy structure" ,  more specifically,  a suitable "homoto-%
py interval structure" (in the sense of section 51).  In section 97
(p. 355) we reviewed the three standard homotopy interval structures
which may be introduced on a category $A^{\hand{\wedge}}$,  and their relationships -
the main impression remaining was that,  in case A is a \uline{contractor}
(cf. section 95),  those three structures coincide,  and may be defined
also in terms of a \uline{contractibility structure} on $A^{\hand{\wedge}}$ (section 51), \var{\inserted{the latter}}{\hand{the latter}}{insertion}{ps_134-8_006-v0008} giving
rise (via the general construction of section 79) to the usual notion
of weak equivalence $W_{A}$ in $A^{\hand{\wedge}}$.  Thus we may hope that the feature just
mentioned  for $\hand{\Delta}^{\hand{\wedge}}$ may be valid too for $A^{\hand{\wedge}}$,  whenever A is a contractor.

In case A is only a test category,  \var{\xout{the reasonable}}{}{typewriter deletion}{ps_134-8_006-v0009} then the most reaso-%
nable homotopy interval structure on $A^{\hand{\wedge}}$ I would think,  which possibly
may still yiled the desired "extra feature",  is the one defined in \var{\xout{\unsure{term}}}{}{typewriter deletion}{ps_134-8_006-v0010}
\margin{$\lfloor$}terms of the set $W_{A}$  of weak equivalences in $A^{\hand{\wedge}}$ as in section 54 (namely
$h_{W_{A}}$ ) .

\margin{$\lceil$}\hand{135.}  This handful of questions (some of which we met with before) \var{\xout{in}}{}{typewriter deletion}{ps_134-8_006-v0011}
\var{\xout{a slightly different light}}{}{typewriter deletion}{ps_134-8_006-v0012} is mainly a way of coming into touch again
with \var{\xout{some}}{}{typewriter deletion}{ps_134-8_006-v0013} the abelianization story,  which has been becoming somewhat
remote during the previous two months.  I am not sure I am going to

\margin{6}

% Visual collation: both supplied image halves inspected; see CHECK_DECISIONS_134_7_8.tsv.

%% END SOURCE ps_134-8_006
\ifnum\EditionMode=1
\psapparatusstart
\psentry{ps_134-8_006-v0001}{1}{typewriter overwriting}{\over[typed]{h}{h}}{h}{The superposed typed reading replaces the underlying reading.}
\psentry{ps_134-8_006-v0002}{2}{typewriter deletion}{\xout{of}}{}{The typewriter strikes out this material; it leaves the corrected running text.}
\psentry{ps_134-8_006-v0003}{3}{writing over}{\over{A}{a}}{\hand{a}}{The superposed handwritten reading replaces the underlying reading.}
\psentry{ps_134-8_006-v0004}{4}{typewriter overwriting}{\over[typed]{h}{h}}{h}{The superposed typed reading replaces the underlying reading.}
\psentry{ps_134-8_006-v0005}{5}{deletion}{\deleted{\xout{for}}}{}{The author cancels this material by hand; it leaves the corrected running text.}
\psentry{ps_134-8_006-v0006}{6}{typed insertion}{\inserted[typed]{replacing \hand{$\Delta$} by}}{replacing \hand{$\Delta$} by}{The typed insertion is incorporated at its indicated position; insertion brackets leave the running text.}
\psentry{ps_134-8_006-v0007}{7}{typed insertion}{\inserted[typed]{it would seem,}}{it would seem,}{The typed insertion is incorporated at its indicated position; insertion brackets leave the running text.}
\psentry{ps_134-8_006-v0008}{8}{insertion}{\inserted{the latter}}{\hand{the latter}}{The handwritten insertion is incorporated at its indicated position; insertion brackets leave the running text.}
\psentry{ps_134-8_006-v0009}{9}{typewriter deletion}{\xout{the reasonable}}{}{The typewriter strikes out this material; it leaves the corrected running text.}
\psentry{ps_134-8_006-v0010}{10}{typewriter deletion}{\xout{\unsure{term}}}{}{The typewriter strikes out this material; it leaves the corrected running text.}
\psentry{ps_134-8_006-v0011}{11}{typewriter deletion}{\xout{in}}{}{The typewriter strikes out this material; it leaves the corrected running text.}
\psentry{ps_134-8_006-v0012}{12}{typewriter deletion}{\xout{a slightly different light}}{}{The typewriter strikes out this material; it leaves the corrected running text.}
\psentry{ps_134-8_006-v0013}{13}{typewriter deletion}{\xout{some}}{}{The typewriter strikes out this material; it leaves the corrected running text.}
\psapparatusend
\fi
\pssource{ps_134-8_007}{134-8/7}{https://grothendieck.umontpellier.fr/web/viewer.html?file=../134-8.pdf\#page=8}
\psnumber{ps_134-8_007-v0001}{1}
\psnumber{ps_134-8_007-v0002}{2}
\psnumber{ps_134-8_007-v0003}{3}
\psnumber{ps_134-8_007-v0004}{4}
\psnumber{ps_134-8_007-v0005}{5}
\psnumber{ps_134-8_007-v0006}{6}
\psnumber{ps_134-8_007-v0007}{7}
\psnumber{ps_134-8_007-v0008}{8}
\psnumber{ps_134-8_007-v0009}{9}
\psnumber{ps_134-8_007-v0010}{10}
\psnumber{ps_134-8_007-v0011}{11}
\psnumber{ps_134-8_007-v0012}{12}
\psnumber{ps_134-8_007-v0013}{13}
%% BEGIN SOURCE ps_134-8_007
% Pursuing Stacks (cotes 134-2 to 134-8): archive cote 134-8, site page 7 (PDF page 8), typescript page 561.
% typescript. Draft compared: ps7.tex:297-376.
% AI-transcribed by Claude Opus 5.5 (claude-opus-5-5, adaptive thinking, effort high), Message Batches API, batch ps4 (msgbatch_012Aoc1NuskATPsyBeKSYK7J), 2026-11-12; reviewed against both supplied scan halves; see CHECK_DECISIONS.tsv.
\scanpage{134-8/7}
\mspage{561}

make any attempt,  now or later,  to come to an answer.  The last one
anyhow seems closely related to the formalism of closed model structures
on categories $A^{\hand{\wedge}}$,  and the proper place for dealing with it would seem
to be rather next chapter VII,  where such structures are going to be
\var{\xout{introduced}}{}{typewriter deletion}{ps_134-8_007-v0001} studied.  What I would like to do however here-and-now,  is
to establish at last the long promised relationship ``integrators are
abelianizators'',  which we have kept t\var{\over[typed]{\illegible{1}}{ur}}{ur}{typewriter overwriting}{ps_134-8_007-v0002}ning around and postponing
ever since section 99,  when those integrators were finally introduced
mainly for this purpose (of furnishing us with ``abelianizators'').

\var{\inserted[typed]{First of all,}}{First of all,}{typed insertion}{ps_134-8_007-v0003} I should be more outspoken than I have been before,  \var{\xout{first of}}{}{typewriter deletion}{ps_134-8_007-v0004} in
defining the ``\uline{abelianization functor}'' \var{\inserted{(or ``\uline{absolute Whitehead functor}''):}}{\hand{(or ``\uline{absolute Whitehead functor}''):}}{insertion}{ps_134-8_007-v0005}
\[ (1) \qquad \hand{\uline{Wh}} \ : \ (\mathrm{Hot}\hand{)} \overset{\mathrm{def}}{=} \underline{W}^{-1}(\mathrm{Cat}) \hand{\longrightarrow} \mathrm{Hotab} \overset{\mathrm{def}}{=} \mathrm{D.(Ab)} \quad , \]
without any use of the semisimplicial machinery\var{\deleted{,}}{}{deletion}{ps_134-8_007-v0006} which\hand{,} at the beginning
of our reflections\hand{,} ha\var{\over{\unsure{s}}{d}}{\hand{d}}{writing over}{ps_134-8_007-v0007} rather obscured the picture\var{\deleted{,}}{}{deletion}{ps_134-8_007-v0008} (section 92).  Defi-%
ning such a functor amounts to defining a functor
\[ (2) \qquad \mathrm{LH.} : \ (\mathrm{Cat}\hand{)} \hand{\longrightarrow} \mathrm{D.(Ab)} \quad , \]
or ``total homology functor'' ,  which should take weak equivalences into
isomorphisms.  For an\var{\xout{\illegible{1}}}{}{typewriter deletion}{ps_134-8_007-v0009} object A in (Cat),  we define
\[ (3) \qquad \mathrm{LH.}(A) = \mathrm{LH.}(A,\underline{Z}_A) = \mathrm{L}p^{ab}_{A!}\,(\underline{Z}_A) \ , \]
where $\underline{Z}_A$ is the constant abelian presheaf on A with value $\underline{Z}$ ,  and
\[ p_A : A \hand{\longrightarrow} e \]
is the map to the final object of (Cat).  We have to define the functo-%
rial dependance on A ,  More generally ,  for pairs
\[ (A,F) \ , \ \text{with } \var{\over[typed]{a}{A}}{A}{typewriter overwriting}{ps_134-8_007-v0010} \text{ in (Cat)} \ , \ F \text{ in } A^{\hand{\wedge}}_{ab} \quad , \]
the expression
\[ \mathrm{LH.}(A,F) \ \text{ in } \ \mathrm{D.(Ab)} \]
is functorial with respect to the pair (A,F),  where a map
\[ (A,F) \hand{\longrightarrow} (A',F') \]
is defined to be a pair  \var{\xout{$f\ \illegible{1}\ A\ \illegible{1}\ A',$}}{}{typewriter deletion}{ps_134-8_007-v0011}  (f,u),  where
\[ (4) \qquad f: A \hand{\to} A' \ , \ u: F \hand{\to} u^{\hand{*}}(F') \ , \]
the composition of \var{\inserted[typed]{maps}}{maps}{typed insertion}{ps_134-8_007-v0012} \var{\xout{arrows}}{}{typewriter deletion}{ps_134-8_007-v0013} being the obvious  one.  To see that such
a pair  defines a map
\[ (5) \qquad \mathrm{LH.}(f,u): \ \mathrm{LH.}(A,F) \hand{\longrightarrow} \mathrm{LH.}(A',F') \ , \]

\margin{7}

% Visual collation: both supplied image halves inspected; see CHECK_DECISIONS_134_7_8.tsv.

%% END SOURCE ps_134-8_007
\ifnum\EditionMode=1
\psapparatusstart
\psentry{ps_134-8_007-v0001}{1}{typewriter deletion}{\xout{introduced}}{}{The typewriter strikes out this material; it leaves the corrected running text.}
\psentry{ps_134-8_007-v0002}{2}{typewriter overwriting}{\over[typed]{\illegible{1}}{ur}}{ur}{The superposed typed reading replaces the underlying reading.}
\psentry{ps_134-8_007-v0003}{3}{typed insertion}{\inserted[typed]{First of all,}}{First of all,}{The typed insertion is incorporated at its indicated position; insertion brackets leave the running text.}
\psentry{ps_134-8_007-v0004}{4}{typewriter deletion}{\xout{first of}}{}{The typewriter strikes out this material; it leaves the corrected running text.}
\psentry{ps_134-8_007-v0005}{5}{insertion}{\inserted{(or ``\uline{absolute Whitehead functor}''):}}{\hand{(or ``\uline{absolute Whitehead functor}''):}}{The handwritten insertion is incorporated at its indicated position; insertion brackets leave the running text.}
\psentry{ps_134-8_007-v0006}{6}{deletion}{\deleted{,}}{}{The author cancels this material by hand; it leaves the corrected running text.}
\psentry{ps_134-8_007-v0007}{7}{writing over}{\over{\unsure{s}}{d}}{\hand{d}}{The superposed handwritten reading replaces the underlying reading.}
\psentry{ps_134-8_007-v0008}{8}{deletion}{\deleted{,}}{}{The author cancels this material by hand; it leaves the corrected running text.}
\psentry{ps_134-8_007-v0009}{9}{typewriter deletion}{\xout{\illegible{1}}}{}{The typewriter strikes out this material; it leaves the corrected running text.}
\psentry{ps_134-8_007-v0010}{10}{typewriter overwriting}{$\over[typed]{a}{A}$}{$A$}{The superposed typed reading replaces the underlying reading.}
\psentry{ps_134-8_007-v0011}{11}{typewriter deletion}{\xout{$f\ \illegible{1}\ A\ \illegible{1}\ A',$}}{}{The typewriter strikes out this material; it leaves the corrected running text.}
\psentry{ps_134-8_007-v0012}{12}{typed insertion}{\inserted[typed]{maps}}{maps}{The typed insertion is incorporated at its indicated position; insertion brackets leave the running text.}
\psentry{ps_134-8_007-v0013}{13}{typewriter deletion}{\xout{arrows}}{}{The typewriter strikes out this material; it leaves the corrected running text.}
\psapparatusend
\fi
\pssource{ps_134-8_008}{134-8/8}{https://grothendieck.umontpellier.fr/web/viewer.html?file=../134-8.pdf\#page=9}
\psnumber{ps_134-8_008-v0001}{1}
\psnumber{ps_134-8_008-v0002}{2}
\psnumber{ps_134-8_008-v0003}{3}
\psnumber{ps_134-8_008-v0004}{4}
\psnumber{ps_134-8_008-v0005}{5}
\psnumber{ps_134-8_008-v0006}{6}
\psnumber{ps_134-8_008-v0007}{7}
\psnumber{ps_134-8_008-v0008}{8}
\psnumber{ps_134-8_008-v0009}{9}
\psnumber{ps_134-8_008-v0010}{10}
%% BEGIN SOURCE ps_134-8_008
% Pursuing Stacks (cotes 134-2 to 134-8): archive cote 134-8, site page 8 (PDF page 9), typescript page 562.
% typescript. Draft compared: ps7.tex:350-423.
% AI-transcribed by Claude Opus 5.5 (claude-opus-5-5, adaptive thinking, effort high), Message Batches API, batch ps4 (msgbatch_012Aoc1NuskATPsyBeKSYK7J), 2026-11-12; reviewed against both supplied scan halves; see CHECK_DECISIONS.tsv.
\scanpage{134-8/8}
\mspage{562}

we use
\[ p_A = p_{A'} \hand{\circ} f \quad , \quad \var{\xout{\text{hence }Lp}}{}{typewriter deletion}{ps_134-8_008-v0001} \]
hence
\[ \var{\xout{\illegible{3}}}{}{typewriter deletion}{ps_134-8_008-v0002} \]
\[ \hand{(*)} \qquad (p_A)_!^{ab} \hand{\simeq} (p_{A'})_!^{ab} \hand{\circ} f_!^{ab} \quad , \quad \text{and} \quad L(p_A)_!^{ab} \hand{\simeq} L(p_{A'})_!^{ab} \hand{\circ} Lf_!^{ab} \quad , \]
taking into account that $f_!^{ab}$ maps projectives to projectives . Hence
we get
\[ LH.(A,F) \overset{def}{=} L(p_A)_!^{ab}(F) \hand{\simeq} L(p^{\var{\deleted{\illegible{1}}}{}{deletion}{ps_134-8_008-v0003}}_{A'})_!^{ab}(Lf_!^{ab}(F)) \overset{def}{=} LH.(A',Lf_!^{ab}(F)) \unsure{.} \]
Thus, in order to get (5), we need only define a map in $D^-(A'{}^{\hand{\wedge}}_{ab})$
\[ (6) \qquad Lf_!^{ab}(F) \hand{\longrightarrow} F' \ , \]
which will be obtained as the composition
\[ (6') \qquad Lf_!^{ab}(F) \hand{\longrightarrow} f_!^{ab}(F) \hand{\longrightarrow} F' \ , \]
where the first map in (6') is the canonical augmentation map towards
the $H_o$ object, and where the second corresponds to u in (4) by adjuncti%
on. \var{\over[typed]{T}{T}}{T}{typewriter overwriting}{ps_134-8_008-v0004}his defines the map (5), and compatibility with compositions
should be a tautology. Hence the functor (2). To get (1), we still
have to check that when
\[ f: \ A \hand{\longrightarrow} A' \]
is a weak equivalence, then
\[ LH.(f): \ LH.(A) \hand{\longrightarrow} LH.(A') \]
is an isomorphism in D.(Ab). It amounts to the same to check that for
any object $K^\bullet$ in $D^+(Ab)$, the corresponding map between the RHom's
with values in $K^\bullet$ is \var{\xout{an quasi iso}}{}{typewriter deletion}{ps_134-8_008-v0005} an isomorphism in $D^+(Ab)$. B\var{\over[typed]{u}{u}}{u}{typewriter overwriting}{ps_134-8_008-v0006}t \var{\xout{the}}{}{typewriter deletion}{ps_134-8_008-v0007}
the latter map can be identified with the map for \uline{cohomology}
\[ RH^\bullet(A',K^\bullet_{A'} \ ) \hand{\longrightarrow} RH^\bullet(A, K^\bullet_A \ ) \quad , \]
with coefficients in the constant complex of presheaves defined by $K^\bullet$
on $A'$ and on $A$ , which is an isomorphism, by the very definition
of weak equivalences in (Cat) via cohomology.

Now the statement ``an integrator \var{\xout{for A}}{}{typewriter deletion}{ps_134-8_008-v0008} is an abelianizator'' may
be rephrased rather evidently, without any reference to a given inte-%
grator, as \var{\xout{bein}}{}{typewriter deletion}{ps_134-8_008-v0009} merely the commutativity, up to canonical isomorphism,
of the \var{\xout{corres}}{}{typewriter deletion}{ps_134-8_008-v0010} following diagram for a given A in (Cat) :

\margin{8}

% Visual collation: both supplied image halves inspected; see CHECK_DECISIONS_134_7_8.tsv.

%% END SOURCE ps_134-8_008
\ifnum\EditionMode=1
\psapparatusstart
\psentry{ps_134-8_008-v0001}{1}{typewriter deletion}{$\xout{\text{hence }Lp}$}{}{The typewriter strikes out this material; it leaves the corrected running text.}
\psentry{ps_134-8_008-v0002}{2}{typewriter deletion}{$\xout{\illegible{3}}$}{}{The typewriter strikes out this material; it leaves the corrected running text.}
\psentry{ps_134-8_008-v0003}{3}{deletion}{$\deleted{\illegible{1}}$}{}{The author cancels this material by hand; it leaves the corrected running text.}
\psentry{ps_134-8_008-v0004}{4}{typewriter overwriting}{\over[typed]{T}{T}}{T}{The superposed typed reading replaces the underlying reading.}
\psentry{ps_134-8_008-v0005}{5}{typewriter deletion}{\xout{an quasi iso}}{}{The typewriter strikes out this material; it leaves the corrected running text.}
\psentry{ps_134-8_008-v0006}{6}{typewriter overwriting}{\over[typed]{u}{u}}{u}{The superposed typed reading replaces the underlying reading.}
\psentry{ps_134-8_008-v0007}{7}{typewriter deletion}{\xout{the}}{}{The typewriter strikes out this material; it leaves the corrected running text.}
\psentry{ps_134-8_008-v0008}{8}{typewriter deletion}{\xout{for A}}{}{The typewriter strikes out this material; it leaves the corrected running text.}
\psentry{ps_134-8_008-v0009}{9}{typewriter deletion}{\xout{bein}}{}{The typewriter strikes out this material; it leaves the corrected running text.}
\psentry{ps_134-8_008-v0010}{10}{typewriter deletion}{\xout{corres}}{}{The typewriter strikes out this material; it leaves the corrected running text.}
\psapparatusend
\fi
\pssource{ps_134-8_009}{134-8/9}{https://grothendieck.umontpellier.fr/web/viewer.html?file=../134-8.pdf\#page=10}
\psnumber{ps_134-8_009-v0001}{1}
\psnumber{ps_134-8_009-v0002}{2}
\psnumber{ps_134-8_009-v0003}{3}
\psnumber{ps_134-8_009-v0004}{4}
\psnumber{ps_134-8_009-v0005}{5}
\psnumber{ps_134-8_009-v0006}{6}
\psnumber{ps_134-8_009-v0007}{7}
\psnumber{ps_134-8_009-v0008}{8}
%% BEGIN SOURCE ps_134-8_009
% Pursuing Stacks (cotes 134-2 to 134-8): archive cote 134-8, site page 9 (PDF page 10), typescript page 563.
% typescript. Draft compared: ps7.tex:405-485.
% AI-transcribed by Claude Opus 5.5 (claude-opus-5-5, adaptive thinking, effort high), Message Batches API, batch ps4 (msgbatch_012Aoc1NuskATPsyBeKSYK7J), 2026-11-12; reviewed against both supplied scan halves; see CHECK_DECISIONS.tsv.
\scanpage{134-8/9}
\mspage{563}

\margin{\textcircled{$\varphi_A$}}
\diagram{
% Hand-drawn arrows are marked blue; typed objects remain plain.
\begin{tikzcd}[column sep=large]
A^{\hand{\wedge}} \arrow[color=blue,r, "\hand{\phi}_A"] \arrow[color=blue,d, "Wh_A"'] \& (Hot) \arrow[color=blue,d, "\underline{Wh}"] \\
A^{\hand{\wedge}}_{ab} \arrow[color=blue,r, "LH.(A{,}-)"'] \& (Hotab) = D.(Ab)
\end{tikzcd}
}
(7) \hfill ,

or equivalently,  of the corresponding diagram  where the categories
$A^{\hand{\wedge}}$, $A^{\hand{\wedge}}_{ab}$  are replaced by the relevant localizqtions :

\diagram{
\begin{tikzcd}
Hot_A \arrow[color=blue,r] \arrow[color=blue,d] \& Hot \arrow[color=blue,d] \\
Hotab_A \arrow[color=blue,r] \& H\overset{\unsure{o|u}}{\underset{\unsure{o}}{o}}tab
\end{tikzcd}
}
(7') \hfill .

The left vertical arrow in (7)  is of course the trivial abelianization
functor \hand{in $A^{\wedge}$} :
\[ Wh_A :\ A^{\hand{\wedge}} \hand{\longrightarrow} A^{\hand{\wedge}}_{ab} \quad , \qquad X \hand{\longmapsto} \underline{Z}^{(X)} = (\ a \hand{\mapsto} \underline{Z}^{(X(a))}\ ) \quad . \]
\var{\over{g}{G}}{\hand{G}}{writing over}{ps_134-8_009-v0001}oing back to the definitions,  the commutativity of (7) up to isomor-%
phism,  means that for X in $A^{\hand{\wedge}}$,  there is a canonical isomorphism
\[ (8) \qquad LH.(A_{/X}\ ,\ \underline{Z}\ ) \hand{\simeq} LH.(A,\ \underline{Z}^{(X)}) \quad , \]
functorial with respect to X.   To define (8),  let
\[ \var{\deleted{f_X\ \text{or }}}{}{deletion}{ps_134-8_009-v0002}f :\ A' \overset{def}{=} A_{/X} \hand{\longrightarrow} A \]
be the canonical functor,  then we get (using (*) of previous page,
with the roles of A and A' reversed)
\[ LH.(A',\ \underline{Z}_{A'}) \hand{\simeq} LH.(A,\ Lf^{ab}_{!}(\underline{Z}_{A'})) \quad , \]
and the relation (8) follows from the more precise relation
\[ (9) \qquad Lf^{ab}_{!}(\underline{Z}_{A'}) \xrightarrow{\hand{\sim}} \underline{Z}^{(X)} \quad . \]
To get \var{\xout{the}}{}{typewriter deletion}{ps_134-8_009-v0003} (9),  we remark that \var{\inserted[typed]{we have the tautological relation}}{we have the tautological relation}{typed insertion}{ps_134-8_009-v0004} \var{\xout{the right hand side is just}}{}{typewriter deletion}{ps_134-8_009-v0005}
\[ (9') \qquad f^{ab}_{!}(\underline{Z}_{A'}) \hand{\simeq} \underline{Z}^{(X)} \quad , \]
hence the \var{\xout{arro}}{}{typewriter deletion}{ps_134-8_009-v0006} map (9).  To prove it is an isomorphism amounts to pro-%
ving
\[ \var{\xout{(10)}}{}{typewriter deletion}{ps_134-8_009-v0007} \qquad L_i f^{ab}_{!}(\underline{Z}_{A'}) = 0 \quad \text{for} \ \ i \hand{>} 0 \ , \]
but we have indeed
\[ (10) \qquad L_i f^{ab}_{!} \ = 0 \quad \text{for} \ \ i \hand{>} 0 \ , \ \ \text{i.e.} \ \ f^{ab}_{!} \ \ \text{is \uline{exact}} \quad , \]
not only right exact,  a rather special feature\hand{,} \var{\inserted{valid}}{\hand{valid}}{insertion}{ps_134-8_009-v0008} for a localization

\margin{9}

% Visual collation: both supplied image halves inspected; see CHECK_DECISIONS_134_7_8.tsv.

%% END SOURCE ps_134-8_009
\ifnum\EditionMode=1
\psapparatusstart
\psentry{ps_134-8_009-v0001}{1}{writing over}{\over{g}{G}}{\hand{G}}{The superposed handwritten reading replaces the underlying reading.}
\psentry{ps_134-8_009-v0002}{2}{deletion}{$\deleted{f_X\ \text{or }}$}{}{The author cancels this material by hand; it leaves the corrected running text.}
\psentry{ps_134-8_009-v0003}{3}{typewriter deletion}{\xout{the}}{}{The typewriter strikes out this material; it leaves the corrected running text.}
\psentry{ps_134-8_009-v0004}{4}{typed insertion}{\inserted[typed]{we have the tautological relation}}{we have the tautological relation}{The typed insertion is incorporated at its indicated position; insertion brackets leave the running text.}
\psentry{ps_134-8_009-v0005}{5}{typewriter deletion}{\xout{the right hand side is just}}{}{The typewriter strikes out this material; it leaves the corrected running text.}
\psentry{ps_134-8_009-v0006}{6}{typewriter deletion}{\xout{arro}}{}{The typewriter strikes out this material; it leaves the corrected running text.}
\psentry{ps_134-8_009-v0007}{7}{typewriter deletion}{$\xout{(10)}$}{}{The typewriter strikes out this material; it leaves the corrected running text.}
\psentry{ps_134-8_009-v0008}{8}{insertion}{\inserted{valid}}{\hand{valid}}{The handwritten insertion is incorporated at its indicated position; insertion brackets leave the running text.}
\psapparatusend
\fi
\pssource{ps_134-8_010}{134-8/10}{https://grothendieck.umontpellier.fr/web/viewer.html?file=../134-8.pdf\#page=11}
\psnumber{ps_134-8_010-v0001}{1}
\psnumber{ps_134-8_010-v0002}{2}
\psnumber{ps_134-8_010-v0003}{3}
\psnumber{ps_134-8_010-v0004}{4}
\psnumber{ps_134-8_010-v0005}{5}
\psnumber{ps_134-8_010-v0006}{6}
%% BEGIN SOURCE ps_134-8_010
% Pursuing Stacks (cotes 134-2 to 134-8): archive cote 134-8, site page 10 (PDF page 11), typescript page 564.
% typescript. Draft compared: ps7.tex:464-522.
% AI-transcribed by Claude Opus 5.5 (claude-opus-5-5, adaptive thinking, effort high), Message Batches API, batch ps4 (msgbatch_012Aoc1NuskATPsyBeKSYK7J), 2026-11-12; reviewed against both supplied scan halves; see CHECK_DECISIONS.tsv.
\scanpage{134-8/10}
\mspage{564}

functor like $f : A_{/X} \hand{\to} A$ , namely for a functor which is fibering
(in the s ense of \var{\xout{fibered}}{}{typewriter deletion}{ps_134-8_010-v0001} the theory of ``fibered categories'') and has
discrete fibers. It comes from the explicit description of $f^{ab}_{!}$ ,
as
\[ (11) \qquad f^{ab}_{!}(F) = ( a \hand{\mapsto} \hand{\bigoplus}_{u\hand{\in} X(a)} F(a)_u ) \qquad , \]
where, for a group object $F$ in $A^{\hand{\wedge}}_{/X} \hand{\simeq} (A_{/X})^{\overset{\var{\deleted{\illegible{1}}}{}{deletion}{ps_134-8_010-v0002}}{\hand{\wedge}}}$ , and $u$ in $X(a)$ (defi\var{\xout{n}}{}{typewriter deletion}{ps_134-8_010-v0003}%
ning therefore an object of $A_{/X}$ ) $F(a)_u$ = fiber of $F(a)$ at $u \hand{\in} X(a)$\hand{,}
is the corresponding abelian group. I'll leave the proof of (11) to
the reader, it should be more or less tautological.

Once the \var{\inserted{whole}}{\hand{whole}}{insertion}{ps_134-8_010-v0004} proof is written down, it looks so simple that I feel
rather stupid and can't quite understand why I have turned around it
for so long, rather than writing it down right away more than three
months ago ! T$\overset{\illegible{1}}{\underset{h}{h}}$e reason surely is that I have been accustomed so
strongly to expressing everything via cohomology rather than homology,
that there has been something like a block against doing work homologi-%
cally,when it is homology which is involved. This block has remained
even after I took the trouble of telling myself quite outspokenly
(in section 100) that homology was just as important and meaningful
as cohomology, and more specifically still (in section 103) that the
proof I had in mind first\hand{,} via Quillen's result about $A \hand{\simeq} A^{o}$ in (Hot)
and via cohomology, was an ``awkward'', an ``upside down'' one\unsure{.|,} The
review on abelianization I went into just after made things rather
worse in a sense, as \var{\xout{it kind of}}{}{typewriter deletion}{ps_134-8_010-v0005} there I took great pains to make
the point that homology and cohomology were just one and the same thing
(so why bother about homology !). Still I \uline{did} develop some typically
homologically flavored formalism with the $*_k$ operation, and I hope
that at the end that block of mine is going to recede \var{\xout{some}}{}{typewriter deletion}{ps_134-8_010-v0006} ...

\margin{10}

% Visual collation: both supplied image halves inspected; see CHECK_DECISIONS_134_7_8.tsv.

%% END SOURCE ps_134-8_010
\ifnum\EditionMode=1
\psapparatusstart
\psentry{ps_134-8_010-v0001}{1}{typewriter deletion}{\xout{fibered}}{}{The typewriter strikes out this material; it leaves the corrected running text.}
\psentry{ps_134-8_010-v0002}{2}{deletion}{$\deleted{\illegible{1}}$}{}{The author cancels this material by hand; it leaves the corrected running text.}
\psentry{ps_134-8_010-v0003}{3}{typewriter deletion}{\xout{n}}{}{The typewriter strikes out this material; it leaves the corrected running text.}
\psentry{ps_134-8_010-v0004}{4}{insertion}{\inserted{whole}}{\hand{whole}}{The handwritten insertion is incorporated at its indicated position; insertion brackets leave the running text.}
\psentry{ps_134-8_010-v0005}{5}{typewriter deletion}{\xout{it kind of}}{}{The typewriter strikes out this material; it leaves the corrected running text.}
\psentry{ps_134-8_010-v0006}{6}{typewriter deletion}{\xout{some}}{}{The typewriter strikes out this material; it leaves the corrected running text.}
\psapparatusend
\fi
\pssource{ps_134-8_011}{134-8/11}{https://grothendieck.umontpellier.fr/web/viewer.html?file=../134-8.pdf\#page=12}
\psnumber{ps_134-8_011-v0001}{1}
\psnumber{ps_134-8_011-v0002}{2}
\psnumber{ps_134-8_011-v0003}{3}
\psnumber{ps_134-8_011-v0004}{4}
\psnumber{ps_134-8_011-v0005}{5}
\psnumber{ps_134-8_011-v0006}{6}
\psnumber{ps_134-8_011-v0007}{7}
\psnumber{ps_134-8_011-v0008}{8}
\psnumber{ps_134-8_011-v0009}{9}
\psnumber{ps_134-8_011-v0010}{10}
\psnumber{ps_134-8_011-v0011}{11}
\psnumber{ps_134-8_011-v0012}{12}
\psnumber{ps_134-8_011-v0013}{13}
\psnumber{ps_134-8_011-v0014}{14}
\psnumber{ps_134-8_011-v0015}{15}
\psnumber{ps_134-8_011-v0016}{16}
\psnumber{ps_134-8_011-v0017}{17}
%% BEGIN SOURCE ps_134-8_011
% Pursuing Stacks (cotes 134-2 to 134-8): archive cote 134-8, site page 11 (PDF page 12), typescript page 565.
% typescript. Draft compared: ps7.tex:504-567.
% AI-transcribed by Claude Opus 5.5 (claude-opus-5-5, adaptive thinking, effort high), Message Batches API, batch ps4 (msgbatch_012Aoc1NuskATPsyBeKSYK7J), 2026-11-12; reviewed against both supplied scan halves; see CHECK_DECISIONS.tsv.
\scanpage{134-8/11}
\mspage{565}

27.10.

136.    I just spent a couple of hours,  after reading the notes of last
night,  trying to get a better feeling of the basic homology operation
in the context of the basic localizer (Cat),  namely taking the left
derived functor  $Lf_!^{ab}$ of the ``unusual'' direct image functor for abelian
presheaves,  associated to a map
\[ f: A \hand{\longrightarrow} B \]
in (Cat).  This led me to read again the notes of section 92,  when I
unsuspectingly started on an ``afterthought''\hand{,} \var{\inserted{later gradually}}{\hand{later gradually}}{insertion}{ps_134-8_011-v0001} turning into a systematic
reflection on abelianization.  With a distance of nearly four months,
what strikes me most now in these notes is awkwardness of the approach\textsubscript{x}
f\var{\over[typed]{l}{o}}{o}{typewriter overwriting}{ps_134-8_011-v0002}llowed at start,  when yielding \var{\inserted[typed]{to}}{to}{typed insertion}{ps_134-8_011-v0003} the reflex of lazyness of describing
abelianization of homotopy types via the semisimplicial grindmill.  The
uneas\var{\over{\unsure{i}}{i}}{\hand{i}}{writing over}{ps_134-8_011-v0004}ness in these notes is obvious throughout  --  I kind of knew
``au fond'' that dragging in the particular test category \hand{$\Delta$} was rather
silly.  In section 100 only does it get clear that the best\hand{,}\var{\xout{\illegible{1} \unsure{intrinsic}}}{}{typewriter deletion}{ps_134-8_011-v0005}
description\hand{,}\var{\inserted{for abelianization}}{\hand{for abelianization}}{insertion}{ps_134-8_011-v0006} within the modelizer (Cat),\var{\deleted{for abelianization}}{}{deletion}{ps_134-8_011-v0007} is via the
unusual direct image $p_!^{ab}$ corre\var{\over[typed]{\illegible{1}}{s}}{s}{typewriter overwriting}{ps_134-8_011-v0008}ponding to the projection
\[ p=p_A \; : A \hand{\longrightarrow} e \]
of the ``model'' A in (Cat) to the final object (formulae (11) and (12)
page 569),  by applying $p_!^{ab}$ to a ``cointegrator''  $L^{A}.$ for A  namely to a
projective resolution of  $\underline{Z}_A$  (where A is written B by the way,  as \var{\xout{we}}{}{typewriter deletion}{ps_134-8_011-v0009}
I had been led before to replace a given A in (Cat) by it's ``dual'' or
opposite B=A$^{o}$ ,  bound as \var{\deleted{we were}}{}{deletion}{ps_134-8_011-v0010} \var{\inserted{I was}}{\hand{I was}}{insertion}{ps_134-8_011-v0011} for interpreting ``integrators'' for A
in terms of ``cointegrators'' for B ...);  and in the next section the
step is finally taken (against the ``block'' !)  to write
\[ p_!^{ab}(L^{B}.) = Lp_!^{ab}(\underline{Z}_B) \quad , \]
which inserts abelianization into the familiar formalism of (left)
derived functors of stan\-dar\var{\over{\unsure{t|d}}{d}}{\hand{d}}{writing over}{ps_134-8_011-v0012} functors.  The reasonable next thing to do
was of course \var{\xout{to}}{}{typewriter deletion}{ps_134-8_011-v0013} what I finally did \var{\inserted{only}}{\hand{only}}{insertion}{ps_134-8_011-v0014} yesterday,  namely c\var{\over{\illegible{1}}{h}}{\hand{h}}{writing over}{ps_134-8_011-v0015}eck the commu-%
tativity of the diagram (7) of p. 563,  namely compatibility of this
notion of abelianization (or homology) with Whitehead's abelianization,
when working \var{\xout{in}}{}{typewriter deletion}{ps_134-8_011-v0016} with models coming from A$^{\hand{\wedge}}$,  A any object in (Cat).
This by the way,  when applied to the case  A = \hand{$\Delta$} ,  gives at once
the equivalence of the intrinsic definition of abelianization,  with
the semisimplicial one we started with  --  provided we remark that in
this case,  $Lp_!^{ab}$ \var{\inserted{(rather, it's restriction to $\hat{\Delta}_{ab}$)}}{\hand{(rather, it's restriction to $\hat{\Delta}_{ab}$)}}{insertion}{ps_134-8_011-v0017} may be equally interpreted as the Kan-Dold-Puppe
\margin{11}

% Visual collation: both supplied image halves inspected; see CHECK_DECISIONS_134_7_8.tsv.

%% END SOURCE ps_134-8_011
\ifnum\EditionMode=1
\psapparatusstart
\psentry{ps_134-8_011-v0001}{1}{insertion}{\inserted{later gradually}}{\hand{later gradually}}{The handwritten insertion is incorporated at its indicated position; insertion brackets leave the running text.}
\psentry{ps_134-8_011-v0002}{2}{typewriter overwriting}{\over[typed]{l}{o}}{o}{The superposed typed reading replaces the underlying reading.}
\psentry{ps_134-8_011-v0003}{3}{typed insertion}{\inserted[typed]{to}}{to}{The typed insertion is incorporated at its indicated position; insertion brackets leave the running text.}
\psentry{ps_134-8_011-v0004}{4}{writing over}{\over{\unsure{i}}{i}}{\hand{i}}{The superposed handwritten reading replaces the underlying reading.}
\psentry{ps_134-8_011-v0005}{5}{typewriter deletion}{\xout{\illegible{1} \unsure{intrinsic}}}{}{The typewriter strikes out this material; it leaves the corrected running text.}
\psentry{ps_134-8_011-v0006}{6}{insertion}{\inserted{for abelianization}}{\hand{for abelianization}}{The handwritten insertion is incorporated at its indicated position; insertion brackets leave the running text.}
\psentry{ps_134-8_011-v0007}{7}{deletion}{\deleted{for abelianization}}{}{The author cancels this material by hand; it leaves the corrected running text.}
\psentry{ps_134-8_011-v0008}{8}{typewriter overwriting}{\over[typed]{\illegible{1}}{s}}{s}{The superposed typed reading replaces the underlying reading.}
\psentry{ps_134-8_011-v0009}{9}{typewriter deletion}{\xout{we}}{}{The typewriter strikes out this material; it leaves the corrected running text.}
\psentry{ps_134-8_011-v0010}{10}{deletion}{\deleted{we were}}{}{The author cancels this material by hand; it leaves the corrected running text.}
\psentry{ps_134-8_011-v0011}{11}{insertion}{\inserted{I was}}{\hand{I was}}{The handwritten insertion is incorporated at its indicated position; insertion brackets leave the running text.}
\psentry{ps_134-8_011-v0012}{12}{writing over}{\over{\unsure{t|d}}{d}}{\hand{d}}{The superposed handwritten reading replaces the underlying reading.}
\psentry{ps_134-8_011-v0013}{13}{typewriter deletion}{\xout{to}}{}{The typewriter strikes out this material; it leaves the corrected running text.}
\psentry{ps_134-8_011-v0014}{14}{insertion}{\inserted{only}}{\hand{only}}{The handwritten insertion is incorporated at its indicated position; insertion brackets leave the running text.}
\psentry{ps_134-8_011-v0015}{15}{writing over}{\over{\illegible{1}}{h}}{\hand{h}}{The superposed handwritten reading replaces the underlying reading.}
\psentry{ps_134-8_011-v0016}{16}{typewriter deletion}{\xout{in}}{}{The typewriter strikes out this material; it leaves the corrected running text.}
\psentry{ps_134-8_011-v0017}{17}{insertion}{\inserted{(rather, it's restriction to $\hat{\Delta}_{ab}$)}}{\hand{(rather, it's restriction to $\hat{\Delta}_{ab}$)}}{The handwritten insertion is incorporated at its indicated position; insertion brackets leave the running text.}
\psapparatusend
\fi
\pssource{ps_134-8_012}{134-8/12}{https://grothendieck.umontpellier.fr/web/viewer.html?file=../134-8.pdf\#page=13}
\psnumber{ps_134-8_012-v0001}{1}
\psnumber{ps_134-8_012-v0002}{2}
\psnumber{ps_134-8_012-v0003}{3}
\psnumber{ps_134-8_012-v0004}{4}
\psnumber{ps_134-8_012-v0005}{5}
\psnumber{ps_134-8_012-v0006}{6}
\psnumber{ps_134-8_012-v0007}{7}
\psnumber{ps_134-8_012-v0008}{8}
\psnumber{ps_134-8_012-v0009}{9}
\psnumber{ps_134-8_012-v0010}{10}
\psnumber{ps_134-8_012-v0011}{11}
\psnumber{ps_134-8_012-v0012}{12}
\psnumber{ps_134-8_012-v0013}{13}
\psnumber{ps_134-8_012-v0014}{14}
\psnumber{ps_134-8_012-v0015}{15}
\psnumber{ps_134-8_012-v0016}{16}
\psnumber{ps_134-8_012-v0017}{17}
%% BEGIN SOURCE ps_134-8_012
% Pursuing Stacks (cotes 134-2 to 134-8): archive cote 134-8, site page 12 (PDF page 13), typescript page 566.
% typescript. Draft compared: ps7.tex:555-613.
% AI-transcribed by Claude Opus 5.5 (claude-opus-5-5, adaptive thinking, effort high), Message Batches API, batch ps4 (msgbatch_012Aoc1NuskATPsyBeKSYK7J), 2026-11-12; reviewed against both supplied scan halves; see CHECK_DECISIONS.tsv.
\scanpage{134-8/12}
\mspage{566}
functor (more accurately, the composition of the latter with the
localization functor $\mathrm{Ch}.(\mathrm{Ab})\hand{\longrightarrow}\allowbreak \mathrm{D}.(\mathrm{Ab}) = \mathrm{Hotab}$ ). T\var{\over[typed]{h}{h}}{h}{typewriter overwriting}{ps_134-8_012-v0001}e equivalence of
the two definitions \var{\inserted{of abelianization}}{\hand{of abelianization}}{insertion}{ps_134-8_012-v0002} is mentioned on the same p. 569, somewhat as a
chore I did'nt really feel then to dive into. \var{\xout{The}}{}{typewriter deletion}{ps_134-8_012-v0003} Besides the ``block''
against homology, the picture was being obscured too by the ``computati-%
onal'' idea I kept in mind, of expressing $\mathrm{Lp}_{!}^{ab}(F)$ for F in $A^{\hand{\wedge}}_{ab}$ in
terms of an ``integrator'' for A i.e. as $F *_{\underline{Z}} L^{B}.$ - whereas I should
have known best myself that for establishing formal properties relating
various derived functors, the particular \var{\xout{computational}}{}{typewriter deletion}{ps_134-8_012-v0004} approach\hand{es} used
for``co\var{\inserted[typed]{m}}{m}{typed insertion}{ps_134-8_012-v0005}puting'' them more or less eleg\var{\xout{g}}{}{typewriter deletion}{ps_134-8_012-v0006}antly \var{\over{\illegible{1}}{are}}{\hand{are}}{writing over}{ps_134-8_012-v0007} wholly irrelevant ...

O$\overset{\illegible{1}}{\underset{n}{n}}$e teaching I am getting out of all this, is that when expres-%
sing abelianization, or presumably any other kind of notion or operati-%
on of significance for homotopy types, one should be careful, for
any modelizer one chooses to work in, to dig out the description which
fits smoothly those particular models. Clearly, when working with
semisimplicial models, the description via tautological abelianization
$\mathrm{Wh}_{\hand{\Delta}}$ and using Kan-Dold-Puppe is the best. When working within the
modelizer (Cat) though, making the detour through $\hand{\Delta}^{\hand{\wedge}}$ is awkward and
makes us \var{\inserted{just}}{\hand{just}}{insertion}{ps_134-8_012-v0008} miss \var{\deleted{the understanding of}}{}{deletion}{ps_134-8_012-v0009} the relevant facts. Once we got th\var{\over{e}{is}}{\hand{is}}{writing over}{ps_134-8_012-v0010}
\var{\deleted{understanding}}{}{deletion}{ps_134-8_012-v0011}, it gets clear too what to do when $\hand{\Delta}$ is replaced by any
other object A in (Cat) \var{\xout{(which needs not even be a pseudotest category)}}{}{typewriter deletion}{ps_134-8_012-v0012}
when taking models in $A^{\hand{\wedge}}$ (never minding even whether $A^{\hand{\wedge}}$ is indeed a
``modelizer'' ): namely, apply $\mathrm{Wh}_{A}$ , and take total homology !

This brings to my mind another example, very similar indeed.
Some time after I got across Thomason's nice paper, showing that
(Cat) is a closed model category (see comments in section 87), I got
from Tim Porter another \var{\xout{re}}{}{typewriter deletion}{ps_134-8_012-v0013} reprint of Thomason's, where he gives the
description of ``homotopy colimits'' (or ``integration'', as I call it)
within the modelizer (Cat),\var{\xout{(compare}}{}{typewriter deletion}{ps_134-8_012-v0014} in terms of the total category
associated to a fibered category (compare section 69). There he grinds
through a tedious, highly technical proof, whereas the direct proof
when describing ``integration'' as the left adjoint functor to the
tautological ``inverse image'' functor, is more or less tautological
too. T\var{\over[typed]{h}{h}}{h}{typewriter overwriting}{ps_134-8_012-v0015}e reason for this awkwardness is again that, rather than
being content to work with the models as they are, Thomason refers to
the Bousfield-Kan description of colimits in the semisimplicial set-up, \var{\inserted{which he takes as his definition for colimits.}}{\hand{which he takes as his definition for colimits.}}{insertion}{ps_134-8_012-v0016}
I suspect that (Cat) is the one modelizer most ideally suited for expres%
-sing the ``integration'' ope\var{\over[typed]{\illegible{2}}{ra}}{ra}{typewriter overwriting}{ps_134-8_012-v0017}tion, and that the Bousfield-Kan descripti-%
on is just the obvious, not-quite-as-simple one which can be deduced
\margin{12}

% Visual collation: both supplied image halves inspected; see CHECK_DECISIONS_134_7_8.tsv.

%% END SOURCE ps_134-8_012
\ifnum\EditionMode=1
\psapparatusstart
\psentry{ps_134-8_012-v0001}{1}{typewriter overwriting}{\over[typed]{h}{h}}{h}{The superposed typed reading replaces the underlying reading.}
\psentry{ps_134-8_012-v0002}{2}{insertion}{\inserted{of abelianization}}{\hand{of abelianization}}{The handwritten insertion is incorporated at its indicated position; insertion brackets leave the running text.}
\psentry{ps_134-8_012-v0003}{3}{typewriter deletion}{\xout{The}}{}{The typewriter strikes out this material; it leaves the corrected running text.}
\psentry{ps_134-8_012-v0004}{4}{typewriter deletion}{\xout{computational}}{}{The typewriter strikes out this material; it leaves the corrected running text.}
\psentry{ps_134-8_012-v0005}{5}{typed insertion}{\inserted[typed]{m}}{m}{The typed insertion is incorporated at its indicated position; insertion brackets leave the running text.}
\psentry{ps_134-8_012-v0006}{6}{typewriter deletion}{\xout{g}}{}{The typewriter strikes out this material; it leaves the corrected running text.}
\psentry{ps_134-8_012-v0007}{7}{writing over}{\over{\illegible{1}}{are}}{\hand{are}}{The superposed handwritten reading replaces the underlying reading.}
\psentry{ps_134-8_012-v0008}{8}{insertion}{\inserted{just}}{\hand{just}}{The handwritten insertion is incorporated at its indicated position; insertion brackets leave the running text.}
\psentry{ps_134-8_012-v0009}{9}{deletion}{\deleted{the understanding of}}{}{The author cancels this material by hand; it leaves the corrected running text.}
\psentry{ps_134-8_012-v0010}{10}{writing over}{\over{e}{is}}{\hand{is}}{The superposed handwritten reading replaces the underlying reading.}
\psentry{ps_134-8_012-v0011}{11}{deletion}{\deleted{understanding}}{}{The author cancels this material by hand; it leaves the corrected running text.}
\psentry{ps_134-8_012-v0012}{12}{typewriter deletion}{\xout{(which needs not even be a pseudotest category)}}{}{The typewriter strikes out this material; it leaves the corrected running text.}
\psentry{ps_134-8_012-v0013}{13}{typewriter deletion}{\xout{re}}{}{The typewriter strikes out this material; it leaves the corrected running text.}
\psentry{ps_134-8_012-v0014}{14}{typewriter deletion}{\xout{(compare}}{}{The typewriter strikes out this material; it leaves the corrected running text.}
\psentry{ps_134-8_012-v0015}{15}{typewriter overwriting}{\over[typed]{h}{h}}{h}{The superposed typed reading replaces the underlying reading.}
\psentry{ps_134-8_012-v0016}{16}{insertion}{\inserted{which he takes as his definition for colimits.}}{\hand{which he takes as his definition for colimits.}}{The handwritten insertion is incorporated at its indicated position; insertion brackets leave the running text.}
\psentry{ps_134-8_012-v0017}{17}{typewriter overwriting}{\over[typed]{\illegible{2}}{ra}}{ra}{The superposed typed reading replaces the underlying reading.}
\psapparatusend
\fi
\pssource{ps_134-8_013}{134-8/13}{https://grothendieck.umontpellier.fr/web/viewer.html?file=../134-8.pdf\#page=14}
\psnumber{ps_134-8_013-v0001}{1}
\psnumber{ps_134-8_013-v0002}{2}
\psnumber{ps_134-8_013-v0003}{3}
\psnumber{ps_134-8_013-v0004}{4}
\psnumber{ps_134-8_013-v0005}{5}
\psnumber{ps_134-8_013-v0006}{6}
\psnumber{ps_134-8_013-v0007}{7}
\psnumber{ps_134-8_013-v0008}{8}
\psnumber{ps_134-8_013-v0009}{9}
\psnumber{ps_134-8_013-v0010}{10}
\psnumber{ps_134-8_013-v0011}{11}
\psnumber{ps_134-8_013-v0012}{12}
\psnumber{ps_134-8_013-v0013}{13}
\psnumber{ps_134-8_013-v0014}{14}
\psnumber{ps_134-8_013-v0015}{15}
\psnumber{ps_134-8_013-v0016}{16}
\psnumber{ps_134-8_013-v0017}{17}
%% BEGIN SOURCE ps_134-8_013
% Pursuing Stacks (cotes 134-2 to 134-8): archive cote 134-8, site page 13 (PDF page 14), typescript page 567.
% typescript. Draft compared: ps7.tex:600-659.
% AI-transcribed by Claude Opus 5.5 (claude-opus-5-5, adaptive thinking, effort high), Message Batches API, batch ps4 (msgbatch_012Aoc1NuskATPsyBeKSYK7J), 2026-11-12; reviewed against both supplied scan halves; see CHECK_DECISIONS.tsv.
\scanpage{134-8/13}
\mspage{567}

from the former,  using the relevant two functors between (Cat) and \hand{$\Delta$}$^{\hand{\wedge}}$
which allow to pass from one type of models to another. (Soon or late
I should check in Bousfi\var{\over[typed]{l}{e}}{e}{typewriter overwriting}{ps_134-8_013-v0001}ld-Kan's book\var{\xout{\unsure{s}}}{}{typewriter deletion}{ps_134-8_013-v0002} whether this is so or not ...)
Thomason however did the opposite,  and it is quite natural that he had
to pay for it by a fair amount of sweat!  \var{\xout{\illegible{1}}}{}{typewriter deletion}{ps_134-8_013-v0003} (Reference of the paper:
Homotopy colimits in the category of small categories,  Math. Proc.
Camb. Phil. Soc. (1979)\textsuperscript{\var{\xout{\unsure{t}}}{}{typewriter deletion}{ps_134-8_013-v0004}}, 85, 91,  p. 91 - 109.)

T$\overset{\illegible{1}}{\underset{h}{h}}$e proof written down yesterday for compatibility of abelianizati-%
on of homotopy types with the Whitehead abelianization functor within
a category A$^{\hand{\wedge}}$,  still goes through when replacing abelianization by
k-linearization,  with respect to an arbitrary ground ring k (not even
commutative).  It was'nt really worth while though to introduce a ring
k,  as the general result should follow at once from the case k=$\underline{Z}$ ,
by the formula
\[ (1) \qquad Lp^{k}_{!}(k_{A}) \hand{\simeq} Lp^{\underline{Z}}_{!}(\underline{Z}_{A}) \overset{L}{\otimes}_{\underline{Z}} k \quad , \]
where in the \var{\xout{first}}{}{typewriter deletion}{ps_134-8_013-v0005} \var{\inserted[typed]{left}}{left}{typed insertion}{ps_134-8_013-v0006} hand side we are taking the \var{\inserted[typed]{left}}{left}{typed insertion}{ps_134-8_013-v0007} derived functor for the
functor
\[ p^{k}_{!} \ : \ A^{\hand{\wedge}}_{k} \hand{\longrightarrow} (Ab_{k}) \]
generalizing $p^{ab}_{!} = p^{\underline{Z}}_{!}$ , \hand{(}\var{\over{where}{with}}{\hand{with}}{writing over}{ps_134-8_013-v0008} $p\colon A \hand{\to} e$ \var{\deleted{is}}{}{deletion}{ps_134-8_013-v0009} as above\hand{)}, and where in
the \var{\xout{second member}}{}{typewriter deletion}{ps_134-8_013-v0010} right hand side we are using the bing extension
functor
\[ \overset{L}{\otimes}_{\underline{Z}} \ : \ D\overset{\var{\deleted{\illegible{1}}}{}{deletion}{ps_134-8_013-v0011}}{.}(Ab) \hand{\longrightarrow} D.(Ab_{k}) \]
for the relevant derived categories.  This reminds me of the need of
developping \var{\xout{more}}{}{typewriter deletion}{ps_134-8_013-v0012} a more or less exhaustive formulaire around the
basic operations
\[ (2) \qquad Lf_{!} \ , \ f^{*} \ , \ Rf_{*} \quad , \]
including the familiar one for  the two latter,  valid \var{\xout{for}}{}{typewriter deletion}{ps_134-8_013-v0013} more general-%
ly for maps between ringed topoi,  and including also \var{\inserted{a}}{\hand{a}}{insertion}{ps_134-8_013-v0014} \hand{``}\uline{projection
formula}'' generalizing (1).  Such a formula will be no surprise surely
to a reader familiar with a duality context (such as étale cohomology,
of ``coherent'' cohomology of noetherian schemes say )\hand{,} where a formalism
of the ``four variances''  $\var{\over[typed]{F}{f}}{f}{typewriter overwriting}{ps_134-8_013-v0015}_{!}$ , $f^{*}$ , $f_{*}$ , $f^{!}$  and the \hand{``}two internal opera-%
tions\hand{''} $\overset{L}{\otimes}$  and R\uline{Hom}  can be developped -  it would seem that the formal
properties of the triple (2),  together with the two internal operations
just referred to,  are very close (for the least) \var{\over{of}{to}}{\hand{to}}{writing over}{ps_134-8_013-v0016} those of the
slightly richer one, \var{\inserted{in duality set-ups,}}{\hand{in duality set-ups,}}{insertion}{ps_134-8_013-v0017} including equally an ``unusual inverse image'' $f^{!}$,
\margin{no, to $Rf_!$} right adjoint to \uline{$Rf_{*}$}  (denoted sometimes simply by $f_{*}$ ).  This similari

\margin{13}

% Visual collation: both supplied image halves inspected; see CHECK_DECISIONS_134_7_8.tsv.

%% END SOURCE ps_134-8_013
\ifnum\EditionMode=1
\psapparatusstart
\psentry{ps_134-8_013-v0001}{1}{typewriter overwriting}{\over[typed]{l}{e}}{e}{The superposed typed reading replaces the underlying reading.}
\psentry{ps_134-8_013-v0002}{2}{typewriter deletion}{\xout{\unsure{s}}}{}{The typewriter strikes out this material; it leaves the corrected running text.}
\psentry{ps_134-8_013-v0003}{3}{typewriter deletion}{\xout{\illegible{1}}}{}{The typewriter strikes out this material; it leaves the corrected running text.}
\psentry{ps_134-8_013-v0004}{4}{typewriter deletion}{\xout{\unsure{t}}}{}{The typewriter strikes out this material; it leaves the corrected running text.}
\psentry{ps_134-8_013-v0005}{5}{typewriter deletion}{\xout{first}}{}{The typewriter strikes out this material; it leaves the corrected running text.}
\psentry{ps_134-8_013-v0006}{6}{typed insertion}{\inserted[typed]{left}}{left}{The typed insertion is incorporated at its indicated position; insertion brackets leave the running text.}
\psentry{ps_134-8_013-v0007}{7}{typed insertion}{\inserted[typed]{left}}{left}{The typed insertion is incorporated at its indicated position; insertion brackets leave the running text.}
\psentry{ps_134-8_013-v0008}{8}{writing over}{\over{where}{with}}{\hand{with}}{The superposed handwritten reading replaces the underlying reading.}
\psentry{ps_134-8_013-v0009}{9}{deletion}{\deleted{is}}{}{The author cancels this material by hand; it leaves the corrected running text.}
\psentry{ps_134-8_013-v0010}{10}{typewriter deletion}{\xout{second member}}{}{The typewriter strikes out this material; it leaves the corrected running text.}
\psentry{ps_134-8_013-v0011}{11}{deletion}{$\deleted{\illegible{1}}$}{}{The author cancels this material by hand; it leaves the corrected running text.}
\psentry{ps_134-8_013-v0012}{12}{typewriter deletion}{\xout{more}}{}{The typewriter strikes out this material; it leaves the corrected running text.}
\psentry{ps_134-8_013-v0013}{13}{typewriter deletion}{\xout{for}}{}{The typewriter strikes out this material; it leaves the corrected running text.}
\psentry{ps_134-8_013-v0014}{14}{insertion}{\inserted{a}}{\hand{a}}{The handwritten insertion is incorporated at its indicated position; insertion brackets leave the running text.}
\psentry{ps_134-8_013-v0015}{15}{typewriter overwriting}{$\over[typed]{F}{f}$}{$f$}{The superposed typed reading replaces the underlying reading.}
\psentry{ps_134-8_013-v0016}{16}{writing over}{\over{of}{to}}{\hand{to}}{The superposed handwritten reading replaces the underlying reading.}
\psentry{ps_134-8_013-v0017}{17}{insertion}{\inserted{in duality set-ups,}}{\hand{in duality set-ups,}}{The handwritten insertion is incorporated at its indicated position; insertion brackets leave the running text.}
\psapparatusend
\fi
\pssource{ps_134-8_014}{134-8/14}{https://grothendieck.umontpellier.fr/web/viewer.html?file=../134-8.pdf\#page=15}
\psnumber{ps_134-8_014-v0001}{1}
\psnumber{ps_134-8_014-v0002}{2}
\psnumber{ps_134-8_014-v0003}{3}
\psnumber{ps_134-8_014-v0004}{4}
\psnumber{ps_134-8_014-v0005}{5}
\psnumber{ps_134-8_014-v0006}{6}
\psnumber{ps_134-8_014-v0007}{7}
\psnumber{ps_134-8_014-v0008}{8}
\psnumber{ps_134-8_014-v0009}{9}
\psnumber{ps_134-8_014-v0010}{10}
\psnumber{ps_134-8_014-v0011}{11}
\psnumber{ps_134-8_014-v0012}{12}
\psnumber{ps_134-8_014-v0013}{13}
\psnumber{ps_134-8_014-v0014}{14}
\psnumber{ps_134-8_014-v0015}{15}
\psnumber{ps_134-8_014-v0016}{16}
\psnumber{ps_134-8_014-v0017}{17}
%% BEGIN SOURCE ps_134-8_014
% Pursuing Stacks (cotes 134-2 to 134-8): archive cote 134-8, site page 14 (PDF page 15), typescript page 568.
% typescript. Draft compared: ps7.tex:647-707.
% AI-transcribed by Claude Opus 5.5 (claude-opus-5-5, adaptive thinking, effort high), Message Batches API, batch ps4 (msgbatch_012Aoc1NuskATPsyBeKSYK7J), 2026-11-12; reviewed against both supplied scan halves; see CHECK_DECISIONS.tsv.
\scanpage{134-8/14}
\mspage{568}

\var{\xout{\unsure{it}}}{}{typewriter deletion}{ps_134-8_014-v0001} ty is a matter of course,  as far as the two last among the functors
(2),  together with the \var{\inserted{two}}{\hand{two}}{insertion}{ps_134-8_014-v0002} internal operations,  are concerned,  as in
both contexts ( homology and cbhomology formalism within (Cat) on the
one hand,  and the ``sweeping duality formalism'' on the other) \var{\xout{these}}{}{typewriter deletion}{ps_134-8_014-v0003}
\var{\xout{operations}}{}{typewriter deletion}{ps_134-8_014-v0004} the formulaire concerning these four operations
\[ f^{*} \ , \ \mathrm{R}f_{*} \ , \ \overset{L}{\otimes} \ , \ \mathrm{R}\underline{\mathrm{Hom}} \]
is no more, no less than just the relevant formulaire \var{\xout{for}}{}{typewriter deletion}{ps_134-8_014-v0005} in the context
of \var{\inserted{arbitrary}}{\hand{arbitrary}}{insertion}{ps_134-8_014-v0006} (commutatively) ringed topoi,  and maps between such.  The common
notation  $f_{!}$  ( occuring in $\mathrm{L}f_{!}$  in the (Cat) context,  in $\mathrm{R}f_{!}$ in the
``duality'' context)  is a very suggestive one,  for the least,  and I am
rather confident that \var{\inserted[typed]{most}}{most}{typed insertion}{ps_134-8_014-v0007} \var{\xout{the}}{}{typewriter deletion}{ps_134-8_014-v0008} reflexes (concerning formal behaviour of $f_{!}$
with respect to the other operations)  coming from one context,  should
be OK too in the other.  Whether this is just a mere formal analogy,
or whether there is a deeper \var{\inserted[typed]{relationship}}{relationship}{typed insertion}{ps_134-8_014-v0009} \var{\xout{identity}}{}{typewriter deletion}{ps_134-8_014-v0010} between the two kind of contexts,
I am at a loss at present to say.  It does'nt seem at all unlikely to
me that among arbitrary maps in (Cat),  one can single out some\hand{,} by
some kind of ``finite type'' condition,  for which a functor
\[ (*) \qquad f^{!} \ : \ D^{+}(B^{\wedge}_{k}) \hand{\longrightarrow} D^{+}(A^{\wedge}_{k}) \]
\margin{\illegible{1}}
can be defined (where $f: A \hand{\longrightarrow} B$),  right adjoint to the familiar $\mathrm{R}f_{*}$
functor ,  so that (2) can be completed to a sequence of \uline{four} functors
\[ (3) \qquad \mathrm{L}f_{!} \ , \ f^{*} \ , \ \mathrm{R}f_{*} \ , \ f^{!} \ , \]
forming a se\var{\over{q}{q}}{\hand{q}}{writing over}{ps_134-8_014-v0011}uence of mutually adjoiht functors between derived catego-%
ries,  in the usual sense (the functor immediately to the right of
another being it's right adjoint).  I faintly remember that Verdier
worked out such a formalism within the context of discrete or profinite
groups, or both,  in a Bourbaki talk \var{\inserted{he gave}}{\hand{he gave}}{insertion}{ps_134-8_014-v0012}, this being inspired by the simi-%
lar work he did within the context of usual topological spaces.  In the
latter,  the finiteness condition required for a map \var{\xout{f}}{}{typewriter deletion}{ps_134-8_014-v0013}
\[ f: X \hand{\longrightarrow} Y \]
of topological spaces to give rise to a functor  $f^{!}$  (\var{\xout{from}}{}{typewriter deletion}{ps_134-8_014-v0014} between\hand{,} say\hand{,}
the derived categories of the categories of abelian sheaves on Y and X)
is \var{\inserted{mainly}}{\hand{mainly}}{insertion}{ps_134-8_014-v0015} that X should be locally embeddable in a product  $Y\mathrm{x}\underline{R}^{d}$  -  a rather
natural condition indeed!   As we are using (Cat) as a kind of algebrai\unsure{c}
% NOTE: the last letter(s) of "algebrai..." run off the right edge of the scan
paradigm for the category of topological spaces,  this last example
for ``sweeping duality''  makes it rather plausible \var{\deleted{indeed}}{}{deletion}{ps_134-8_014-v0016} that something
of the same kind should exist indeed in (Cat)  -  and likewise for the
context of \var{\over{g}{g}}{\hand{g}}{writing over}{ps_134-8_014-v0017}roups  (which we may view as just particular cases of

\margin{14}

% Visual collation: both supplied image halves inspected; see CHECK_DECISIONS_134_7_8.tsv.

%% END SOURCE ps_134-8_014
\ifnum\EditionMode=1
\psapparatusstart
\psentry{ps_134-8_014-v0001}{1}{typewriter deletion}{\xout{\unsure{it}}}{}{The typewriter strikes out this material; it leaves the corrected running text.}
\psentry{ps_134-8_014-v0002}{2}{insertion}{\inserted{two}}{\hand{two}}{The handwritten insertion is incorporated at its indicated position; insertion brackets leave the running text.}
\psentry{ps_134-8_014-v0003}{3}{typewriter deletion}{\xout{these}}{}{The typewriter strikes out this material; it leaves the corrected running text.}
\psentry{ps_134-8_014-v0004}{4}{typewriter deletion}{\xout{operations}}{}{The typewriter strikes out this material; it leaves the corrected running text.}
\psentry{ps_134-8_014-v0005}{5}{typewriter deletion}{\xout{for}}{}{The typewriter strikes out this material; it leaves the corrected running text.}
\psentry{ps_134-8_014-v0006}{6}{insertion}{\inserted{arbitrary}}{\hand{arbitrary}}{The handwritten insertion is incorporated at its indicated position; insertion brackets leave the running text.}
\psentry{ps_134-8_014-v0007}{7}{typed insertion}{\inserted[typed]{most}}{most}{The typed insertion is incorporated at its indicated position; insertion brackets leave the running text.}
\psentry{ps_134-8_014-v0008}{8}{typewriter deletion}{\xout{the}}{}{The typewriter strikes out this material; it leaves the corrected running text.}
\psentry{ps_134-8_014-v0009}{9}{typed insertion}{\inserted[typed]{relationship}}{relationship}{The typed insertion is incorporated at its indicated position; insertion brackets leave the running text.}
\psentry{ps_134-8_014-v0010}{10}{typewriter deletion}{\xout{identity}}{}{The typewriter strikes out this material; it leaves the corrected running text.}
\psentry{ps_134-8_014-v0011}{11}{writing over}{\over{q}{q}}{\hand{q}}{The superposed handwritten reading replaces the underlying reading.}
\psentry{ps_134-8_014-v0012}{12}{insertion}{\inserted{he gave}}{\hand{he gave}}{The handwritten insertion is incorporated at its indicated position; insertion brackets leave the running text.}
\psentry{ps_134-8_014-v0013}{13}{typewriter deletion}{\xout{f}}{}{The typewriter strikes out this material; it leaves the corrected running text.}
\psentry{ps_134-8_014-v0014}{14}{typewriter deletion}{\xout{from}}{}{The typewriter strikes out this material; it leaves the corrected running text.}
\psentry{ps_134-8_014-v0015}{15}{insertion}{\inserted{mainly}}{\hand{mainly}}{The handwritten insertion is incorporated at its indicated position; insertion brackets leave the running text.}
\psentry{ps_134-8_014-v0016}{16}{deletion}{\deleted{indeed}}{}{The author cancels this material by hand; it leaves the corrected running text.}
\psentry{ps_134-8_014-v0017}{17}{writing over}{\over{g}{g}}{\hand{g}}{The superposed handwritten reading replaces the underlying reading.}
\psapparatusend
\fi
\pssource{ps_134-8_015}{134-8/15}{https://grothendieck.umontpellier.fr/web/viewer.html?file=../134-8.pdf\#page=16}
\psnumber{ps_134-8_015-v0001}{1}
\psnumber{ps_134-8_015-v0002}{2}
\psnumber{ps_134-8_015-v0003}{3}
\psnumber{ps_134-8_015-v0004}{4}
\psnumber{ps_134-8_015-v0005}{5}
\psnumber{ps_134-8_015-v0006}{6}
\psnumber{ps_134-8_015-v0007}{7}
\psnumber{ps_134-8_015-v0008}{8}
\psnumber{ps_134-8_015-v0009}{9}
\psnumber{ps_134-8_015-v0010}{10}
\psnumber{ps_134-8_015-v0011}{11}
\psnumber{ps_134-8_015-v0012}{12}
\psnumber{ps_134-8_015-v0013}{13}
\psnumber{ps_134-8_015-v0014}{14}
\psnumber{ps_134-8_015-v0015}{15}
\psnumber{ps_134-8_015-v0016}{16}
\psnumber{ps_134-8_015-v0017}{17}
\psnumber{ps_134-8_015-v0018}{18}
%% BEGIN SOURCE ps_134-8_015
% Pursuing Stacks (cotes 134-2 to 134-8): archive cote 134-8, site page 15 (PDF page 16), typescript page 569.
% typescript. Draft compared: ps7.tex:694-748.
% AI-transcribed by Claude Opus 5.5 (claude-opus-5-5, adaptive thinking, effort high), Message Batches API, batch ps4 (msgbatch_012Aoc1NuskATPsyBeKSYK7J), 2026-11-12; reviewed against both supplied scan halves; see CHECK_DECISIONS.tsv.
\scanpage{134-8/15}
\mspage{569}

models in (Cat)).

I am sorry I was a bit confused, when describing $f^!$ as a right
adjoint to $Rf_*$ - I was thinking of the analogy with a map of schemes,
or spaces, which is not only "of finite type" in a suitable sense,
b\var{\over[typed]{y}{u}}{u}{typewriter overwriting}{ps_134-8_015-v0001}t moreover \uline{proper} - in which case, in those duality contexts,
$Rf_*$ is canonically isomorphic with the functor denoted by $f_!$ or $Rf_!$ .
Otherwise, the ("non trivial")"duality theorem" will assert, rather,
that \var{\xout{if}}{}{typewriter deletion}{ps_134-8_015-v0002} the pair
\[ \var{\over{\unsure{L}}{R}}{\hand{R}}{writing over}{ps_134-8_015-v0003}f_! \ , \ f^! \]
is a pair of adjoint functors, just as \var{\inserted{is}}{\hand{is}}{insertion}{ps_134-8_015-v0004} the pair
\[ f^* \ , \ Rf_* \quad \text{(or simply } f_* \text{ )} \quad . \]
But even when f is assumed to be proper, the sequence (3) is'nt a
sequence of adjoint functors in the stand\var{\over[typed]{\illegible{2}}{ar}}{ar}{typewriter overwriting}{ps_134-8_015-v0005}t duality contexts, namely
$\hand{R}f_!$ is by no means left adjoint to $f^*$ \var{\inserted{i.e. $f^*$ is'nt isomorphic to $f^!$}}{\hand{i.e. $f^*$ is'nt isomorphic to $f^!$}}{insertion}{ps_134-8_015-v0006} (except in extremely special cases,
\var{\xout{such as}}{}{typewriter deletion}{ps_134-8_015-v0007} practically I would think only étale maps are OK , which in
the context of (Cat) would correspond to \var{\xout{functors isomorphic}}{}{typewriter deletion}{ps_134-8_015-v0008} maps in
(Cat) isomorphic to a map $A_{/X} \hand{\longrightarrow} A$ for X in $A^{\hand{\wedge}}$, namely \var{\xout{which}}{}{typewriter deletion}{ps_134-8_015-v0009} maps
which are fibering with discrete fibers). This \uline{does} make an important
discrepancy indeed, between the two kind of contexts - and increases
my perplexity, as to whether or not one should expect a "four variance
duality\var{\xout{"}}{}{typewriter deletion}{ps_134-8_015-v0010} formalism" to make a sense in (Cat). If so, presumably the
$f_!$ or $Rf_!$ it will involve (perhaps via a suitable notion of"proper"
maps in (Cat), as already refer\var{\over{\unsure{e}}{\unsure{r}}}{\hand{\unsure{r}}}{writing over}{ps_134-8_015-v0011}ed to earlier (section 70)) will be
different after all from the $Lf_!$ we have been working with lately,
embodying homology properties. \var{\xout{(}}{}{typewriter deletion}{ps_134-8_015-v0012}But so does $Rf_!$ too , in a rather
strong sense,\var{\xout{(}}{}{typewriter deletion}{ps_134-8_015-v0013} via the\hand{``}duality theorem\hand{''}!\var{\xout{)}}{}{typewriter deletion}{ps_134-8_015-v0014}

In the various duality contexts, a basic part is played by the
three particular classes of maps: proper maps, smooth maps, and im-%
mersions , and factorisation of maps into an immersion followed by
either a proper, or a smooth map. In the context of (Cat), there
is a very natural way indeed to define the three class\var{\inserted[typed]{es}}{es}{typed insertion}{ps_134-8_015-v0015} of maps, as
we'll see in next chapter presumably - so natural indeed, that it is
hard to believe that there may be any other reasonable \var{\over{\illegible{1}}{choice}}{\hand{choice}}{writing over}{ps_134-8_015-v0016}! O$\overset{\illegible{1}}{n}$e very
striking feature (already mentioned in section 70) is that the two \var{\inserted{first}}{\hand{first}}{insertion}{ps_134-8_015-v0017}
notions are "dual" to each other in the rather tautological sense,
namely \var{\inserted{that}}{\hand{that}}{insertion}{ps_134-8_015-v0018} a map $f: A \hand{\longrightarrow} B$ in (Cat) is proper iff the corresponding map
$f^o : A^o \hand{\longrightarrow} B^o$ is smooth (whereas the notion of an immersion is auto-%

\margin{15}

% Visual collation: both supplied image halves inspected; see CHECK_DECISIONS_134_7_8.tsv.

%% END SOURCE ps_134-8_015
\ifnum\EditionMode=1
\psapparatusstart
\psentry{ps_134-8_015-v0001}{1}{typewriter overwriting}{\over[typed]{y}{u}}{u}{The superposed typed reading replaces the underlying reading.}
\psentry{ps_134-8_015-v0002}{2}{typewriter deletion}{\xout{if}}{}{The typewriter strikes out this material; it leaves the corrected running text.}
\psentry{ps_134-8_015-v0003}{3}{writing over}{$\over{\unsure{L}}{R}$}{$\hand{R}$}{The superposed handwritten reading replaces the underlying reading.}
\psentry{ps_134-8_015-v0004}{4}{insertion}{\inserted{is}}{\hand{is}}{The handwritten insertion is incorporated at its indicated position; insertion brackets leave the running text.}
\psentry{ps_134-8_015-v0005}{5}{typewriter overwriting}{\over[typed]{\illegible{2}}{ar}}{ar}{The superposed typed reading replaces the underlying reading.}
\psentry{ps_134-8_015-v0006}{6}{insertion}{\inserted{i.e. $f^*$ is'nt isomorphic to $f^!$}}{\hand{i.e. $f^*$ is'nt isomorphic to $f^!$}}{The handwritten insertion is incorporated at its indicated position; insertion brackets leave the running text.}
\psentry{ps_134-8_015-v0007}{7}{typewriter deletion}{\xout{such as}}{}{The typewriter strikes out this material; it leaves the corrected running text.}
\psentry{ps_134-8_015-v0008}{8}{typewriter deletion}{\xout{functors isomorphic}}{}{The typewriter strikes out this material; it leaves the corrected running text.}
\psentry{ps_134-8_015-v0009}{9}{typewriter deletion}{\xout{which}}{}{The typewriter strikes out this material; it leaves the corrected running text.}
\psentry{ps_134-8_015-v0010}{10}{typewriter deletion}{\xout{"}}{}{The typewriter strikes out this material; it leaves the corrected running text.}
\psentry{ps_134-8_015-v0011}{11}{writing over}{\over{\unsure{e}}{\unsure{r}}}{\hand{\unsure{r}}}{The superposed handwritten reading replaces the underlying reading.}
\psentry{ps_134-8_015-v0012}{12}{typewriter deletion}{\xout{(}}{}{The typewriter strikes out this material; it leaves the corrected running text.}
\psentry{ps_134-8_015-v0013}{13}{typewriter deletion}{\xout{(}}{}{The typewriter strikes out this material; it leaves the corrected running text.}
\psentry{ps_134-8_015-v0014}{14}{typewriter deletion}{\xout{)}}{}{The typewriter strikes out this material; it leaves the corrected running text.}
\psentry{ps_134-8_015-v0015}{15}{typed insertion}{\inserted[typed]{es}}{es}{The typed insertion is incorporated at its indicated position; insertion brackets leave the running text.}
\psentry{ps_134-8_015-v0016}{16}{writing over}{\over{\illegible{1}}{choice}}{\hand{choice}}{The superposed handwritten reading replaces the underlying reading.}
\psentry{ps_134-8_015-v0017}{17}{insertion}{\inserted{first}}{\hand{first}}{The handwritten insertion is incorporated at its indicated position; insertion brackets leave the running text.}
\psentry{ps_134-8_015-v0018}{18}{insertion}{\inserted{that}}{\hand{that}}{The handwritten insertion is incorporated at its indicated position; insertion brackets leave the running text.}
\psapparatusend
\fi
\pssource{ps_134-8_016}{134-8/16}{https://grothendieck.umontpellier.fr/web/viewer.html?file=../134-8.pdf\#page=17}
\psnumber{ps_134-8_016-v0001}{1}
\psnumber{ps_134-8_016-v0002}{2}
\psnumber{ps_134-8_016-v0003}{3}
\psnumber{ps_134-8_016-v0004}{4}
\psnumber{ps_134-8_016-v0005}{5}
\psnumber{ps_134-8_016-v0006}{6}
\psnumber{ps_134-8_016-v0007}{7}
\psnumber{ps_134-8_016-v0008}{8}
\psnumber{ps_134-8_016-v0009}{9}
\psnumber{ps_134-8_016-v0010}{10}
\psnumber{ps_134-8_016-v0011}{11}
\psnumber{ps_134-8_016-v0012}{12}
\psnumber{ps_134-8_016-v0013}{13}
\psnumber{ps_134-8_016-v0014}{14}
\psnumber{ps_134-8_016-v0015}{15}
\psnumber{ps_134-8_016-v0016}{16}
\psnumber{ps_134-8_016-v0017}{17}
\psnumber{ps_134-8_016-v0018}{18}
\psnumber{ps_134-8_016-v0019}{19}
\psnumber{ps_134-8_016-v0020}{20}
%% BEGIN SOURCE ps_134-8_016
% Pursuing Stacks (cotes 134-2 to 134-8): archive cote 134-8, site page 16 (PDF page 17), typescript page 570.
% typescript. Draft compared: ps7.tex:736-797.
% AI-transcribed by Claude Opus 5.5 (claude-opus-5-5, adaptive thinking, effort high), Message Batches API, batch ps4 (msgbatch_012Aoc1NuskATPsyBeKSYK7J), 2026-11-12; reviewed against both supplied scan halves; see CHECK_DECISIONS.tsv.
\scanpage{134-8/16}
\mspage{570}

dual). How does this fit with the expectation of d\var{\over{\unsure{e}}{e}}{\hand{e}}{writing over}{ps_134-8_016-v0001}velopping a ``four
variance'' duality formalism within (Cat) ? It rather heightens perple-%
xity, at first sight ! Proper maps include cofibrations (in the sense
of category theory, not in \var{\xout{in}}{}{typewriter deletion}{ps_134-8_016-v0002} Kan-Quillen's sense !); dually,\var{\xout{s}}{}{typewriter deletion}{ps_134-8_016-v0003} smooth
maps include fibrations\var{\deleted{;}}{}{deletion}{ps_134-8_016-v0004}\hand{.} \var{\over{c}{C}}{\hand{C}}{writing over}{ps_134-8_016-v0005}onsequently, maps which are both smooth and
proper include bifibrations, and hence are not too uncommon. Now, how
strong a restriction is it for a map $f\colon A\hand{\longrightarrow} B$ in(\hand{(}Cat) to factor into
\[ f = p\circ i \qquad , \]
where $i$ is an immersion $A\hand{\longrightarrow} A'$ (namely a functor identifying $A$ to
a full subcategory of $A'$, whose essential image includes, \var{\xout{for}}{}{typewriter deletion}{ps_134-8_016-v0006} \var{\inserted[typed]{with}}{with}{typed insertion}{ps_134-8_016-v0007} any
two objects $x,y$ , any other $z$ which is ``in between'' $x\hand{\to} z\hand{\to} y$ ),
and $p$ is both proper and smooth (a bifibration say) ?

I \var{\xout{feel a little}}{}{typewriter deletion}{ps_134-8_016-v0008} start feeling like a \var{\xout{wari}}{}{typewriter deletion}{ps_134-8_016-v0009} battle horse \var{\xout{\illegible{1}}}{}{typewriter deletion}{ps_134-8_016-v0010}sc\var{\over{\unsure{e}}{e}}{\hand{e}}{writing over}{ps_134-8_016-v0011}nting
gunpowder again - still I do'nt think I'll run into it. Surely there
is something to be cleared up, and perhaps once again a beautiful
duality formalism with the six operations and all will emerge out of
darkness - but this time I will not do the pulling. Maybe someone
else will - if he is'nt discouraged beforehand, because the bigshots
all seem kind of blasé with \var{\xout{dualit}}{}{typewriter deletion}{ps_134-8_016-v0012} ``big duality'', derived categories
and all that. As for my present undertaking, I feel that the question
is'nt really \var{\deleted{a question}}{}{deletion}{ps_134-8_016-v0013} about homotopy models, or about foundations
of homotopy and cohomology formalism - at any rate, that I definitely
do'nt need this kind of stuff, for the program I have been out for.
I should'nt refrain of course to pause on the way every now and then
and have a look at a landscape, however remote \var{\xout{and}}{}{typewriter deletion}{ps_134-8_016-v0014} or misty - but
I am not going to forget I am bound for a journey, and that the jour-%
ney should not be an \var{\deleted{endless}}{}{deletion}{ps_134-8_016-v0015} \var{\inserted{unending}}{\hand{unending}}{insertion}{ps_134-8_016-v0016} one \var{\xout{!}}{}{typewriter deletion}{ps_134-8_016-v0017}..

27.10.

\hand{137.} \var{\xout{Maybe}}{}{typewriter deletion}{ps_134-8_016-v0018} It occured to me that I have been a little rash yesterday,
when asserting that the notions of ``smoothness'' and ``properness'' for
maps in (Cat) which I hit upon last spring is the only ``reasonable''
one. Initially, I referred to these notions by the names ``cohomologi-%
cally smooth'',\var{\xout{\illegible{1}}}{}{typewriter deletion}{ps_134-8_016-v0019} ``cohomologically proper'', as a measure of caution -
they were defined by properties of commutation of base change to
formation of the Leray sheaves $R^if_*$ (i.e. essentially, to\hand{``}cointegra-%
tion\hand{''}\hand{),} which were familiar to me for smooth resp. proper maps in the
context of schemes, or ordinary topological spaces. \var{\deleted{Yesterday}}{}{deletion}{ps_134-8_016-v0020} These

\margin{16}

% Visual collation: both supplied image halves inspected; see CHECK_DECISIONS_134_7_8.tsv.

%% END SOURCE ps_134-8_016
\ifnum\EditionMode=1
\psapparatusstart
\psentry{ps_134-8_016-v0001}{1}{writing over}{\over{\unsure{e}}{e}}{\hand{e}}{The superposed handwritten reading replaces the underlying reading.}
\psentry{ps_134-8_016-v0002}{2}{typewriter deletion}{\xout{in}}{}{The typewriter strikes out this material; it leaves the corrected running text.}
\psentry{ps_134-8_016-v0003}{3}{typewriter deletion}{\xout{s}}{}{The typewriter strikes out this material; it leaves the corrected running text.}
\psentry{ps_134-8_016-v0004}{4}{deletion}{\deleted{;}}{}{The author cancels this material by hand; it leaves the corrected running text.}
\psentry{ps_134-8_016-v0005}{5}{writing over}{\over{c}{C}}{\hand{C}}{The superposed handwritten reading replaces the underlying reading.}
\psentry{ps_134-8_016-v0006}{6}{typewriter deletion}{\xout{for}}{}{The typewriter strikes out this material; it leaves the corrected running text.}
\psentry{ps_134-8_016-v0007}{7}{typed insertion}{\inserted[typed]{with}}{with}{The typed insertion is incorporated at its indicated position; insertion brackets leave the running text.}
\psentry{ps_134-8_016-v0008}{8}{typewriter deletion}{\xout{feel a little}}{}{The typewriter strikes out this material; it leaves the corrected running text.}
\psentry{ps_134-8_016-v0009}{9}{typewriter deletion}{\xout{wari}}{}{The typewriter strikes out this material; it leaves the corrected running text.}
\psentry{ps_134-8_016-v0010}{10}{typewriter deletion}{\xout{\illegible{1}}}{}{The typewriter strikes out this material; it leaves the corrected running text.}
\psentry{ps_134-8_016-v0011}{11}{writing over}{\over{\unsure{e}}{e}}{\hand{e}}{The superposed handwritten reading replaces the underlying reading.}
\psentry{ps_134-8_016-v0012}{12}{typewriter deletion}{\xout{dualit}}{}{The typewriter strikes out this material; it leaves the corrected running text.}
\psentry{ps_134-8_016-v0013}{13}{deletion}{\deleted{a question}}{}{The author cancels this material by hand; it leaves the corrected running text.}
\psentry{ps_134-8_016-v0014}{14}{typewriter deletion}{\xout{and}}{}{The typewriter strikes out this material; it leaves the corrected running text.}
\psentry{ps_134-8_016-v0015}{15}{deletion}{\deleted{endless}}{}{The author cancels this material by hand; it leaves the corrected running text.}
\psentry{ps_134-8_016-v0016}{16}{insertion}{\inserted{unending}}{\hand{unending}}{The handwritten insertion is incorporated at its indicated position; insertion brackets leave the running text.}
\psentry{ps_134-8_016-v0017}{17}{typewriter deletion}{\xout{!}}{}{The typewriter strikes out this material; it leaves the corrected running text.}
\psentry{ps_134-8_016-v0018}{18}{typewriter deletion}{\xout{Maybe}}{}{The typewriter strikes out this material; it leaves the corrected running text.}
\psentry{ps_134-8_016-v0019}{19}{typewriter deletion}{\xout{\illegible{1}}}{}{The typewriter strikes out this material; it leaves the corrected running text.}
\psentry{ps_134-8_016-v0020}{20}{deletion}{\deleted{Yesterday}}{}{The author cancels this material by hand; it leaves the corrected running text.}
\psapparatusend
\fi
\pssource{ps_134-8_017}{134-8/17}{https://grothendieck.umontpellier.fr/web/viewer.html?file=../134-8.pdf\#page=18}
\psnumber{ps_134-8_017-v0001}{1}
\psnumber{ps_134-8_017-v0002}{2}
\psnumber{ps_134-8_017-v0003}{3}
\psnumber{ps_134-8_017-v0004}{4}
\psnumber{ps_134-8_017-v0005}{5}
\psnumber{ps_134-8_017-v0006}{6}
\psnumber{ps_134-8_017-v0007}{7}
\psnumber{ps_134-8_017-v0008}{8}
\psnumber{ps_134-8_017-v0009}{9}
\psnumber{ps_134-8_017-v0010}{10}
\psnumber{ps_134-8_017-v0011}{11}
\psnumber{ps_134-8_017-v0012}{12}
\psnumber{ps_134-8_017-v0013}{13}
\psnumber{ps_134-8_017-v0014}{14}
\psnumber{ps_134-8_017-v0015}{15}
\psnumber{ps_134-8_017-v0016}{16}
\psnumber{ps_134-8_017-v0017}{17}
%% BEGIN SOURCE ps_134-8_017
% Pursuing Stacks (cotes 134-2 to 134-8): archive cote 134-8, site page 17 (PDF page 18), typescript page 571.
% typescript. Draft compared: ps7.tex:785-840.
% AI-transcribed by Claude Opus 5.5 (claude-opus-5-5, adaptive thinking, effort high), Message Batches API, batch ps4 (msgbatch_012Aoc1NuskATPsyBeKSYK7J), 2026-11-12; reviewed against both supplied scan halves; see CHECK_DECISIONS.tsv.
\scanpage{134-8/17}
\mspage{571}

cohomological counterparts of smoothness and properness fi\var{\over[typed]{g}{t}}{t}{typewriter overwriting}{ps_134-8_017-v0001} very neatly
into the homology and cohomology formalism,  and I played around enough
with them, \var{\xout{\unsure{this}}}{}{typewriter deletion}{ps_134-8_017-v0002} \var{\inserted[typed]{last}}{last}{typed insertion}{ps_134-8_017-v0003} spring as well as more than twenty years ago when
developping étale cohomology of schemes, \var{\over{s}{f}}{\hand{f}}{writing over}{ps_134-8_017-v0004}o\hand{r} there \var{\deleted{is indeed}}{}{deletion}{ps_134-8_017-v0005} \var{\inserted{being}}{\hand{being}}{insertion}{ps_134-8_017-v0006} no doubt
left in my mind that these notions are relevant indeed.  However I
\var{\xout{forgot}}{}{typewriter deletion}{ps_134-8_017-v0007} was rather rash yesterday,  while forgetting that the\hand{se} cohomologi%
-cal versions of smoothness and properness are considerably weaker than
the usual notions.  Thus in the context of schemes over a ground field,
the product of any two schemes is cohomologically smooth over it's
factors  -  or equivalently,  any scheme over a field \var{\inserted[typed]{k}}{k}{typed insertion}{ps_134-8_017-v0008} is cohomological-%
ly smooth over k !  Similarly,  in the context of (Cat),  any \var{\xout{category}}{}{typewriter deletion}{ps_134-8_017-v0009}
object in (Cat)\hand{,} namely any small category\hand{,} is both cohomologically
smooth and proper over the final object $e$  (as it is trivially \hand{``}bifi-%
bered\hand{''} over $e$ ).  On the other hand,  it is'nt reasonable\hand{,} of course to
expect any kind of Poincaré-like duality to hold for the cohomology
(with twisted coefficients say) of an arbitrary object A in \hand{(Cat).} To be
more specific,  it is easy to see that in many cases,  the functor
\[ Rf_{*}\ \text{\var{\xout{(\illegible{1}Rf\illegible{1}}}{}{typewriter deletion}{ps_134-8_017-v0010}} : D^{+}(A^{\hand{\wedge}}_{ab}) \hand{\longrightarrow} D^{+}(Ab) \]
does not admit a right adjoint (which we would call $f^{!}$ ).  For instance
when A is discrete,  then $f^{!}$  exists (and may then be identified with
$f^{*}$ ) iff A is moreover \uline{finite}  -  a rather natural condition indeed,
when we keep in mind the topological significance of the usual notion
of properness !  This immediately brings to mind some further proper-%
ties besides base change properties,  which go with the intuitions
around properness  -  for instance we would expect for proper f the
functors $f_{*}$ and $Rf_{*}$ to commute to filtering direct limits, and the
same expectation goes with the assumption that $Rf_{*}$ should admit a
right adjoint.  This exactness property is not satisfied of course
when A is di\var{\xout{\illegible{1}}}{}{typewriter deletion}{ps_134-8_017-v0011}screte infinite. \var{\xout{Another}}{}{typewriter deletion}{ps_134-8_017-v0012}  We \var{\inserted{now}}{\hand{now}}{insertion}{ps_134-8_017-v0013} may think (still in case of
\var{\over{\illegible{1}}{a}}{\hand{a}}{writing over}{ps_134-8_017-v0014} target category equal to $e$) to impose the drastic condition that the
category A is finite.  Such restriction however looks in some respects
too weak,  in others too strong.  Thus it will include categories de-%
fined by finite groups,\var{\deleted{\unsure{f}}}{}{deletion}{ps_134-8_017-v0015} which goes against the rather natural expec-%
tation that properness + smoothness,  or any kind of Poincaré duality,
should go with \uline{finite cohomological dimension}.  On the other hand,
there are beautiful infinite groups (such as the fundamental group\var{\xout{s}}{}{typewriter deletion}{ps_134-8_017-v0016}
of a  com\var{\over[typed]{\illegible{1}}{p}}{p}{typewriter overwriting}{ps_134-8_017-v0017}act surface,  or of any other compact variety which is a
$K(\hand{\pi},1)$ space ...) which satisfy Poincaré duality.

\margin{17}

% Visual collation: both supplied image halves inspected; see CHECK_DECISIONS_134_7_8.tsv.

%% END SOURCE ps_134-8_017
\ifnum\EditionMode=1
\psapparatusstart
\psentry{ps_134-8_017-v0001}{1}{typewriter overwriting}{\over[typed]{g}{t}}{t}{The superposed typed reading replaces the underlying reading.}
\psentry{ps_134-8_017-v0002}{2}{typewriter deletion}{\xout{\unsure{this}}}{}{The typewriter strikes out this material; it leaves the corrected running text.}
\psentry{ps_134-8_017-v0003}{3}{typed insertion}{\inserted[typed]{last}}{last}{The typed insertion is incorporated at its indicated position; insertion brackets leave the running text.}
\psentry{ps_134-8_017-v0004}{4}{writing over}{\over{s}{f}}{\hand{f}}{The superposed handwritten reading replaces the underlying reading.}
\psentry{ps_134-8_017-v0005}{5}{deletion}{\deleted{is indeed}}{}{The author cancels this material by hand; it leaves the corrected running text.}
\psentry{ps_134-8_017-v0006}{6}{insertion}{\inserted{being}}{\hand{being}}{The handwritten insertion is incorporated at its indicated position; insertion brackets leave the running text.}
\psentry{ps_134-8_017-v0007}{7}{typewriter deletion}{\xout{forgot}}{}{The typewriter strikes out this material; it leaves the corrected running text.}
\psentry{ps_134-8_017-v0008}{8}{typed insertion}{\inserted[typed]{k}}{k}{The typed insertion is incorporated at its indicated position; insertion brackets leave the running text.}
\psentry{ps_134-8_017-v0009}{9}{typewriter deletion}{\xout{category}}{}{The typewriter strikes out this material; it leaves the corrected running text.}
\psentry{ps_134-8_017-v0010}{10}{typewriter deletion}{\xout{(\illegible{1}Rf\illegible{1}}}{}{The typewriter strikes out this material; it leaves the corrected running text.}
\psentry{ps_134-8_017-v0011}{11}{typewriter deletion}{\xout{\illegible{1}}}{}{The typewriter strikes out this material; it leaves the corrected running text.}
\psentry{ps_134-8_017-v0012}{12}{typewriter deletion}{\xout{Another}}{}{The typewriter strikes out this material; it leaves the corrected running text.}
\psentry{ps_134-8_017-v0013}{13}{insertion}{\inserted{now}}{\hand{now}}{The handwritten insertion is incorporated at its indicated position; insertion brackets leave the running text.}
\psentry{ps_134-8_017-v0014}{14}{writing over}{\over{\illegible{1}}{a}}{\hand{a}}{The superposed handwritten reading replaces the underlying reading.}
\psentry{ps_134-8_017-v0015}{15}{deletion}{\deleted{\unsure{f}}}{}{The author cancels this material by hand; it leaves the corrected running text.}
\psentry{ps_134-8_017-v0016}{16}{typewriter deletion}{\xout{s}}{}{The typewriter strikes out this material; it leaves the corrected running text.}
\psentry{ps_134-8_017-v0017}{17}{typewriter overwriting}{\over[typed]{\illegible{1}}{p}}{p}{The superposed typed reading replaces the underlying reading.}
\psapparatusend
\fi
\pssource{ps_134-8_018}{134-8/18}{https://grothendieck.umontpellier.fr/web/viewer.html?file=../134-8.pdf\#page=19}
\psnumber{ps_134-8_018-v0001}{1}
\psnumber{ps_134-8_018-v0002}{2}
\psnumber{ps_134-8_018-v0003}{3}
\psnumber{ps_134-8_018-v0004}{4}
\psnumber{ps_134-8_018-v0005}{5}
\psnumber{ps_134-8_018-v0006}{6}
\psnumber{ps_134-8_018-v0007}{7}
\psnumber{ps_134-8_018-v0008}{8}
\psnumber{ps_134-8_018-v0009}{9}
\psnumber{ps_134-8_018-v0010}{10}
\psnumber{ps_134-8_018-v0011}{11}
\psnumber{ps_134-8_018-v0012}{12}
\psnumber{ps_134-8_018-v0013}{13}
\psnumber{ps_134-8_018-v0014}{14}
\psnumber{ps_134-8_018-v0015}{15}
\psnumber{ps_134-8_018-v0016}{16}
\psnumber{ps_134-8_018-v0017}{17}
\psnumber{ps_134-8_018-v0018}{18}
\psnumber{ps_134-8_018-v0019}{19}
\psnumber{ps_134-8_018-v0020}{20}
\psnumber{ps_134-8_018-v0021}{21}
\psnumber{ps_134-8_018-v0022}{22}
\psnumber{ps_134-8_018-v0023}{23}
\psnumber{ps_134-8_018-v0024}{24}
\psnumber{ps_134-8_018-v0025}{25}
\psnumber{ps_134-8_018-v0026}{26}
\psnumber{ps_134-8_018-v0027}{27}
\psnumber{ps_134-8_018-v0028}{28}
\psnumber{ps_134-8_018-v0029}{29}
\psnumber{ps_134-8_018-v0030}{30}
\psnumber{ps_134-8_018-v0031}{31}
\psnumber{ps_134-8_018-v0032}{32}
%% BEGIN SOURCE ps_134-8_018
% Pursuing Stacks (cotes 134-2 to 134-8): archive cote 134-8, site page 18 (PDF page 19), typescript page 572.
% typescript. Draft compared: ps7.tex:825-887.
% AI-transcribed by Claude Opus 5.5 (claude-opus-5-5, adaptive thinking, effort high), Message Batches API, batch ps4 (msgbatch_012Aoc1NuskATPsyBeKSYK7J), 2026-11-12; reviewed against both supplied scan halves; see CHECK_DECISIONS.tsv.
\scanpage{134-8/18}
\mspage{572}

These reflections make it quite clear that the \var{\xout{relevant}}{}{typewriter deletion}{ps_134-8_018-v0001} notions
of properness and \var{\inserted[typed]{of}}{of}{typed insertion}{ps_134-8_018-v0002} smoothness for maps in (Cat), relevant for \var{\deleted{working
out}}{}{deletion}{ps_134-8_018-v0003} a duality formalism, have still to be worked out. Two basic requi-%
rements to be kept in mind are the following: 1) for a proper map \var{\xout{f\illegible{1}}}{}{typewriter deletion}{ps_134-8_018-v0004}
f: A $\hand{\longrightarrow}$ B, and any ring of coefficients k, the functor
\[ Rf_{*} : D^{+}(A^{\wedge}_{k}) \hand{\longrightarrow} D^{+}(B^{\wedge}_{k}) \]
should admit a right adjoint $f^{!}$ , and 2) for a smooth map f factored
as $f = g\circ i$ , with g proper and i an "open immersion" , the compositi-%
% NOTE: the end of this line runs off the right edge of the scan; one character after "compositi" is cut
on \var{\deleted{f}}{}{deletion}{ps_134-8_018-v0005} \var{\inserted{\fbox{$f^{!}\overset{\mathrm{def}}{=}$}}}{\hand{\fbox{$f^{!}\overset{\mathrm{def}}{=}$}}}{insertion}{ps_134-8_018-v0006}$i^{*}\circ g^{!}$\hand{$\leftarrow$} should be \var{\xout{isomorphic to}}{}{typewriter deletion}{ps_134-8_018-v0007} expressible as
\margin{\tikz[baseline=(n.base)]{\node[draw,ellipse,inner sep=2pt] (n) {$f^{!}\overset{\mathrm{def}}{=} i^{*}\circ g^{!}$};}}
\[ (5\unsure{)} \qquad f^{!} : \quad K^{\bullet} \hand{\mapsto} f^{*}(K)\overset{L}{\otimes}_{k} T_{f}(k)\hand{[}d_{f}\hand{]} , \]
where $T_{f}(k)$ is a presheaf of k-modules on A locally isomorphic to
the constant \var{\xout{sheaf}}{}{typewriter deletion}{ps_134-8_018-v0008} presheaf $k_{A}$ ($T_{f}(k)$ may be called the \uline{orientation}
\uline{sheaf} for f, with coefficients in k), and \var{$d_{f}$\hand{|}is}{$d_{f}$ is}{authorial spacing instruction}{ps_134-8_018-v0009} a\var{\xout{n}}{}{typewriter deletion}{ps_134-8_018-v0010} natural integer
(which may be called the \uline{relative dimension} of A over B, or of f).
\hand{(}For simplicity I assume in 2) that A is connected, otherwise $d_{f}$
should be viewed as a function on the set of connected components of
A .\hand{)} This again should give the correct relationship, for f \var{\xout{smooth}}{}{typewriter deletion}{ps_134-8_018-v0011} \var{\inserted[typed]{as above}}{as above}{typed insertion}{ps_134-8_018-v0012}\hand{,}
between the (for the time being hypothetical) $Rf_{!}$ ($\overset{\mathrm{def}}{=} Rg_{*}\circ Li_{!}$ ) and
\var{\xout{the}}{}{typewriter deletion}{ps_134-8_018-v0013} our anodyne $Lf_{!}$ , for an argument L. in $D^{\var{\over{\bullet}{b}}{\hand{b}}{writing over}{ps_134-8_018-v0014}}(A^{\wedge}_{k})$ say :
\[ (6) \qquad Rf_{!}(L.) \hand{\simeq} Lf_{!}(L.\otimes T_{f}^{-1})\hand{[}-d_{f}\hand{]} , \]
where the \var{\xout{first}}{}{typewriter deletion}{ps_134-8_018-v0015} left hand side is just $Rf_{*}(L.)$, if we assume moreover
f to be proper.

This precise relationship between the two possible versions of
an\var{\xout{f}}{}{typewriter deletion}{ps_134-8_018-v0016} $f_{!}$ operation \var{\xout{at th}}{}{typewriter deletion}{ps_134-8_018-v0017} between derived categories, namely $Lf_{!}$
embodying homology, defined for any map f in (Cat), and $Rf_{!}$ embodying
"cohomology with proper supports" , defined for a map which may be
factored as $g\circ i$ with \var{g\hand{/}"proper"and}{g "proper"and}{authorial spacing instruction}{ps_134-8_018-v0018} i an (open, if we like) immersion,
relation valid if f is moreover assumed to be"smooth", greatly relaxes
\var{\xout{the}}{}{typewriter deletion}{ps_134-8_018-v0019} \var{\inserted[typed]{yesterday's}}{yesterday's}{typed insertion}{ps_134-8_018-v0020} perplexity, coming from a \var{\inserted[typed]{partial}}{partial}{typed insertion}{ps_134-8_018-v0021} confusion in my mind between the operati-%
ons $Rf_{!}$ and $Lf_{!}$ . (\var{\deleted{where}}{}{deletion}{ps_134-8_018-v0022} \var{\inserted{Beware}}{\hand{Beware}}{insertion}{ps_134-8_018-v0023} \var{\over{\unsure{T|t}}{t}}{\hand{t}}{writing over}{ps_134-8_018-v0024}he notation $Rf_{!}$ \var{\deleted{being}}{}{deletion}{ps_134-8_018-v0025} \var{\inserted{is}}{\hand{is}}{insertion}{ps_134-8_018-v0026} an abuse,\var{\xout{of not}}{}{typewriter deletion}{ps_134-8_018-v0027} as
\var{\inserted[typed]{\xout{in the present context}}}{}{typed insertion}{ps_134-8_018-v0028} \var{\inserted[typed]{here}}{here}{typed insertion}{ps_134-8_018-v0029} it does'nt mean at all anything like the \var{\xout{derived}}{}{typewriter deletion}{ps_134-8_018-v0030} right derived functor
of the functor $f_{!}$ \hand{!}) At the same time, I feel a lot less dubious now
about the existence of a \var{\xout{duality}}{}{typewriter deletion}{ps_134-8_018-v0031} "six operations" duality formalism in
the (Cat) context - I am pretty much convinced now that such formalism
exists indeed. T\var{\over[typed]{h}{h}}{h}{typewriter overwriting}{ps_134-8_018-v0032}e main specific work ahead is to get hold of the
relevant notions of proper and smooth maps. THe demands we have on

\margin{18}

% Visual collation: both supplied image halves inspected; see CHECK_DECISIONS_134_7_8.tsv.

%% END SOURCE ps_134-8_018
\ifnum\EditionMode=1
\psapparatusstart
\psentry{ps_134-8_018-v0001}{1}{typewriter deletion}{\xout{relevant}}{}{The typewriter strikes out this material; it leaves the corrected running text.}
\psentry{ps_134-8_018-v0002}{2}{typed insertion}{\inserted[typed]{of}}{of}{The typed insertion is incorporated at its indicated position; insertion brackets leave the running text.}
\psentry{ps_134-8_018-v0003}{3}{deletion}{\deleted{working
out}}{}{The author cancels this material by hand; it leaves the corrected running text.}
\psentry{ps_134-8_018-v0004}{4}{typewriter deletion}{\xout{f\illegible{1}}}{}{The typewriter strikes out this material; it leaves the corrected running text.}
\psentry{ps_134-8_018-v0005}{5}{deletion}{\deleted{f}}{}{The author cancels this material by hand; it leaves the corrected running text.}
\psentry{ps_134-8_018-v0006}{6}{insertion}{\inserted{\fbox{$f^{!}\overset{\mathrm{def}}{=}$}}}{\hand{\fbox{$f^{!}\overset{\mathrm{def}}{=}$}}}{The handwritten insertion is incorporated at its indicated position; insertion brackets leave the running text.}
\psentry{ps_134-8_018-v0007}{7}{typewriter deletion}{\xout{isomorphic to}}{}{The typewriter strikes out this material; it leaves the corrected running text.}
\psentry{ps_134-8_018-v0008}{8}{typewriter deletion}{\xout{sheaf}}{}{The typewriter strikes out this material; it leaves the corrected running text.}
\psentry{ps_134-8_018-v0009}{9}{authorial spacing instruction}{$d_{f}$\hand{|}is}{$d_{f}$ is}{The thin vertical pen stroke separates the symbol from is. A reading variant applies that separation.}
\psentry{ps_134-8_018-v0010}{10}{typewriter deletion}{\xout{n}}{}{The typewriter strikes out this material; it leaves the corrected running text.}
\psentry{ps_134-8_018-v0011}{11}{typewriter deletion}{\xout{smooth}}{}{The typewriter strikes out this material; it leaves the corrected running text.}
\psentry{ps_134-8_018-v0012}{12}{typed insertion}{\inserted[typed]{as above}}{as above}{The typed insertion is incorporated at its indicated position; insertion brackets leave the running text.}
\psentry{ps_134-8_018-v0013}{13}{typewriter deletion}{\xout{the}}{}{The typewriter strikes out this material; it leaves the corrected running text.}
\psentry{ps_134-8_018-v0014}{14}{writing over}{$\over{\bullet}{b}$}{$\hand{b}$}{The superposed handwritten reading replaces the underlying reading.}
\psentry{ps_134-8_018-v0015}{15}{typewriter deletion}{\xout{first}}{}{The typewriter strikes out this material; it leaves the corrected running text.}
\psentry{ps_134-8_018-v0016}{16}{typewriter deletion}{\xout{f}}{}{The typewriter strikes out this material; it leaves the corrected running text.}
\psentry{ps_134-8_018-v0017}{17}{typewriter deletion}{\xout{at th}}{}{The typewriter strikes out this material; it leaves the corrected running text.}
\psentry{ps_134-8_018-v0018}{18}{authorial spacing instruction}{g\hand{/}"proper"and}{g "proper"and}{The slanting stroke between g and the opening quotation requests separation. A reading variant replaces it by a space.}
\psentry{ps_134-8_018-v0019}{19}{typewriter deletion}{\xout{the}}{}{The typewriter strikes out this material; it leaves the corrected running text.}
\psentry{ps_134-8_018-v0020}{20}{typed insertion}{\inserted[typed]{yesterday's}}{yesterday's}{The typed insertion is incorporated at its indicated position; insertion brackets leave the running text.}
\psentry{ps_134-8_018-v0021}{21}{typed insertion}{\inserted[typed]{partial}}{partial}{The typed insertion is incorporated at its indicated position; insertion brackets leave the running text.}
\psentry{ps_134-8_018-v0022}{22}{deletion}{\deleted{where}}{}{The author cancels this material by hand; it leaves the corrected running text.}
\psentry{ps_134-8_018-v0023}{23}{insertion}{\inserted{Beware}}{\hand{Beware}}{The handwritten insertion is incorporated at its indicated position; insertion brackets leave the running text.}
\psentry{ps_134-8_018-v0024}{24}{writing over}{\over{\unsure{T|t}}{t}}{\hand{t}}{The superposed handwritten reading replaces the underlying reading.}
\psentry{ps_134-8_018-v0025}{25}{deletion}{\deleted{being}}{}{The author cancels this material by hand; it leaves the corrected running text.}
\psentry{ps_134-8_018-v0026}{26}{insertion}{\inserted{is}}{\hand{is}}{The handwritten insertion is incorporated at its indicated position; insertion brackets leave the running text.}
\psentry{ps_134-8_018-v0027}{27}{typewriter deletion}{\xout{of not}}{}{The typewriter strikes out this material; it leaves the corrected running text.}
\psentry{ps_134-8_018-v0028}{28}{typed insertion}{\inserted[typed]{\xout{in the present context}}}{}{The typed insertion is incorporated at its indicated position; insertion brackets leave the running text.}
\psentry{ps_134-8_018-v0029}{29}{typed insertion}{\inserted[typed]{here}}{here}{The typed insertion is incorporated at its indicated position; insertion brackets leave the running text.}
\psentry{ps_134-8_018-v0030}{30}{typewriter deletion}{\xout{derived}}{}{The typewriter strikes out this material; it leaves the corrected running text.}
\psentry{ps_134-8_018-v0031}{31}{typewriter deletion}{\xout{duality}}{}{The typewriter strikes out this material; it leaves the corrected running text.}
\psentry{ps_134-8_018-v0032}{32}{typewriter overwriting}{\over[typed]{h}{h}}{h}{The superposed typed reading replaces the underlying reading.}
\psapparatusend
\fi
\pssource{ps_134-8_019}{134-8/19}{https://grothendieck.umontpellier.fr/web/viewer.html?file=../134-8.pdf\#page=20}
\psnumber{ps_134-8_019-v0001}{1}
\psnumber{ps_134-8_019-v0002}{2}
\psnumber{ps_134-8_019-v0003}{3}
\psnumber{ps_134-8_019-v0004}{4}
\psnumber{ps_134-8_019-v0005}{5}
\psnumber{ps_134-8_019-v0006}{6}
\psnumber{ps_134-8_019-v0007}{7}
\psnumber{ps_134-8_019-v0008}{8}
\psnumber{ps_134-8_019-v0009}{9}
\psnumber{ps_134-8_019-v0010}{10}
\psnumber{ps_134-8_019-v0011}{11}
\psnumber{ps_134-8_019-v0012}{12}
\psnumber{ps_134-8_019-v0013}{13}
\psnumber{ps_134-8_019-v0014}{14}
\psnumber{ps_134-8_019-v0015}{15}
\psnumber{ps_134-8_019-v0016}{16}
\psnumber{ps_134-8_019-v0017}{17}
\psnumber{ps_134-8_019-v0018}{18}
\psnumber{ps_134-8_019-v0019}{19}
\psnumber{ps_134-8_019-v0020}{20}
\psnumber{ps_134-8_019-v0021}{21}
\psnumber{ps_134-8_019-v0022}{22}
%% BEGIN SOURCE ps_134-8_019
% Pursuing Stacks (cotes 134-2 to 134-8): archive cote 134-8, site page 19 (PDF page 20), typescript page 573.
% typescript. Draft compared: ps7.tex:876-933.
% AI-transcribed by Claude Opus 5.5 (claude-opus-5-5, adaptive thinking, effort high), Message Batches API, batch ps4 (msgbatch_012Aoc1NuskATPsyBeKSYK7J), 2026-11-12; reviewed against both supplied scan halves; see CHECK_DECISIONS.tsv.
\scanpage{134-8/19}
\mspage{573}

these,  besides the relevant base change properties,  are so precise,
one feels,  that they may almost be taken as a definition!  Maybe even
the ``almost'' could be dropped  -  namely that a comprehensive axiomatic
set up for the duality formalism could be worked out,  in a way applicab%
le to the known instances as well as to the pr\var{\over[typed]{\illegible{1}}{e}}{e}{typewriter overwriting}{ps_134-8_019-v0001}sent\var{\xout{\unsure{a}}}{}{typewriter deletion}{ps_134-8_019-v0002}ly still unknown
one of (Cat), \var{\inserted[typed]{by}}{by}{typed insertion}{ps_134-8_019-v0003} going a little further still than Deligne's exposition
in SGA 5 (where the notions of ``smooth'' and ``proper'' maps were supposed
to be given beforehand,  satisfying suitable properties). Before diving
into such axiomatization game,  one should get a better feeling though,
through a fait number of examples (not all w\var{\over{\unsure{i}}{i}}{\hand{i}}{writing over}{ps_134-8_019-v0004}th $e$ as the target category
moreover),  of how the proper,  the smooth and the proper-and-smooth
maps in (Cat) actually look like.  Here presumably Verdier's work in
the context of discrete infinite groups should give useful clues.

Other important clues should come from the opposite side so to say
- namely ordered sets.  Such \var{\inserted[typed]{a}}{a}{typed insertion}{ps_134-8_019-v0005} set\var{\xout{s}}{}{typewriter deletion}{ps_134-8_019-v0006} \var{\inserted[typed]{I}}{I}{typed insertion}{ps_134-8_019-v0007},  besides defining in the usual way
a \var{\inserted[typed]{small}}{small}{typed insertion}{ps_134-8_019-v0008} category and hence a topos, \var{\inserted[typed]{may}}{may}{typed insertion}{ps_134-8_019-v0009} equally \var{\xout{define}}{}{typewriter deletion}{ps_134-8_019-v0010} be viewed as a topological
space,  more accurately the topos it defines may be viewed as being asso%
ciated to a topological space,  admitting I as it's underlying set
(cf. section 22, p.18).  This topological space is noetherian iff the
ordered set I is -  for instance if I is finite.  In such a case, an \var{\xout{al}}{}{typewriter deletion}{ps_134-8_019-v0011}
old algebraic geometer like me will feel in known territory,  which
maybe is a delusion however - at any rate,  I doubt the duality formali-%
sm for topological spaces (using factorizations of maps $X\hand{\longrightarrow}Y$ via embed-%
dings in \var{\over{S}{s}}{\hand{s}}{writing over}{ps_134-8_019-v0012}paces $Y\times\underline{R}^{d}$ ) makes much sense for such non-separated spaces.
However,  as we saw in section 22,  when I satisfies some mild \var{\inserted[typed]{\hand{``}local}}{\hand{``}local}{typed insertion}{ps_134-8_019-v0013} finite-%
ness'' requirement (for instance when I is finite),  we may associate to
it a \uline{geometrical realization} \hand{$|$}I\hand{$|$} ,  which is a locally compact space
(a compact one indeed if I is finite) endowed with a \hand{``}conical \var{\xout{cellular}}{}{typewriter deletion}{ps_134-8_019-v0014}
subdivision\hand{''} \hand{(}indexed by \var{\inserted{the opposite ordered set}}{\hand{the opposite ordered set}}{insertion}{ps_134-8_019-v0015} I$^{\hand{o}}$\hand{)} hence \var{\over[typed]{\illegible{3}}{can}}{can}{typewriter overwriting}{ps_134-8_019-v0016}onically triangulated via the \hand{``}barycen-%
tric subdivision\hand{''}\var{\deleted{)}}{}{deletion}{ps_134-8_019-v0017}.  The homotopy type of this space is canonically iso-%
morphic to the homotopy type of I,  viewed as a ``model'' in (Cat).  What
is more important here,  is that a \var{\inserted[typed]{\hand{(}pre\hand{)}}}{\hand{(}pre\hand{)}}{typed insertion}{ps_134-8_019-v0018}sheaf of sets (say) on the category
I  may be interpreted as being essentially the same as a sheaf of sets
on the geometric realization \hand{$|$}I\hand{$|$} \uline{which is locally constant} (\uline{and} \uline{hence}
\uline{constant}) \uline{on each of the} ``open'' \uline{strata or \hand{``}cones\hand{''} of}  I .  This descripti-%
on then carries over to sheaves of $k$-modules,  The ``clue'' I had in mind
is that \uline{within the context of} \var{\xout{locally}}{}{typewriter deletion}{ps_134-8_019-v0019} \var{\inserted{locally}}{\hand{locally}}{insertion}{ps_134-8_019-v0020} \uline{finite ordered sets}\var{\xout{\illegible{1}}}{}{typewriter deletion}{ps_134-8_019-v0021}\hand{,} \var{\deleted{say,}}{}{deletion}{ps_134-8_019-v0022}  the
\uline{looked for} \uline{``six operations duality formalism''} \uline{should be no more,} \uline{no less}
\uline{than the accurate reflection of the same formalism within the context}

\margin{19}

% Visual collation: both supplied image halves inspected; see CHECK_DECISIONS_134_7_8.tsv.

%% END SOURCE ps_134-8_019
\ifnum\EditionMode=1
\psapparatusstart
\psentry{ps_134-8_019-v0001}{1}{typewriter overwriting}{\over[typed]{\illegible{1}}{e}}{e}{The superposed typed reading replaces the underlying reading.}
\psentry{ps_134-8_019-v0002}{2}{typewriter deletion}{\xout{\unsure{a}}}{}{The typewriter strikes out this material; it leaves the corrected running text.}
\psentry{ps_134-8_019-v0003}{3}{typed insertion}{\inserted[typed]{by}}{by}{The typed insertion is incorporated at its indicated position; insertion brackets leave the running text.}
\psentry{ps_134-8_019-v0004}{4}{writing over}{\over{\unsure{i}}{i}}{\hand{i}}{The superposed handwritten reading replaces the underlying reading.}
\psentry{ps_134-8_019-v0005}{5}{typed insertion}{\inserted[typed]{a}}{a}{The typed insertion is incorporated at its indicated position; insertion brackets leave the running text.}
\psentry{ps_134-8_019-v0006}{6}{typewriter deletion}{\xout{s}}{}{The typewriter strikes out this material; it leaves the corrected running text.}
\psentry{ps_134-8_019-v0007}{7}{typed insertion}{\inserted[typed]{I}}{I}{The typed insertion is incorporated at its indicated position; insertion brackets leave the running text.}
\psentry{ps_134-8_019-v0008}{8}{typed insertion}{\inserted[typed]{small}}{small}{The typed insertion is incorporated at its indicated position; insertion brackets leave the running text.}
\psentry{ps_134-8_019-v0009}{9}{typed insertion}{\inserted[typed]{may}}{may}{The typed insertion is incorporated at its indicated position; insertion brackets leave the running text.}
\psentry{ps_134-8_019-v0010}{10}{typewriter deletion}{\xout{define}}{}{The typewriter strikes out this material; it leaves the corrected running text.}
\psentry{ps_134-8_019-v0011}{11}{typewriter deletion}{\xout{al}}{}{The typewriter strikes out this material; it leaves the corrected running text.}
\psentry{ps_134-8_019-v0012}{12}{writing over}{\over{S}{s}}{\hand{s}}{The superposed handwritten reading replaces the underlying reading.}
\psentry{ps_134-8_019-v0013}{13}{typed insertion}{\inserted[typed]{\hand{``}local}}{\hand{``}local}{The typed insertion is incorporated at its indicated position; insertion brackets leave the running text.}
\psentry{ps_134-8_019-v0014}{14}{typewriter deletion}{\xout{cellular}}{}{The typewriter strikes out this material; it leaves the corrected running text.}
\psentry{ps_134-8_019-v0015}{15}{insertion}{\inserted{the opposite ordered set}}{\hand{the opposite ordered set}}{The handwritten insertion is incorporated at its indicated position; insertion brackets leave the running text.}
\psentry{ps_134-8_019-v0016}{16}{typewriter overwriting}{\over[typed]{\illegible{3}}{can}}{can}{The superposed typed reading replaces the underlying reading.}
\psentry{ps_134-8_019-v0017}{17}{deletion}{\deleted{)}}{}{The author cancels this material by hand; it leaves the corrected running text.}
\psentry{ps_134-8_019-v0018}{18}{typed insertion}{\inserted[typed]{\hand{(}pre\hand{)}}}{\hand{(}pre\hand{)}}{The typed insertion is incorporated at its indicated position; insertion brackets leave the running text.}
\psentry{ps_134-8_019-v0019}{19}{typewriter deletion}{\xout{locally}}{}{The typewriter strikes out this material; it leaves the corrected running text.}
\psentry{ps_134-8_019-v0020}{20}{insertion}{\inserted{locally}}{\hand{locally}}{The handwritten insertion is incorporated at its indicated position; insertion brackets leave the running text.}
\psentry{ps_134-8_019-v0021}{21}{typewriter deletion}{\xout{\illegible{1}}}{}{The typewriter strikes out this material; it leaves the corrected running text.}
\psentry{ps_134-8_019-v0022}{22}{deletion}{\deleted{say,}}{}{The author cancels this material by hand; it leaves the corrected running text.}
\psapparatusend
\fi
\pssource{ps_134-8_020}{134-8/20}{https://grothendieck.umontpellier.fr/web/viewer.html?file=../134-8.pdf\#page=21}
\psnumber{ps_134-8_020-v0001}{1}
\psnumber{ps_134-8_020-v0002}{2}
\psnumber{ps_134-8_020-v0003}{3}
\psnumber{ps_134-8_020-v0004}{4}
\psnumber{ps_134-8_020-v0005}{5}
\psnumber{ps_134-8_020-v0006}{6}
\psnumber{ps_134-8_020-v0007}{7}
\psnumber{ps_134-8_020-v0008}{8}
\psnumber{ps_134-8_020-v0009}{9}
\psnumber{ps_134-8_020-v0010}{10}
%% BEGIN SOURCE ps_134-8_020
% Pursuing Stacks (cotes 134-2 to 134-8): archive cote 134-8, site page 20 (PDF page 21), typescript page 574.
% typescript. Draft compared: ps7.tex:921-985.
% AI-transcribed by Claude Opus 5.5 (claude-opus-5-5, adaptive thinking, effort high), Message Batches API, batch ps4 (msgbatch_012Aoc1NuskATPsyBeKSYK7J), 2026-11-12; reviewed against both supplied scan halves; see CHECK_DECISIONS.tsv.
\scanpage{134-8/20}
\mspage{574}

\uline{of} (\uline{locally compac}t) \uline{topological spaces}, as worked out by Verdier in one
of his Bourbaki talks - it being understood that when applying the
latter formalism to spaces endowed with conical stratifications, maps
between these which are compatible with the stratifications (in a suitab%
% NOTE: the right edge of the scan cuts off the ends of some lines (after "suitab", "satis", "one"); a hyphen or letter may be hidden there
le sense which should still be pinned down), and to sheaves of modules
which \var{\xout{satisfy the}}{}{typewriter deletion}{ps_134-8_020-v0001} are compatible with the stratifications in the sense
above, these will give rise (via the ``six operations'') to sheaves satis%
\hand{-}fying the same compatibility. (NB when speaking of``sheaves'', I really
mean complexes of sheaves, \var{\inserted[typed]{$K^\bullet$}}{$K^\bullet$}{typed insertion}{ps_134-8_020-v0002} and the compatibility condition should be
understood as a condition \var{\xout{at}}{}{typewriter deletion}{ps_134-8_020-v0003} \var{\inserted[typed]{for}}{for}{typed insertion}{ps_134-8_020-v0004} the ordinary sheaves of modules $\underline{H}^i(K^\bullet)$.)
This remark should allow to work out quite explicitly, in purely alge-%
braic (or ``combinatorial'') terms, the ``six operations'' in the context
of locally finite ordered sets, at any rate.

This interpretation suggests that an ordered set I should be viewed
as a ``proper'' object of (Cat) iff I is finite. In the same vein, when-%
ever I is ``locally proper'' namely locally finite, and moreover it's to-%
pological realization I is a ``smooth'' topological space in the usual
sense, namely is a topological \uline{variety} (for which it is enough that
for every x in I, the topological realization
\hand{(7)}
\[ \hand{|}I_{\hand{>}x}\hand{|} \]
of the set of elements y such that y \hand{$>$} x should be a sphere), we would
consider I as a ``smooth'' object in (Cat). If A is any object in (Cat)
and I is an ordered set which is finite resp. ``smooth'' in the sense
above, we will surely expect AxI \hand{$\longrightarrow$} A to be``proper'' resp. ``smooth''
for the duality formalism we wish to develop in (Cat).

It should be kept in mind that for the \var{\inserted[typed]{algebraic}}{algebraic}{typed insertion}{ps_134-8_020-v0005} interpretation above \var{\inserted[typed]{to hold}}{to hold}{typed insertion}{ps_134-8_020-v0006}, for
sheaves on a conically stratified space \hand{$|$}I\hand{$|$} in terms of an ordered set
I , we had to take on the indexing set \var{\inserted[typed]{I}}{I}{typed insertion}{ps_134-8_020-v0007} for the strata the order relati-%
on \uline{opposite} to the inclusion relation between (closed) strata - other-%
wise the correct interpretation of sheaves constant on the open strata
is via \uline{covariant} functors I \hand{$\longrightarrow$} (Sets), i.e. \uline{presheaves on} $I^o$ (not I).
At any rate, as there is a canonical homeomorphism
\hand{(8)}
\[ \hand{|}I\hand{|} \hand{\simeq} \hand{|}I^o\hand{|} \]
respecting \var{\inserted[typed]{the canonical}}{the canonical}{typed insertion}{ps_134-8_020-v0008} barycentric subdivisions of both sides, notions \var{\xout{\unsure{f}}}{}{typewriter deletion}{ps_134-8_020-v0009}for I
(such as properness, or smoothness) which are expressed as intrinsic \var{\xout{p}}{}{typewriter deletion}{ps_134-8_020-v0010}
properties of the corresponding topological space \hand{$|$}I\hand{$|$} (independently of
it's subdivision) are autodual - they hold for I iff they do for $\hand{I^o}$.

\margin{20}

% Visual collation: both supplied image halves inspected; see CHECK_DECISIONS_134_7_8.tsv.

%% END SOURCE ps_134-8_020
\ifnum\EditionMode=1
\psapparatusstart
\psentry{ps_134-8_020-v0001}{1}{typewriter deletion}{\xout{satisfy the}}{}{The typewriter strikes out this material; it leaves the corrected running text.}
\psentry{ps_134-8_020-v0002}{2}{typed insertion}{\inserted[typed]{$K^\bullet$}}{$K^\bullet$}{The typed insertion is incorporated at its indicated position; insertion brackets leave the running text.}
\psentry{ps_134-8_020-v0003}{3}{typewriter deletion}{\xout{at}}{}{The typewriter strikes out this material; it leaves the corrected running text.}
\psentry{ps_134-8_020-v0004}{4}{typed insertion}{\inserted[typed]{for}}{for}{The typed insertion is incorporated at its indicated position; insertion brackets leave the running text.}
\psentry{ps_134-8_020-v0005}{5}{typed insertion}{\inserted[typed]{algebraic}}{algebraic}{The typed insertion is incorporated at its indicated position; insertion brackets leave the running text.}
\psentry{ps_134-8_020-v0006}{6}{typed insertion}{\inserted[typed]{to hold}}{to hold}{The typed insertion is incorporated at its indicated position; insertion brackets leave the running text.}
\psentry{ps_134-8_020-v0007}{7}{typed insertion}{\inserted[typed]{I}}{I}{The typed insertion is incorporated at its indicated position; insertion brackets leave the running text.}
\psentry{ps_134-8_020-v0008}{8}{typed insertion}{\inserted[typed]{the canonical}}{the canonical}{The typed insertion is incorporated at its indicated position; insertion brackets leave the running text.}
\psentry{ps_134-8_020-v0009}{9}{typewriter deletion}{\xout{\unsure{f}}}{}{The typewriter strikes out this material; it leaves the corrected running text.}
\psentry{ps_134-8_020-v0010}{10}{typewriter deletion}{\xout{p}}{}{The typewriter strikes out this material; it leaves the corrected running text.}
\psapparatusend
\fi
\pssource{ps_134-8_021}{134-8/21}{https://grothendieck.umontpellier.fr/web/viewer.html?file=../134-8.pdf\#page=22}
\psnumber{ps_134-8_021-v0001}{1}
\psnumber{ps_134-8_021-v0002}{2}
\psnumber{ps_134-8_021-v0003}{3}
\psnumber{ps_134-8_021-v0004}{4}
\psnumber{ps_134-8_021-v0005}{5}
\psnumber{ps_134-8_021-v0006}{6}
\psnumber{ps_134-8_021-v0007}{7}
\psnumber{ps_134-8_021-v0008}{8}
\psnumber{ps_134-8_021-v0009}{9}
\psnumber{ps_134-8_021-v0010}{10}
\psnumber{ps_134-8_021-v0011}{11}
\psnumber{ps_134-8_021-v0012}{12}
\psnumber{ps_134-8_021-v0013}{13}
\psnumber{ps_134-8_021-v0014}{14}
\psnumber{ps_134-8_021-v0015}{15}
%% BEGIN SOURCE ps_134-8_021
% Pursuing Stacks (cotes 134-2 to 134-8): archive cote 134-8, site page 21 (PDF page 22), typescript page 575.
% typescript. Draft compared: ps7.tex:968-1030.
% AI-transcribed by Claude Opus 5.5 (claude-opus-5-5, adaptive thinking, effort high), Message Batches API, batch ps4 (msgbatch_012Aoc1NuskATPsyBeKSYK7J), 2026-11-12; reviewed against both supplied scan halves; see CHECK_DECISIONS.tsv.
\scanpage{134-8/21}
\mspage{575}

This is in sharp contrast with the more naïve notions of ``cohomological''
properness and smoothness via base change properties, which are inter-%
changed by duality (which was part of yesterday's perplexities,\var{\over{n}{\unsure{N|n}}}{\hand{\unsure{N|n}}}{writing over}{ps_134-8_021-v0001}ow
straightening out ...).

I feel I should be a little more outspoken about the relevant
notion of ``proper maps'' between ordered sets, which should be the alge-%
braic expression of the geometric notion alluded to above\hand{,} of a map
between conically stratified spaces being ``compatible with the conical
stratifications''. In order for the corresponding direct image functor
for sheaves to take sheaves compatible with the stratification above to
sheaves of s ame type below, we'll have to insist that the image by f
of a strata above should be a whole strata below. Now it is clear that
any map
\[ (9) \qquad f: I \hand{\longrightarrow} J \]
between ordered sets \var{\xout{defines}}{}{typewriter deletion}{ps_134-8_021-v0002} takes flags into flags, and hence induces
a map
\[ (10) \qquad \hand{|}f\hand{|}:\hand{|}I\hand{|}\hand{\to}\hand{|}J\hand{|} \]
between the geometric realizations, compatible with the barycentric
triangulations and \var{\deleted{hence}}{}{deletion}{ps_134-8_021-v0003} \var{\inserted{thus}}{\hand{thus}}{insertion}{ps_134-8_021-v0004} taking simplices into simplices. But even
when I and J are finite (hence ``proper''), the latter map does not al-%
ways satisfy the condition above\var{\over{,}{.}}{\hand{.}}{writing over}{ps_134-8_021-v0005} \var{\xout{even when}}{}{typewriter deletion}{ps_134-8_021-v0006} Thus when I is reduced to
just one point, hence (9) is defined \var{\xout{\illegible{1}i}}{}{typewriter deletion}{ps_134-8_021-v0007} by the image $\var{\over{\illegible{1}}{y}}{\hand{y}}{writing over}{ps_134-8_021-v0008}\hand{\in} J$ of the lat-%
ter, the corresponding map (10) maps the unique point of $\hand{|}I\hand{|}$ to the
barycenter of the strat\var{\over{a}{um}}{\hand{um}}{writing over}{ps_134-8_021-v0009} $C_{\var{\over{\illegible{1}}{y}}{\hand{y}}{writing over}{ps_134-8_021-v0010}}$ in $\hand{|}J\hand{|}$ , which is a strat\var{\over{a}{um}}{\hand{um}}{writing over}{ps_134-8_021-v0011} of $\hand{|}J\hand{|}$ iff
\var{\xout{\illegible{1} is minim}}{}{typewriter deletion}{ps_134-8_021-v0012} $C_{\var{\over{\illegible{1}}{y}}{\hand{y}}{writing over}{ps_134-8_021-v0013}}$ is a minimal strata, i.e. $\var{\over{\illegible{1}}{y}}{\hand{y}}{writing over}{ps_134-8_021-v0014}$ is a \uline{maximal} element in J.
A natural algebraic condition to impose upon f, in order to insure
that $\hand{|}f\hand{|}$ satisfy the simple geometric condition above, is that for
any x in I, and any y' in J such that
\[ y' \hand{>} y \overset{\text{def}}{=} f(x) \qquad , \]
there should exist an x' in I satisfying
\[ \hand{(11)} \qquad x' \hand{>} x \quad \text{and} \quad f(x') = y' \quad . \]
This condition strongly reminds of the valuative criterion for proper-%
ness in the context of schemes, where the relation $y' \hand{\geq} y$ \var{\inserted[typed]{or $y\hand{\to} y'$}}{or $y\hand{\to} y'$}{typed insertion}{ps_134-8_021-v0015}\hand{,} say\hand{,} should
be interpreted as meaning that y' is a \uline{specialization} of \hand{$y$} . However,
in the valuative criterion for properness (for a map of preschemes of
finite type over a noetherian prescheme say), if one wants actual pro\unsure{p}-%

\margin{21}

% Visual collation: both supplied image halves inspected; see CHECK_DECISIONS_134_7_8.tsv.

%% END SOURCE ps_134-8_021
\ifnum\EditionMode=1
\psapparatusstart
\psentry{ps_134-8_021-v0001}{1}{writing over}{\over{n}{\unsure{N|n}}}{\hand{\unsure{N|n}}}{The superposed handwritten reading replaces the underlying reading.}
\psentry{ps_134-8_021-v0002}{2}{typewriter deletion}{\xout{defines}}{}{The typewriter strikes out this material; it leaves the corrected running text.}
\psentry{ps_134-8_021-v0003}{3}{deletion}{\deleted{hence}}{}{The author cancels this material by hand; it leaves the corrected running text.}
\psentry{ps_134-8_021-v0004}{4}{insertion}{\inserted{thus}}{\hand{thus}}{The handwritten insertion is incorporated at its indicated position; insertion brackets leave the running text.}
\psentry{ps_134-8_021-v0005}{5}{writing over}{\over{,}{.}}{\hand{.}}{The superposed handwritten reading replaces the underlying reading.}
\psentry{ps_134-8_021-v0006}{6}{typewriter deletion}{\xout{even when}}{}{The typewriter strikes out this material; it leaves the corrected running text.}
\psentry{ps_134-8_021-v0007}{7}{typewriter deletion}{\xout{\illegible{1}i}}{}{The typewriter strikes out this material; it leaves the corrected running text.}
\psentry{ps_134-8_021-v0008}{8}{writing over}{$\over{\illegible{1}}{y}$}{$\hand{y}$}{The superposed handwritten reading replaces the underlying reading.}
\psentry{ps_134-8_021-v0009}{9}{writing over}{\over{a}{um}}{\hand{um}}{The superposed handwritten reading replaces the underlying reading.}
\psentry{ps_134-8_021-v0010}{10}{writing over}{$\over{\illegible{1}}{y}$}{$\hand{y}$}{The superposed handwritten reading replaces the underlying reading.}
\psentry{ps_134-8_021-v0011}{11}{writing over}{\over{a}{um}}{\hand{um}}{The superposed handwritten reading replaces the underlying reading.}
\psentry{ps_134-8_021-v0012}{12}{typewriter deletion}{\xout{\illegible{1} is minim}}{}{The typewriter strikes out this material; it leaves the corrected running text.}
\psentry{ps_134-8_021-v0013}{13}{writing over}{$\over{\illegible{1}}{y}$}{$\hand{y}$}{The superposed handwritten reading replaces the underlying reading.}
\psentry{ps_134-8_021-v0014}{14}{writing over}{$\over{\illegible{1}}{y}$}{$\hand{y}$}{The superposed handwritten reading replaces the underlying reading.}
\psentry{ps_134-8_021-v0015}{15}{typed insertion}{\inserted[typed]{or $y\hand{\to} y'$}}{or $y\hand{\to} y'$}{The typed insertion is incorporated at its indicated position; insertion brackets leave the running text.}
\psapparatusend
\fi
\pssource{ps_134-8_022}{134-8/22}{https://grothendieck.umontpellier.fr/web/viewer.html?file=../134-8.pdf\#page=23}
\psnumber{ps_134-8_022-v0001}{1}
\psnumber{ps_134-8_022-v0002}{2}
\psnumber{ps_134-8_022-v0003}{3}
\psnumber{ps_134-8_022-v0004}{4}
\psnumber{ps_134-8_022-v0005}{5}
\psnumber{ps_134-8_022-v0006}{6}
\psnumber{ps_134-8_022-v0007}{7}
\psnumber{ps_134-8_022-v0008}{8}
\psnumber{ps_134-8_022-v0009}{9}
\psnumber{ps_134-8_022-v0010}{10}
\psnumber{ps_134-8_022-v0011}{11}
\psnumber{ps_134-8_022-v0012}{12}
\psnumber{ps_134-8_022-v0013}{13}
\psnumber{ps_134-8_022-v0014}{14}
\psnumber{ps_134-8_022-v0015}{15}
\psnumber{ps_134-8_022-v0016}{16}
\psnumber{ps_134-8_022-v0017}{17}
\psnumber{ps_134-8_022-v0018}{18}
%% BEGIN SOURCE ps_134-8_022
% Pursuing Stacks (cotes 134-2 to 134-8): archive cote 134-8, site page 22 (PDF page 23), typescript page 576.
% typescript. Draft compared: ps7.tex:1018-1080.
% AI-transcribed by Claude Opus 5.5 (claude-opus-5-5, adaptive thinking, effort high), Message Batches API, batch ps4 (msgbatch_012Aoc1NuskATPsyBeKSYK7J), 2026-11-12; reviewed against both supplied scan halves; see CHECK_DECISIONS.tsv.
\scanpage{134-8/22}
\mspage{576}

perness indeed (including separation of f,  not just that f is universal%
ly closed),  one has to insist that the x' above should be unique:
every specialization y' of y= f(x)  lifts to a \uline{unique} specialization
x' of x.  If we applied this literally in the present context,  this
would translate into the condition that \hand{$|$}f\hand{$|$}  should map \uline{injectively}
each closed strat\var{\over{a}{um}}{\hand{um}}{writing over}{ps_134-8_022-v0001} $C_x$ of \hand{$|$}I\hand{$|$} -  which would exclude such basic maps as
the projection I \hand{$\longrightarrow$} e to just one point !

O\var{\over[typed]{n}{n}}{n}{typewriter overwriting}{ps_134-8_022-v0002}e may wonder why trouble about the analogy with algebraic geo-%
metry and any extra condition on f besddes the one we got.  T\var{\over[typed]{h}{h}}{h}{typewriter overwriting}{ps_134-8_022-v0003}e point
however is that we would like the map
\[ \text{(12)} \qquad C_x = \hand{|} I_{\hand{\geqslant}x} \hand{|} \hand{\longrightarrow} \ C_y = \hand{|} I_{\hand{\geqslant}y} \hand{|} \]
between corresponding strata induced by \hand{$|$}f\hand{$|$} to be ``cohomologically
trivial'' in a \uline{suitable sense},  not only surjective  -  in analogy say
\margin{namely: aspheric fibers!}with the usual notion of maps between triangulated spaces;  if we don't
have some condition of this type,  we will have no control over the
structure of the direct image and the higher direct images $R^i\hand{|}f\hand{|}_*$
of an ``admissible'' sheaf upstairs.  I did'nt really analyze the situati-%
on carefully,  in terms of what we are after here in the context of
those stratified geometric realizations.  I feel pretty sure though \var{\xout{tha}}{}{typewriter deletion}{ps_134-8_022-v0004}
that\var{\xout{\unsure{t}}}{}{typewriter deletion}{ps_134-8_022-v0005} here the general \unsure{motion|notion} of ``cohomological properness'' which I
worked out last spring fits in just right.,  to give the correct answer.
T\var{\over[typed]{h}{h}}{h}{typewriter overwriting}{ps_134-8_022-v0006}e crit\var{\over[typed]{e}{e}}{e}{typewriter overwriting}{ps_134-8_022-v0007}rion I obtained (necessary \var{\inserted[typed]{and sufficient}}{and sufficient}{typed insertion}{ps_134-8_022-v0008} for the relevant \var{\xout{base change}}{}{typewriter deletion}{ps_134-8_022-v0009} \var{\inserted[typed]{compatibility}}{compatibility}{typed insertion}{ps_134-8_022-v0010} proper%
ty of $R^if_*$  with arbitrary base change J' \hand{$\to$} J in (Cat)) reads as fol-%
lows:  let
\[ \text{(13)} \qquad I(f,x,y') = \hand{\{}\, x'\hand{\in} I \hand{|} \ x' \hand{>} x \ \text{ and } f(x') = y' \hand{\}} \]
be the subset of I satisfying the conditions (11) above.  I\var{\over[typed]{n}{n}}{n}{typewriter overwriting}{ps_134-8_022-v0011}stead
of demanding only that this set be non empty,  or going as far as
demanding that it should \var{\xout{have}}{}{typewriter deletion}{ps_134-8_022-v0012} be reduced to just one point,  we'll
demand that this set should be \uline{aspheric}.  (When we are concerned with
cohomology with commutative coefficients ohly,  presumab\var{\over[typed]{l}{l}}{l}{typewriter overwriting}{ps_134-8_022-v0013}y it should
be enough to demand only that this set be \uline{acyclic} - which should be
enough for the sake of \var{\inserted[typed]{a}}{a}{typed insertion}{ps_134-8_022-v0014} mere duality formalism ...).  In more geometric
terms,  this should mean,  I guess,  tha\var{\over[typed]{t}{t}}{t}{typewriter overwriting}{ps_134-8_022-v0015} the inverse image\hand{,} by the
\margin{$\times$}map (12) above between closed strata\hand{,} of any closed strat\var{\over{a}{um}}{\hand{um}}{writing over}{ps_134-8_022-v0016} below,
\margin{\unsure{,}}\unsure{shoubd|should} be aspheric.  I doubt this condition holds under the  mere \var{\xout{assu}}{}{typewriter deletion}{ps_134-8_022-v0017}
assumption that the sets (13) be non empty,  i.e. (12) be surjecti\var{\over[typed]{\unsure{f}}{v}}{v}{typewriter overwriting}{ps_134-8_022-v0018}e
-  I did'nt sit down though to try and make an example.

\margin{22}

% Visual collation: both supplied image halves inspected; see CHECK_DECISIONS_134_7_8.tsv.

%% END SOURCE ps_134-8_022
\ifnum\EditionMode=1
\psapparatusstart
\psentry{ps_134-8_022-v0001}{1}{writing over}{\over{a}{um}}{\hand{um}}{The superposed handwritten reading replaces the underlying reading.}
\psentry{ps_134-8_022-v0002}{2}{typewriter overwriting}{\over[typed]{n}{n}}{n}{The superposed typed reading replaces the underlying reading.}
\psentry{ps_134-8_022-v0003}{3}{typewriter overwriting}{\over[typed]{h}{h}}{h}{The superposed typed reading replaces the underlying reading.}
\psentry{ps_134-8_022-v0004}{4}{typewriter deletion}{\xout{tha}}{}{The typewriter strikes out this material; it leaves the corrected running text.}
\psentry{ps_134-8_022-v0005}{5}{typewriter deletion}{\xout{\unsure{t}}}{}{The typewriter strikes out this material; it leaves the corrected running text.}
\psentry{ps_134-8_022-v0006}{6}{typewriter overwriting}{\over[typed]{h}{h}}{h}{The superposed typed reading replaces the underlying reading.}
\psentry{ps_134-8_022-v0007}{7}{typewriter overwriting}{\over[typed]{e}{e}}{e}{The superposed typed reading replaces the underlying reading.}
\psentry{ps_134-8_022-v0008}{8}{typed insertion}{\inserted[typed]{and sufficient}}{and sufficient}{The typed insertion is incorporated at its indicated position; insertion brackets leave the running text.}
\psentry{ps_134-8_022-v0009}{9}{typewriter deletion}{\xout{base change}}{}{The typewriter strikes out this material; it leaves the corrected running text.}
\psentry{ps_134-8_022-v0010}{10}{typed insertion}{\inserted[typed]{compatibility}}{compatibility}{The typed insertion is incorporated at its indicated position; insertion brackets leave the running text.}
\psentry{ps_134-8_022-v0011}{11}{typewriter overwriting}{\over[typed]{n}{n}}{n}{The superposed typed reading replaces the underlying reading.}
\psentry{ps_134-8_022-v0012}{12}{typewriter deletion}{\xout{have}}{}{The typewriter strikes out this material; it leaves the corrected running text.}
\psentry{ps_134-8_022-v0013}{13}{typewriter overwriting}{\over[typed]{l}{l}}{l}{The superposed typed reading replaces the underlying reading.}
\psentry{ps_134-8_022-v0014}{14}{typed insertion}{\inserted[typed]{a}}{a}{The typed insertion is incorporated at its indicated position; insertion brackets leave the running text.}
\psentry{ps_134-8_022-v0015}{15}{typewriter overwriting}{\over[typed]{t}{t}}{t}{The superposed typed reading replaces the underlying reading.}
\psentry{ps_134-8_022-v0016}{16}{writing over}{\over{a}{um}}{\hand{um}}{The superposed handwritten reading replaces the underlying reading.}
\psentry{ps_134-8_022-v0017}{17}{typewriter deletion}{\xout{assu}}{}{The typewriter strikes out this material; it leaves the corrected running text.}
\psentry{ps_134-8_022-v0018}{18}{typewriter overwriting}{\over[typed]{\unsure{f}}{v}}{v}{The superposed typed reading replaces the underlying reading.}
\psapparatusend
\fi
\pssource{ps_134-8_023}{134-8/23}{https://grothendieck.umontpellier.fr/web/viewer.html?file=../134-8.pdf\#page=24}
\psnumber{ps_134-8_023-v0001}{1}
\psnumber{ps_134-8_023-v0002}{2}
\psnumber{ps_134-8_023-v0003}{3}
\psnumber{ps_134-8_023-v0004}{4}
\psnumber{ps_134-8_023-v0005}{5}
\psnumber{ps_134-8_023-v0006}{6}
\psnumber{ps_134-8_023-v0007}{7}
\psnumber{ps_134-8_023-v0008}{8}
\psnumber{ps_134-8_023-v0009}{9}
\psnumber{ps_134-8_023-v0010}{10}
\psnumber{ps_134-8_023-v0011}{11}
\psnumber{ps_134-8_023-v0012}{12}
\psnumber{ps_134-8_023-v0013}{13}
\psnumber{ps_134-8_023-v0014}{14}
\psnumber{ps_134-8_023-v0015}{15}
\psnumber{ps_134-8_023-v0016}{16}
\psnumber{ps_134-8_023-v0017}{17}
\psnumber{ps_134-8_023-v0018}{18}
\psnumber{ps_134-8_023-v0019}{19}
%% BEGIN SOURCE ps_134-8_023
% Pursuing Stacks (cotes 134-2 to 134-8): archive cote 134-8, site page 23 (PDF page 24), typescript page 577.
% typescript. Draft compared: ps7.tex:1068-1126.
% AI-transcribed by Claude Opus 5.5 (claude-opus-5-5, adaptive thinking, effort high), Message Batches API, batch ps4 (msgbatch_012Aoc1NuskATPsyBeKSYK7J), 2026-11-12; reviewed against both supplied scan halves; see CHECK_DECISIONS.tsv.
\scanpage{134-8/23}
\mspage{577}

Thus we see that when working with the notion of properness of
maps, even for such simple gadgets as finite ordered sets, which \var{\xout{shoul}}{}{typewriter deletion}{ps_134-8_023-v0001}
should be viewed as "proper" (or "compact") objects by themselves,
this notion is far from being a wholly trivial one - for instance it
does not hold t\var{\over[typed]{t}{r}}{r}{typewriter overwriting}{ps_134-8_023-v0002}ue that any map between such "p\var{\over[typed]{m}{r}}{r}{typewriter overwriting}{ps_134-8_023-v0003}oper" objects is again
"proper". This now seems to me just a mathematical "fact of life",
which we may not disregard when working with finite ordered sets say,
in view of expressing in algebraic terms some standard operations \var{\over[typed]{o}{i}}{i}{typewriter overwriting}{ps_134-8_023-v0004}n
the cohomology theory of sheaves on topological spaces, endowed with
suita$\underset{b}{b}$le conical stratifications. The fault surely is not with the
notion of conical stratification itself, which may be felt by some
as being ad hoc, awkward and what not. I know the notion is just
right - but at any rate, even when working with strata which are
perfect topological cells (so that nobody could possibly object to
them), exactly the same facts of life are there - \var{\inserted{\deleted{\illegible{1}}}}{\hand{}}{insertion}{ps_134-8_023-v0005}\hand{(}not every map \var{\xout{of}}{}{typewriter deletion}{ps_134-8_023-v0006}
between such cellular decompositions, mapping cells \uline{onto} cells, will
fit into a"combinatorial"description,when it comes to describing the
standard operations of the cohomology of sheaves, for sheaves compatib-%
le with the stratifications ...

To sum up, it seems to me that definitely, there is a very rich
experimental mat\var{\xout{r}}{}{typewriter deletion}{ps_134-8_023-v0007}erial available already at present, in order to come
to a feeling of what duality is like in the context of (Cat), and \var{\xout{for}}{}{typewriter deletion}{ps_134-8_023-v0008}
for developping some of the basic intuitions needed for working out,
hopefully\var{\xout{\illegible{1}}}{}{typewriter deletion}{ps_134-8_023-v0009}, "the" full-\var{\over[typed]{\illegible{2}}{fl}}{fl}{typewriter overwriting}{ps_134-8_023-v0010}edged duality theory which should hold in
(Cat). Before leaving this digression, I would like to point out still
one \var{\xout{more}}{}{typewriter deletion}{ps_134-8_023-v0011} other property connected with the \var{\xout{notion}}{}{typewriter deletion}{ps_134-8_023-v0012} intuitions around
"properness" - more gen\var{\over[typed]{m}{e}}{e}{typewriter overwriting}{ps_134-8_023-v0013}rally, around maps "of finite type" in a
suitable sense, which sometimes may translate into: factorizable as
a composition $f = g\circ i$ , with g proper and i an immersion. This is about
stability of "constructibility" o\var{\over[typed]{r}{f}}{f}{typewriter overwriting}{ps_134-8_023-v0014} "finiteness" conditions for (comple-%
xes\var{\xout{\illegible{1}}}{}{typewriter deletion}{ps_134-8_023-v0015} of) sheaves of k-modules\hand{,} with respect to the standard operations
$Rf_{*}$ , $Rf_{!}$ , $f^{!}$ (stability by $f^{*}$ being a tautology in any case). Such
stability of course, whenever it holds, is an important feature,
for instance for making "virtual" calculations in suitable "Grothendieck
groups" (where Euler-Poincaré type invariants may be defined). It shoul
% NOTE: the right edge of the scan cuts the end of the line after "shoul"
be recalled however that the six fundamental operations in the duality
formulaire, as well as nearly all of the formulaire itself, make a
sense (and \var{\inserted{formulae}}{\hand{formulae}}{insertion}{ps_134-8_023-v0016} hold\var{\deleted{s}}{}{deletion}{ps_134-8_023-v0017} true) without any finiteness conditions on the complexes
of sheaves we work with, except just \var{\deleted{\xout{restrictions} bounds}}{}{deletion}{ps_134-8_023-v0018} \var{\inserted{boundedness conditions}}{\hand{boundedness conditions}}{insertion}{ps_134-8_023-v0019} on the degrees
of those complexes.

\margin{23}

% Visual collation: both supplied image halves inspected; see CHECK_DECISIONS_134_7_8.tsv.

%% END SOURCE ps_134-8_023
\ifnum\EditionMode=1
\psapparatusstart
\psentry{ps_134-8_023-v0001}{1}{typewriter deletion}{\xout{shoul}}{}{The typewriter strikes out this material; it leaves the corrected running text.}
\psentry{ps_134-8_023-v0002}{2}{typewriter overwriting}{\over[typed]{t}{r}}{r}{The superposed typed reading replaces the underlying reading.}
\psentry{ps_134-8_023-v0003}{3}{typewriter overwriting}{\over[typed]{m}{r}}{r}{The superposed typed reading replaces the underlying reading.}
\psentry{ps_134-8_023-v0004}{4}{typewriter overwriting}{\over[typed]{o}{i}}{i}{The superposed typed reading replaces the underlying reading.}
\psentry{ps_134-8_023-v0005}{5}{insertion}{\inserted{\deleted{\illegible{1}}}}{\hand{}}{The handwritten insertion is incorporated at its indicated position; insertion brackets leave the running text.}
\psentry{ps_134-8_023-v0006}{6}{typewriter deletion}{\xout{of}}{}{The typewriter strikes out this material; it leaves the corrected running text.}
\psentry{ps_134-8_023-v0007}{7}{typewriter deletion}{\xout{r}}{}{The typewriter strikes out this material; it leaves the corrected running text.}
\psentry{ps_134-8_023-v0008}{8}{typewriter deletion}{\xout{for}}{}{The typewriter strikes out this material; it leaves the corrected running text.}
\psentry{ps_134-8_023-v0009}{9}{typewriter deletion}{\xout{\illegible{1}}}{}{The typewriter strikes out this material; it leaves the corrected running text.}
\psentry{ps_134-8_023-v0010}{10}{typewriter overwriting}{\over[typed]{\illegible{2}}{fl}}{fl}{The superposed typed reading replaces the underlying reading.}
\psentry{ps_134-8_023-v0011}{11}{typewriter deletion}{\xout{more}}{}{The typewriter strikes out this material; it leaves the corrected running text.}
\psentry{ps_134-8_023-v0012}{12}{typewriter deletion}{\xout{notion}}{}{The typewriter strikes out this material; it leaves the corrected running text.}
\psentry{ps_134-8_023-v0013}{13}{typewriter overwriting}{\over[typed]{m}{e}}{e}{The superposed typed reading replaces the underlying reading.}
\psentry{ps_134-8_023-v0014}{14}{typewriter overwriting}{\over[typed]{r}{f}}{f}{The superposed typed reading replaces the underlying reading.}
\psentry{ps_134-8_023-v0015}{15}{typewriter deletion}{\xout{\illegible{1}}}{}{The typewriter strikes out this material; it leaves the corrected running text.}
\psentry{ps_134-8_023-v0016}{16}{insertion}{\inserted{formulae}}{\hand{formulae}}{The handwritten insertion is incorporated at its indicated position; insertion brackets leave the running text.}
\psentry{ps_134-8_023-v0017}{17}{deletion}{\deleted{s}}{}{The author cancels this material by hand; it leaves the corrected running text.}
\psentry{ps_134-8_023-v0018}{18}{deletion}{\deleted{\xout{restrictions} bounds}}{}{The author cancels this material by hand; it leaves the corrected running text.}
\psentry{ps_134-8_023-v0019}{19}{insertion}{\inserted{boundedness conditions}}{\hand{boundedness conditions}}{The handwritten insertion is incorporated at its indicated position; insertion brackets leave the running text.}
\psapparatusend
\fi
\pssource{ps_134-8_024}{134-8/24}{https://grothendieck.umontpellier.fr/web/viewer.html?file=../134-8.pdf\#page=25}
\psnumber{ps_134-8_024-v0001}{1}
\psnumber{ps_134-8_024-v0002}{2}
\psnumber{ps_134-8_024-v0003}{3}
\psnumber{ps_134-8_024-v0004}{4}
\psnumber{ps_134-8_024-v0005}{5}
\psnumber{ps_134-8_024-v0006}{6}
\psnumber{ps_134-8_024-v0007}{7}
\psnumber{ps_134-8_024-v0008}{8}
\psnumber{ps_134-8_024-v0009}{9}
\psnumber{ps_134-8_024-v0010}{10}
\psnumber{ps_134-8_024-v0011}{11}
\psnumber{ps_134-8_024-v0012}{12}
\psnumber{ps_134-8_024-v0013}{13}
\psnumber{ps_134-8_024-v0014}{14}
\psnumber{ps_134-8_024-v0015}{15}
\psnumber{ps_134-8_024-v0016}{16}
\psnumber{ps_134-8_024-v0017}{17}
\psnumber{ps_134-8_024-v0018}{18}
%% BEGIN SOURCE ps_134-8_024
% Pursuing Stacks (cotes 134-2 to 134-8): archive cote 134-8, site page 24 (PDF page 25), typescript page 578.
% typescript. Draft compared: ps7.tex:1110-1169.
% AI-transcribed by Claude Opus 5.5 (claude-opus-5-5, adaptive thinking, effort high), Message Batches API, batch ps4 (msgbatch_012Aoc1NuskATPsyBeKSYK7J), 2026-11-12; reviewed against both supplied scan halves; see CHECK_DECISIONS.tsv.
\scanpage{134-8/24}
\mspage{578}

\uline{Remark\hand{s}}\hand{, 1)}One of the (rather few) instances in the duality formulaire
where finiteness conditions are clearly needed,  is the so-called
"biduality theorem",  when taking \uline{RHom}'s with values in a so-called
"dualizing complex", \var{\deleted{\xout{\unsure{R}\illegible{1}D$^{\mathrm b}$(\unsure{A}}}}{}{deletion}{ps_134-8_024-v0001} which here should be an object \var{\xout{in}}{}{typewriter deletion}{ps_134-8_024-v0002}
\[ R^{\bullet} \quad \text{in} \quad D^{b}(A^{\wedge}_{k}) \]
(for given A in (Cat) and given ring of coefficients\textsubscript{x} k,  for instance
k=$\underline{Z}$  ).  The question arises here, \var{\deleted{\xout{a}}}{}{deletion}{ps_134-8_024-v0003} for any given A and k,  a) whe-%
ther there exists a dualizing complex R$^{\bullet}$  (which,as usual,  will be
unique up to \var{\inserted{dimension shift and}}{\hand{dimension shift and}}{insertion}{ps_134-8_024-v0004} \hand{"}twist\hand{"} by an invertible sheaf of k-modules on A),  and
b) can such dualizing complex be obtained as
\[ R^{\bullet} = p^{!}(k_{e}) \quad , \]
where  p: A $\hand{\longrightarrow}$ e is the projection to the final \var{\over[typed]{f\illegible{2}}{obj}}{obj}{typewriter overwriting}{ps_134-8_024-v0005}ect of (Cat) ?  It
should be easy enough,  once the basic duality formalism is written down
for ordered sets,  as contemplated above, \var{\inserted{to show}}{\hand{to show}}{insertion}{ps_134-8_024-v0006} that when A is a finite or-%
dered set (or only locally finite of finite combinatorial dimension),
then $p^{!}(e_{k})$  is indeed a dualizing complex.  To refer to something
more "en vogue" at present than those poor dualizing complexes,  it is
clear,  from what I heard from Illusie and Mebkhout about the "complexe
d'intersection" for stratified spaces,  that this complex can be descri-%
bed also in the context of locally finite ordered sets (if I got the
stepwise construction of this complex right).  As the construction
here corresponds to stratifications where the strata are by no means
even dimensional,  I am not too sure though if the complex obtained this
way is really relevant  -  it is'nt a topological invariant of the
topological space at any rate,  independently of it's stratification -
a bad point indeed. Too bad !

2)  Here is a rather evident example showing that the asphericity
condition on (13) is not automatic.  Take  I  with a \var{\xout{largest}}{}{typewriter deletion}{ps_134-8_024-v0007} \var{\inserted[typed]{smallest}}{smallest}{typed insertion}{ps_134-8_024-v0008} element x
(hence  $C_{x} = \hand{|}I\hand{|}$ ),  and J =$\hand{\Delta}_{1}$ \var{\inserted{$= (0 \to 1)$.}}{\hand{$= (0 \to 1)$.}}{insertion}{ps_134-8_024-v0009} \var{\deleted{, consisting of 0 and 1 \inserted[typed]{= y} \xout{\illegible{1}}}}{}{deletion}{ps_134-8_024-v0010}.  To give
a map  I$\hand{\rightarrow}$ J  amounts to give  $I_{o} = f^{-1}(0) \hand{\subset} I$,  which is any \var{\inserted{\unsure{``}}}{\hand{\unsure{``}}}{insertion}{ps_134-8_024-v0011}open subset
of I (i.e. such that $a \hand{\in} I_{o}$ ,  $b \hand{\leqslant} a$  implies $b \hand{\in} I_{o}$).  \var{\inserted{If we take}}{\hand{If we take}}{insertion}{ps_134-8_024-v0012} \var{\deleted{The relation}}{}{deletion}{ps_134-8_024-v0013}
\var{\deleted{f(x)=y just means that \xout{I$_o$\illegible{1}$\hand{\neq}$\illegible{1}I\illegible{1}} I$_o$}}{}{deletion}{ps_134-8_024-v0014}.  $I_{o} = \hand{\{}x\hand{\}}$,  then properness of
f is equivalent with  $I_{1} = I \hand{\smallsetminus} I_{o} = I\hand{\smallsetminus}\hand{\{}x\hand{\}}$ being aspheric,  while the
weaker condition contemplated first (the sets (13) non empty) means only
that $I_{1} \hand{\neq} \hand{\emptyset}$ .  N$\overset{\illegible{1}}{\underset{\unsure{o}}{o}}$w $I_{1}$ \var{\xout{can}}{}{typewriter deletion}{ps_134-8_024-v0015} \var{\inserted[typed]{may}}{may}{typed insertion}{ps_134-8_024-v0016} of course be taken to be any (finite say)
ordered set - the construction made amounting to taking the cone over
the map $\hand{|}I_{1}\hand{|} \hand{\longrightarrow} e$ .   In this example,  all fi\var{\over{\unsure{b}}{b}}{\hand{b}}{writing over}{ps_134-8_024-v0017}ers of $\hand{|}f\hand{|} :\hand{|}I\hand{|}\hand{\rightarrow}\hand{|}J\hand{|}$
are (canonically) homeomorphic to  $I_{1}$  , except the fiber at \var{\inserted{the "barycenter"}}{\hand{the "barycenter"}}{insertion}{ps_134-8_024-v0018} 1, which

\margin{24}

% Visual collation: both supplied image halves inspected; see CHECK_DECISIONS_134_7_8.tsv.

%% END SOURCE ps_134-8_024
\ifnum\EditionMode=1
\psapparatusstart
\psentry{ps_134-8_024-v0001}{1}{deletion}{\deleted{\xout{\unsure{R}\illegible{1}D$^{\mathrm b}$(\unsure{A}}}}{}{The author cancels this material by hand; it leaves the corrected running text.}
\psentry{ps_134-8_024-v0002}{2}{typewriter deletion}{\xout{in}}{}{The typewriter strikes out this material; it leaves the corrected running text.}
\psentry{ps_134-8_024-v0003}{3}{deletion}{\deleted{\xout{a}}}{}{The author cancels this material by hand; it leaves the corrected running text.}
\psentry{ps_134-8_024-v0004}{4}{insertion}{\inserted{dimension shift and}}{\hand{dimension shift and}}{The handwritten insertion is incorporated at its indicated position; insertion brackets leave the running text.}
\psentry{ps_134-8_024-v0005}{5}{typewriter overwriting}{\over[typed]{f\illegible{2}}{obj}}{obj}{The superposed typed reading replaces the underlying reading.}
\psentry{ps_134-8_024-v0006}{6}{insertion}{\inserted{to show}}{\hand{to show}}{The handwritten insertion is incorporated at its indicated position; insertion brackets leave the running text.}
\psentry{ps_134-8_024-v0007}{7}{typewriter deletion}{\xout{largest}}{}{The typewriter strikes out this material; it leaves the corrected running text.}
\psentry{ps_134-8_024-v0008}{8}{typed insertion}{\inserted[typed]{smallest}}{smallest}{The typed insertion is incorporated at its indicated position; insertion brackets leave the running text.}
\psentry{ps_134-8_024-v0009}{9}{insertion}{\inserted{$= (0 \to 1)$.}}{\hand{$= (0 \to 1)$.}}{The handwritten insertion is incorporated at its indicated position; insertion brackets leave the running text.}
\psentry{ps_134-8_024-v0010}{10}{deletion}{\deleted{, consisting of 0 and 1 \inserted[typed]{= y} \xout{\illegible{1}}}}{}{The author cancels this material by hand; it leaves the corrected running text.}
\psentry{ps_134-8_024-v0011}{11}{insertion}{\inserted{\unsure{``}}}{\hand{\unsure{``}}}{The handwritten insertion is incorporated at its indicated position; insertion brackets leave the running text.}
\psentry{ps_134-8_024-v0012}{12}{insertion}{\inserted{If we take}}{\hand{If we take}}{The handwritten insertion is incorporated at its indicated position; insertion brackets leave the running text.}
\psentry{ps_134-8_024-v0013}{13}{deletion}{\deleted{The relation}}{}{The author cancels this material by hand; it leaves the corrected running text.}
\psentry{ps_134-8_024-v0014}{14}{deletion}{\deleted{f(x)=y just means that \xout{I$_o$\illegible{1}$\hand{\neq}$\illegible{1}I\illegible{1}} I$_o$}}{}{The author cancels this material by hand; it leaves the corrected running text.}
\psentry{ps_134-8_024-v0015}{15}{typewriter deletion}{\xout{can}}{}{The typewriter strikes out this material; it leaves the corrected running text.}
\psentry{ps_134-8_024-v0016}{16}{typed insertion}{\inserted[typed]{may}}{may}{The typed insertion is incorporated at its indicated position; insertion brackets leave the running text.}
\psentry{ps_134-8_024-v0017}{17}{writing over}{\over{\unsure{b}}{b}}{\hand{b}}{The superposed handwritten reading replaces the underlying reading.}
\psentry{ps_134-8_024-v0018}{18}{insertion}{\inserted{the "barycenter"}}{\hand{the "barycenter"}}{The handwritten insertion is incorporated at its indicated position; insertion brackets leave the running text.}
\psapparatusend
\fi
\pssource{ps_134-8_025}{134-8/25}{https://grothendieck.umontpellier.fr/web/viewer.html?file=../134-8.pdf\#page=26}
\psnumber{ps_134-8_025-v0001}{1}
\psnumber{ps_134-8_025-v0002}{2}
\psnumber{ps_134-8_025-v0003}{3}
\psnumber{ps_134-8_025-v0004}{4}
\psnumber{ps_134-8_025-v0005}{5}
\psnumber{ps_134-8_025-v0006}{6}
\psnumber{ps_134-8_025-v0007}{7}
\psnumber{ps_134-8_025-v0008}{8}
\psnumber{ps_134-8_025-v0009}{9}
\psnumber{ps_134-8_025-v0010}{10}
\psnumber{ps_134-8_025-v0011}{11}
\psnumber{ps_134-8_025-v0012}{12}
\psnumber{ps_134-8_025-v0013}{13}
\psnumber{ps_134-8_025-v0014}{14}
\psnumber{ps_134-8_025-v0015}{15}
\psnumber{ps_134-8_025-v0016}{16}
\psnumber{ps_134-8_025-v0017}{17}
\psnumber{ps_134-8_025-v0018}{18}
\psnumber{ps_134-8_025-v0019}{19}
\psnumber{ps_134-8_025-v0020}{20}
%% BEGIN SOURCE ps_134-8_025
% Pursuing Stacks (cotes 134-2 to 134-8): archive cote 134-8, site page 25 (PDF page 26), typescript page 579.
% typescript. Draft compared: ps7.tex:1157-1213.
% AI-transcribed by Claude Opus 5.5 (claude-opus-5-5, adaptive thinking, effort high), Message Batches API, batch ps4 (msgbatch_012Aoc1NuskATPsyBeKSYK7J), 2026-11-12; reviewed against both supplied scan halves; see CHECK_DECISIONS.tsv.
\scanpage{134-8/25}
\mspage{579}

is reduced to a point  -  visibly not a very "cellular" behaviour when
$\hand{|}I_1\hand{|}$ is'nt aspheric !  At any rate,  as the sheaves \var{\xout{$R^if_*$ may be compu}}{}{typewriter deletion}{ps_134-8_025-v0001}
$R^i\hand{|}f\hand{|}_*(F)$  (for F a constant sheaf above say) may be computed fiberwise,
we see that if $I_1$  is'nt aspheric,  these sheaves (which are constant
on $\hand{|}J\hand{|}\hand{\setminus\{}1\hand{\}}$ )  are \uline{not} going to be constant on  $\hand{|}J\hand{|}\hand{\setminus\{}0\hand{\}}$, as \var{\xout{we expect}}{}{typewriter deletion}{ps_134-8_025-v0002}
\var{\xout{for}}{}{typewriter deletion}{ps_134-8_025-v0003} they should if we want an algebraic paradigm of operations like
$R\hand{|}f\hand{|}_*$  in terms of sheaves on finite ordered sets.

In this example we could take $\hand{|}I\hand{|}$ to be a perfect n-cell (n$\hand{\geq}$1),
hence  $I_1$  is an (n-1)-sphere,  whereas $\hand{|}J\hand{|}$ is'nt really a (combinatori-%
al) 1-cell,  as it's bou\var{\over[typed]{d}{n}}{n}{typewriter overwriting}{ps_134-8_025-v0004}dary has just \uline{one} point 0 ,  instead of two.
T\var{\over[typed]{\unsure{H|h}}{h}}{h}{typewriter overwriting}{ps_134-8_025-v0005}e 1-cell structure corresponds to the ordered set (\uline{opposite} to the
ordered set formed by the two endpoints and the dimension 1 stratum)
\diagram{\[ J = \begin{tikzcd}[row sep=small, column sep=small] \& \var{\over{y}{y}}{\hand{y}}{writing over}{ps_134-8_025-v0006} \ar[color=blue,dl] \ar[color=blue,dr] \& \\ \var{\over[typed]{0}{0}}{0}{typewriter overwriting}{ps_134-8_025-v0007} \& \& 1 \end{tikzcd} \quad . \]}
% Hand-drawn arrows are marked blue; typed objects remain plain.
For any ordered set I,  to give a map  f: I $\hand{\longrightarrow}$ J  amount to give
the two subsets
\[ I_0 = f^{-1}(0) \ , \ \ \var{\deleted{\xout{\unsure{et}}}}{}{deletion}{ps_134-8_025-v0008}\ I_1 = f^{-1}(1) \qquad , \]
subject\var{\xout{s}}{}{typewriter deletion}{ps_134-8_025-v0009} to the only condition of being open,  and disjoint.  In terms
of strata of $\hand{|}I\hand{|}$ ,  this means that we give \var{\over{to}{two}}{\hand{two}}{writing over}{ps_134-8_025-v0010} sets $I_0,I_1$ of strata,
containing with any stratum any smaller one,  and \var{\over{with}{having}}{\hand{having}}{writing over}{ps_134-8_025-v0011} no stratum \var{\inserted{in}}{\hand{in}}{insertion}{ps_134-8_025-v0012}
common \var{\deleted{to $I_0$ and $I_1$}}{}{deletion}{ps_134-8_025-v0013}.  The condition that the sets (13) should be non-%
empty says that any point in I which is neither in $I_0$ no\var{\over[typed]{t}{r}}{r}{typewriter overwriting}{ps_134-8_025-v0014} $I_1$ , admits
majorants in both  -  or geometrically,  any stratum which is'nt in
$I_0$ nor in $I_1$ meets both $\hand{|}I_0\hand{|}$ and $\hand{|}I_1\hand{|}$ i.e. contains strata which are
in $I_0$ and strata which are in $I_1$ .   Even when $\hand{|}I\hand{|}$ is a combinatorial
2-cell i.e. a polygonal disc,  this condition does not imply asphérici-%
ty of the sets (13) (not even 0-connectedness).  To see this,  we take
a polygon of even order 2m $\hand{\geq}$\var{\over{\illegible{1}}{4}}{\hand{4}}{writing over}{ps_134-8_025-v0015} ,  and we take  $I_0$ and $I_1$ to consist
of the vertices of the polygon of even resp. odd index.  We take the
set (13) where x is the dimension 2 stratum  i.e. the \var{\xout{mini}}{}{typewriter deletion}{ps_134-8_025-v0016} smallest
element of I ,  mapped to y (the smallest in J),  and y' either 0 or 1.
The reader will easily figure out on a drawing the structure of the
map $\hand{|}f\hand{|}$ ,  as I just did myself:  the fi\var{\over[typed]{v}{b}}{b}{typewriter overwriting}{ps_134-8_025-v0017}ers at the endpoints of the
segment $\hand{|}J\hand{|}$ are,  as given,  discrete with cardinal m,  the fiber at
a point different f\var{\over[typed]{o}{r}}{r}{typewriter overwriting}{ps_134-8_025-v0018}\var{\over[typed]{r}{o}}{o}{typewriter overwriting}{ps_134-8_025-v0019}m the endpoints and \var{\inserted{from}}{\hand{from}}{insertion}{ps_134-8_025-v0020} their barycenter are disjoint
sums of m segments (hence homotopic to the former fibers),  whereas
the fiber at the barycenter is the union of m segments meeting in their

\margin{25}

% Visual collation: both supplied image halves inspected; see CHECK_DECISIONS_134_7_8.tsv.

%% END SOURCE ps_134-8_025
\ifnum\EditionMode=1
\psapparatusstart
\psentry{ps_134-8_025-v0001}{1}{typewriter deletion}{\xout{$R^if_*$ may be compu}}{}{The typewriter strikes out this material; it leaves the corrected running text.}
\psentry{ps_134-8_025-v0002}{2}{typewriter deletion}{\xout{we expect}}{}{The typewriter strikes out this material; it leaves the corrected running text.}
\psentry{ps_134-8_025-v0003}{3}{typewriter deletion}{\xout{for}}{}{The typewriter strikes out this material; it leaves the corrected running text.}
\psentry{ps_134-8_025-v0004}{4}{typewriter overwriting}{\over[typed]{d}{n}}{n}{The superposed typed reading replaces the underlying reading.}
\psentry{ps_134-8_025-v0005}{5}{typewriter overwriting}{\over[typed]{\unsure{H|h}}{h}}{h}{The superposed typed reading replaces the underlying reading.}
\psentry{ps_134-8_025-v0006}{6}{writing over}{$\over{y}{y}$}{$\hand{y}$}{The superposed handwritten reading replaces the underlying reading.}
\psentry{ps_134-8_025-v0007}{7}{typewriter overwriting}{$\over[typed]{0}{0}$}{$0$}{The superposed typed reading replaces the underlying reading.}
\psentry{ps_134-8_025-v0008}{8}{deletion}{$\deleted{\xout{\unsure{et}}}$}{}{The author cancels this material by hand; it leaves the corrected running text.}
\psentry{ps_134-8_025-v0009}{9}{typewriter deletion}{\xout{s}}{}{The typewriter strikes out this material; it leaves the corrected running text.}
\psentry{ps_134-8_025-v0010}{10}{writing over}{\over{to}{two}}{\hand{two}}{The superposed handwritten reading replaces the underlying reading.}
\psentry{ps_134-8_025-v0011}{11}{writing over}{\over{with}{having}}{\hand{having}}{The superposed handwritten reading replaces the underlying reading.}
\psentry{ps_134-8_025-v0012}{12}{insertion}{\inserted{in}}{\hand{in}}{The handwritten insertion is incorporated at its indicated position; insertion brackets leave the running text.}
\psentry{ps_134-8_025-v0013}{13}{deletion}{\deleted{to $I_0$ and $I_1$}}{}{The author cancels this material by hand; it leaves the corrected running text.}
\psentry{ps_134-8_025-v0014}{14}{typewriter overwriting}{\over[typed]{t}{r}}{r}{The superposed typed reading replaces the underlying reading.}
\psentry{ps_134-8_025-v0015}{15}{writing over}{\over{\illegible{1}}{4}}{\hand{4}}{The superposed handwritten reading replaces the underlying reading.}
\psentry{ps_134-8_025-v0016}{16}{typewriter deletion}{\xout{mini}}{}{The typewriter strikes out this material; it leaves the corrected running text.}
\psentry{ps_134-8_025-v0017}{17}{typewriter overwriting}{\over[typed]{v}{b}}{b}{The superposed typed reading replaces the underlying reading.}
\psentry{ps_134-8_025-v0018}{18}{typewriter overwriting}{\over[typed]{o}{r}}{r}{The superposed typed reading replaces the underlying reading.}
\psentry{ps_134-8_025-v0019}{19}{typewriter overwriting}{\over[typed]{r}{o}}{o}{The superposed typed reading replaces the underlying reading.}
\psentry{ps_134-8_025-v0020}{20}{insertion}{\inserted{from}}{\hand{from}}{The handwritten insertion is incorporated at its indicated position; insertion brackets leave the running text.}
\psapparatusend
\fi
\pssource{ps_134-8_026}{134-8/26}{https://grothendieck.umontpellier.fr/web/viewer.html?file=../134-8.pdf\#page=27}
\psnumber{ps_134-8_026-v0001}{1}
\psnumber{ps_134-8_026-v0002}{2}
\psnumber{ps_134-8_026-v0003}{3}
\psnumber{ps_134-8_026-v0004}{4}
\psnumber{ps_134-8_026-v0005}{5}
\psnumber{ps_134-8_026-v0006}{6}
\psnumber{ps_134-8_026-v0007}{7}
\psnumber{ps_134-8_026-v0008}{8}
\psnumber{ps_134-8_026-v0009}{9}
\psnumber{ps_134-8_026-v0010}{10}
\psnumber{ps_134-8_026-v0011}{11}
\psnumber{ps_134-8_026-v0012}{12}
\psnumber{ps_134-8_026-v0013}{13}
\psnumber{ps_134-8_026-v0014}{14}
\psnumber{ps_134-8_026-v0015}{15}
\psnumber{ps_134-8_026-v0016}{16}
\psnumber{ps_134-8_026-v0017}{17}
\psnumber{ps_134-8_026-v0018}{18}
\psnumber{ps_134-8_026-v0019}{19}
\psnumber{ps_134-8_026-v0020}{20}
\psnumber{ps_134-8_026-v0021}{21}
\psnumber{ps_134-8_026-v0022}{22}
\psnumber{ps_134-8_026-v0023}{23}
\psnumber{ps_134-8_026-v0024}{24}
\psnumber{ps_134-8_026-v0025}{25}
\psnumber{ps_134-8_026-v0026}{26}
%% BEGIN SOURCE ps_134-8_026
% Pursuing Stacks (cotes 134-2 to 134-8): archive cote 134-8, site page 26 (PDF page 27), typescript page 580.
% typescript. Draft compared: ps7.tex:1201-1263.
% AI-transcribed by Claude Opus 5.5 (claude-opus-5-5, adaptive thinking, effort high), Message Batches API, batch ps4 (msgbatch_012Aoc1NuskATPsyBeKSYK7J), 2026-11-12; reviewed against both supplied scan halves; see CHECK_DECISIONS.tsv.
\scanpage{134-8/26}
\mspage{580}
common middle, hence is contractible. Thus the $R^{0}\hand{|}f\hand{|}_{*}$ of a constant
sheaf on the disc $\hand{|}I\hand{|}$ is by no means constant on the open\hand{,} dimension one
stratum of $\hand{|}J\hand{|}$ .

Th\var{\over[typed]{i}{e}}{e}{typewriter overwriting}{ps_134-8_026-v0001}se examples bring\var{\xout{s}}{}{typewriter deletion}{ps_134-8_026-v0002} to my mind that for any map $f: I \hand{\to} J$ between
finite ordered sets, \var{\xout{satisfying the}}{}{typewriter deletion}{ps_134-8_026-v0003} possibly submitted to the mild
restriction that the sets (13) should \var{\over{\unsure{b}}{b}}{\hand{b}}{writing over}{ps_134-8_026-v0004}e non-empty, the homotopy types
of those \var{\xout{set}}{}{typewriter deletion}{ps_134-8_026-v0005} ordered sets \var{\inserted[typed]{(13)}}{(13)}{typed insertion}{ps_134-8_026-v0006} (for the order relation induced by I) should
be exactly the homotopy types of the fibers of the \var{\over{\unsure{m}}{m}}{\hand{m}}{writing over}{ps_134-8_026-v0007}aps \var{\xout{f\illegible{1}}}{}{typewriter deletion}{ps_134-8_026-v0008} \var{\inserted[typed]{(12) between strata.}}{(12) between strata.}{typed insertion}{ps_134-8_026-v0009}
Thus asphericity of these ordered sets should express no more, no less
than the con\var{\over{\unsure{t}}{t}}{\hand{t}}{writing over}{ps_134-8_026-v0010}\var{\over[typed]{t}{r}}{r}{typewriter overwriting}{ps_134-8_026-v0011}actibility of those fibers.\var{\deleted{(}}{}{deletion}{ps_134-8_026-v0012} This latter condition is
exactly what is needed in order to insure stability by $R\hand{|}f\hand{|}_{*}$ of the
notion of complexes of sheaves compatible with the stratifications.

28.10.

138. Yesterday and the day before\hand{,} I got involved \var{\xout{again}}{}{typewriter deletion}{ps_134-8_026-v0013} in that
unforeseen digression\var{\deleted{\unsure{,}}}{}{deletion}{ps_134-8_026-v0014} around the foreboding of a ``six operations
duality formalism'' in (Cat), and suitable notions for smoothness and
(more important still) of properness for a map in (Cat). This digres-%
sion wholly convinced me that the usual duality formalism should hold
in (Cat) too, Working this out in full detail should be a most pleasant
task indeed, and presumably the best, or even the only way for gaining
complete mastery of the cohomology formalism within (Cat) - or, what
more or less amounts to the same, for topoi which admit sufficiently
many projective objects. In the previous two sections, I referred to
such duality formalism as \var{\inserted{one}}{\hand{one}}{insertion}{ps_134-8_026-v0015} concerned with sheaves of k-modules, for any
fixed ring k - but from the example of étale duality \var{\xout{say, it}}{}{typewriter deletion}{ps_134-8_026-v0016} for sche-%
mes say, it is likely that instead of fixing a ring k, we may as
well take objects A in (Cat) endowed with an arbitrary sheaf of rings
$\underline{O}_{A}$ (which we'll only have to suppose commutative, \var{\over{\unsure{wh}}{wh}}{\hand{wh}}{writing over}{ps_134-8_026-v0017}en concerned with
the two\hand{,} internal operations $\overset{\var{\over{L}{L}}{\hand{L}}{writing over}{ps_134-8_026-v0018}}{\otimes}$ and $R\underline{Hom}$ ), and taking maps of such
ringed objects as the basic maps. In the present context, the \var{\inserted[typed]{usual}}{usual}{typed insertion}{ps_134-8_026-v0019} ``six
operations'' in duality theory will be enriched however by still another
\margin{! $\times$} one, namely $Lf_{!}$ , \var{\xout{valid}}{}{typewriter deletion}{ps_134-8_026-v0020} \var{\inserted[typed]{defined}}{defined}{typed insertion}{ps_134-8_026-v0021} for any map f between such ringed objects\var{\deleted{\unsure{,}}}{}{deletion}{ps_134-8_026-v0022}
\hand{(}not to be confused with the $Rf_{!}$ operation\hand{,} \var{\deleted{(}}{}{deletion}{ps_134-8_026-v0023}defined only under sui-%
table finiteness asum\hand{\illegible{1}}ptions, such as ``properness'', for the underly-%
ing map in (Cat)), whose relationship to the other operations should
be understood and \var{\deleted{included}}{}{deletion}{ps_134-8_026-v0024} \var{\inserted{added}}{\hand{added}}{insertion}{ps_134-8_026-v0025} \var{\over{in}{to}}{\hand{to}}{writing over}{ps_134-8_026-v0026} the standard duality formulaire. One

\margin{26}

% Visual collation: both supplied image halves inspected; see CHECK_DECISIONS_134_7_8.tsv.

%% END SOURCE ps_134-8_026
\ifnum\EditionMode=1
\psapparatusstart
\psentry{ps_134-8_026-v0001}{1}{typewriter overwriting}{\over[typed]{i}{e}}{e}{The superposed typed reading replaces the underlying reading.}
\psentry{ps_134-8_026-v0002}{2}{typewriter deletion}{\xout{s}}{}{The typewriter strikes out this material; it leaves the corrected running text.}
\psentry{ps_134-8_026-v0003}{3}{typewriter deletion}{\xout{satisfying the}}{}{The typewriter strikes out this material; it leaves the corrected running text.}
\psentry{ps_134-8_026-v0004}{4}{writing over}{\over{\unsure{b}}{b}}{\hand{b}}{The superposed handwritten reading replaces the underlying reading.}
\psentry{ps_134-8_026-v0005}{5}{typewriter deletion}{\xout{set}}{}{The typewriter strikes out this material; it leaves the corrected running text.}
\psentry{ps_134-8_026-v0006}{6}{typed insertion}{\inserted[typed]{(13)}}{(13)}{The typed insertion is incorporated at its indicated position; insertion brackets leave the running text.}
\psentry{ps_134-8_026-v0007}{7}{writing over}{\over{\unsure{m}}{m}}{\hand{m}}{The superposed handwritten reading replaces the underlying reading.}
\psentry{ps_134-8_026-v0008}{8}{typewriter deletion}{\xout{f\illegible{1}}}{}{The typewriter strikes out this material; it leaves the corrected running text.}
\psentry{ps_134-8_026-v0009}{9}{typed insertion}{\inserted[typed]{(12) between strata.}}{(12) between strata.}{The typed insertion is incorporated at its indicated position; insertion brackets leave the running text.}
\psentry{ps_134-8_026-v0010}{10}{writing over}{\over{\unsure{t}}{t}}{\hand{t}}{The superposed handwritten reading replaces the underlying reading.}
\psentry{ps_134-8_026-v0011}{11}{typewriter overwriting}{\over[typed]{t}{r}}{r}{The superposed typed reading replaces the underlying reading.}
\psentry{ps_134-8_026-v0012}{12}{deletion}{\deleted{(}}{}{The author cancels this material by hand; it leaves the corrected running text.}
\psentry{ps_134-8_026-v0013}{13}{typewriter deletion}{\xout{again}}{}{The typewriter strikes out this material; it leaves the corrected running text.}
\psentry{ps_134-8_026-v0014}{14}{deletion}{\deleted{\unsure{,}}}{}{The author cancels this material by hand; it leaves the corrected running text.}
\psentry{ps_134-8_026-v0015}{15}{insertion}{\inserted{one}}{\hand{one}}{The handwritten insertion is incorporated at its indicated position; insertion brackets leave the running text.}
\psentry{ps_134-8_026-v0016}{16}{typewriter deletion}{\xout{say, it}}{}{The typewriter strikes out this material; it leaves the corrected running text.}
\psentry{ps_134-8_026-v0017}{17}{writing over}{\over{\unsure{wh}}{wh}}{\hand{wh}}{The superposed handwritten reading replaces the underlying reading.}
\psentry{ps_134-8_026-v0018}{18}{writing over}{$\over{L}{L}$}{$\hand{L}$}{The superposed handwritten reading replaces the underlying reading.}
\psentry{ps_134-8_026-v0019}{19}{typed insertion}{\inserted[typed]{usual}}{usual}{The typed insertion is incorporated at its indicated position; insertion brackets leave the running text.}
\psentry{ps_134-8_026-v0020}{20}{typewriter deletion}{\xout{valid}}{}{The typewriter strikes out this material; it leaves the corrected running text.}
\psentry{ps_134-8_026-v0021}{21}{typed insertion}{\inserted[typed]{defined}}{defined}{The typed insertion is incorporated at its indicated position; insertion brackets leave the running text.}
\psentry{ps_134-8_026-v0022}{22}{deletion}{\deleted{\unsure{,}}}{}{The author cancels this material by hand; it leaves the corrected running text.}
\psentry{ps_134-8_026-v0023}{23}{deletion}{\deleted{(}}{}{The author cancels this material by hand; it leaves the corrected running text.}
\psentry{ps_134-8_026-v0024}{24}{deletion}{\deleted{included}}{}{The author cancels this material by hand; it leaves the corrected running text.}
\psentry{ps_134-8_026-v0025}{25}{insertion}{\inserted{added}}{\hand{added}}{The handwritten insertion is incorporated at its indicated position; insertion brackets leave the running text.}
\psentry{ps_134-8_026-v0026}{26}{writing over}{\over{in}{to}}{\hand{to}}{The superposed handwritten reading replaces the underlying reading.}
\psapparatusend
\fi
\pssource{ps_134-8_027}{134-8/27}{https://grothendieck.umontpellier.fr/web/viewer.html?file=../134-8.pdf\#page=28}
\psnumber{ps_134-8_027-v0001}{1}
\psnumber{ps_134-8_027-v0002}{2}
\psnumber{ps_134-8_027-v0003}{3}
\psnumber{ps_134-8_027-v0004}{4}
\psnumber{ps_134-8_027-v0005}{5}
\psnumber{ps_134-8_027-v0006}{6}
\psnumber{ps_134-8_027-v0007}{7}
\psnumber{ps_134-8_027-v0008}{8}
\psnumber{ps_134-8_027-v0009}{9}
\psnumber{ps_134-8_027-v0010}{10}
\psnumber{ps_134-8_027-v0011}{11}
\psnumber{ps_134-8_027-v0012}{12}
\psnumber{ps_134-8_027-v0013}{13}
%% BEGIN SOURCE ps_134-8_027
% Pursuing Stacks (cotes 134-2 to 134-8): archive cote 134-8, site page 27 (PDF page 28), typescript page 581.
% typescript. Draft compared: ps7.tex:1249-1313.
% AI-transcribed by Claude Opus 5.5 (claude-opus-5-5, adaptive thinking, effort high), Message Batches API, batch ps4 (msgbatch_012Aoc1NuskATPsyBeKSYK7J), 2026-11-12; reviewed against both supplied scan halves; see CHECK_DECISIONS.tsv.
\scanpage{134-8/27}
\mspage{581}

such formula, namely the precise relation between $Lf_!$ and $Rf_!$ for
a smooth $f$, was given yesterday \var{\inserted{\deleted{yesterday}}}{\hand{}}{insertion}{ps_134-8_027-v0001} (p.572 (6)).

The day before\hand{,} I was out for trying to get a better understanding
of the $Rf_!$ operation\var{\deleted{\unsure{,}}}{}{deletion}{ps_134-8_027-v0002} \hand{(}including the case of non-constant sheaves of
rings for the Modules we work with\hand{)}. This operation still remains unfa-%
miliar to me, very unlike my old friend $Rf_*$ - there is a number of
things about it which are not quite clear \var{\xout{about it}}{}{typewriter deletion}{ps_134-8_027-v0003} yet in my mind,
even when just taking the functor $f_!$ between Modules, before taking a
left derived functor. For instance, for general rings of operators
$\underline{O}_A$ and $\underline{O}_B$ , when $f$ is a ringed map
\[ \hand{(1)} \qquad f\colon (A,\underline{O}_A)\hand{\longrightarrow} (B,\underline{O}_B) \quad , \]
it does'nt seem that formation of $f_!$ commutes to restriction of
rings of operators (to the constant rings $\underline{Z}_A$ and $\underline{Z}_B$ say), namely that
for a given $\underline{O}_A$-module $F$, $f_!(A)$ may be interpreted as just $f^{ab}_{o!}(F)$ with
suitable operat\var{\over[typed]{i}{i}}{i}{typewriter overwriting}{ps_134-8_027-v0004}\var{\over[typed]{o}{o}}{o}{typewriter overwriting}{ps_134-8_027-v0005}\var{\over[typed]{h}{n}}{n}{typewriter overwriting}{ps_134-8_027-v0006}s of $\underline{O}_B$ on the latter, where
\[ \hand{(2)} \qquad f_o : A\hand{\longrightarrow} B \]
is the map in (Cat) underlying $f$. When reducing to a \var{\inserted{suitable}}{\hand{suitable}}{insertion}{ps_134-8_027-v0007} ``universal'' case,
this may be expressed by saying that for given map $f_o$ and given abelian
sheaf $F$ on $A$, if we define
\[ \var{\margin{\deleted{\illegible{1}}}}{}{marginal note revised}{ps_134-8_027-v0008} \qquad \var{\over{\underline{O}}{\underline{O}}}{\hand{\underline{O}}}{writing over}{ps_134-8_027-v0009}_A= \underline{End}_{\underline{Z}}(F) \quad , \quad \underline{O}_B = f_*(\underline{O}_A) \ , \]
there does'nt seem to be a natural operation of $\underline{O}_B$ upon $f^{\hand{ab}}_{\hand{o!}}(F)$; all we
can say, it seems, is that the ring of global sections of $\underline{O}_B$ operates
o\var{\over{n}{n}}{\hand{n}}{writing over}{ps_134-8_027-v0010} $f^{ab}_{\hand{o!}}(F)$ . More generally, reverting to the general case of a map $f$
of ringed objects in (Cat), the ring of global sections of $\underline{O}_B$ operates
on $f^{ab}_{o!}(F)$. When $\underline{O}_B$ is a constant sheaf of rings $k_B$ , this implies
of course that $k$ operates on this abelian sheaf on $B$, from which will
follow by an obvious argument that with this structure of a sheaf of
$k$-modules, $f^{ab}_{o!}(F)$ may indeed be identified with $f_!(F)$.

Going over to $Lf_!$ which we would like to express via $Lf^{ab}_{o!}$ , the
situation is worse, as we still have to check (granting \var{\xout{B to be constan}}{}{typewriter deletion}{ps_134-8_027-v0011}
$\underline{O}_B$ to be constant)that for $F$ a projective $\underline{O}_A$-Module, we got
\[ \hand{(3)} \qquad L_i f^{ab}_{o!}(F)=0 \quad \text{for } i \hand{>} \var{\over{\unsure{0}}{0}}{\hand{0}}{writing over}{ps_134-8_027-v0012} \ . \]
This is'nt always true, even when $B$ is the final object in (Cat),
and $A$ has a final object, hence $\underline{O}_A$ is a projective Module over
\var{\deleted{itself}}{}{deletion}{ps_134-8_027-v0013}

\margin{27}

% Visual collation: both supplied image halves inspected; see CHECK_DECISIONS_134_7_8.tsv.

%% END SOURCE ps_134-8_027
\ifnum\EditionMode=1
\psapparatusstart
\psentry{ps_134-8_027-v0001}{1}{insertion}{\inserted{\deleted{yesterday}}}{\hand{}}{The handwritten insertion is incorporated at its indicated position; insertion brackets leave the running text.}
\psentry{ps_134-8_027-v0002}{2}{deletion}{\deleted{\unsure{,}}}{}{The author cancels this material by hand; it leaves the corrected running text.}
\psentry{ps_134-8_027-v0003}{3}{typewriter deletion}{\xout{about it}}{}{The typewriter strikes out this material; it leaves the corrected running text.}
\psentry{ps_134-8_027-v0004}{4}{typewriter overwriting}{\over[typed]{i}{i}}{i}{The superposed typed reading replaces the underlying reading.}
\psentry{ps_134-8_027-v0005}{5}{typewriter overwriting}{\over[typed]{o}{o}}{o}{The superposed typed reading replaces the underlying reading.}
\psentry{ps_134-8_027-v0006}{6}{typewriter overwriting}{\over[typed]{h}{n}}{n}{The superposed typed reading replaces the underlying reading.}
\psentry{ps_134-8_027-v0007}{7}{insertion}{\inserted{suitable}}{\hand{suitable}}{The handwritten insertion is incorporated at its indicated position; insertion brackets leave the running text.}
\psentry{ps_134-8_027-v0008}{8}{marginal note revised}{$\margin{\deleted{\illegible{1}}}$}{}{The author's internal cancellations and replacements in this marginal note are applied at its source-line call.}
\psentry{ps_134-8_027-v0009}{9}{writing over}{$\over{\underline{O}}{\underline{O}}$}{$\hand{\underline{O}}$}{The superposed handwritten reading replaces the underlying reading.}
\psentry{ps_134-8_027-v0010}{10}{writing over}{\over{n}{n}}{\hand{n}}{The superposed handwritten reading replaces the underlying reading.}
\psentry{ps_134-8_027-v0011}{11}{typewriter deletion}{\xout{B to be constan}}{}{The typewriter strikes out this material; it leaves the corrected running text.}
\psentry{ps_134-8_027-v0012}{12}{writing over}{$\over{\unsure{0}}{0}$}{$\hand{0}$}{The superposed handwritten reading replaces the underlying reading.}
\psentry{ps_134-8_027-v0013}{13}{deletion}{\deleted{itself}}{}{The author cancels this material by hand; it leaves the corrected running text.}
\psapparatusend
\fi
\pssource{ps_134-8_028}{134-8/28}{https://grothendieck.umontpellier.fr/web/viewer.html?file=../134-8.pdf\#page=29}
\psnumber{ps_134-8_028-v0001}{1}
\psnumber{ps_134-8_028-v0002}{2}
\psnumber{ps_134-8_028-v0003}{3}
\psnumber{ps_134-8_028-v0004}{4}
\psnumber{ps_134-8_028-v0005}{5}
\psnumber{ps_134-8_028-v0006}{6}
\psnumber{ps_134-8_028-v0007}{7}
\psnumber{ps_134-8_028-v0008}{8}
\psnumber{ps_134-8_028-v0009}{9}
\psnumber{ps_134-8_028-v0010}{10}
%% BEGIN SOURCE ps_134-8_028
% Pursuing Stacks (cotes 134-2 to 134-8): archive cote 134-8, site page 28 (PDF page 29), typescript page 582.
% typescript. Draft compared: ps7.tex:1297-1363.
% AI-transcribed by Claude Opus 5.5 (claude-opus-5-5, adaptive thinking, effort high), Message Batches API, batch ps4 (msgbatch_012Aoc1NuskATPsyBeKSYK7J), 2026-11-12; reviewed against both supplied scan halves; see CHECK_DECISIONS.tsv.
\scanpage{134-8/28}
\mspage{582}
itself,  and (3) reads
\[ H_i(A,\underline{O}_A)= 0 \quad \text{for } i\hand{>}0 \ , \]
which is'nt always true.  (NB if it were for any commutative ring $\underline{O}_A$
on A,  it would be too for any abelian sheaf M on A,  as we see by \var{\xout{takin}}{}{typewriter deletion}{ps_134-8_028-v0001}
taking  $\underline{O}_A = \underline{Z}_A \hand{\oplus} M$ ,  hence A would be of \var{\deleted{an}}{}{deletion}{ps_134-8_028-v0002} homological dimension 0,\var{\deleted{\illegible{1}}}{}{deletion}{ps_134-8_028-v0003}
a drastic restriction indeed \var{\xout{beyond t}}{}{typewriter deletion}{ps_134-8_028-v0004} even when A has a final object.)

When however $\underline{O}_A$ is equally a constant sheaf of rings,  \var{\xout{then the}}{}{typewriter deletion}{ps_134-8_028-v0005}
say  $\underline{O}_A = k'_A$ ,  then the relation (3) holds for any projective Module
on A .  We need only check it for
\[ F = \underline{O}_A{}^{(a)} = k'{}^{(a)} \ , \]
with a in A ,  then (3) follows from
\[ \text{(4)} \qquad L_i f^{ab}_{o!}(M^{(a)}) = 0 \quad \text{for } i\hand{>}0 \ , \quad M \text{ in } (Ab) \ , \]
To chedk (4),  we take a proje\var{\over{\unsure{ct}}{ct}}{\hand{ct}}{writing over}{ps_134-8_028-v0006}ive resolution of $M^{(a)}$ ,  by using a
projective resolution L. of M in (Ab) and taking  $L.^{(a)}$  (using the
fact that the functor
\[ L \hand{\mapsto} L^{(a)} \ : \ (Ab) \hand{\to} A^{\wedge}_{ab} \]
is exact).  We then get
\[ Lf_{o!}(M^{(a)}) = f_{o!}(L.^{(a)}) = L.^{(b)} \hand{\simeq} M^{(b)} \quad , \qquad \text{\var{\xout{b=f(a)}}{}{typewriter deletion}{ps_134-8_028-v0007}} \]
where $b=f_o(a)$,and where the last equality stems from exactness of $L\hand{\mapsto} L^{(b)}$.
(NB  T\var{\over[typed]{h}{h}}{h}{typewriter overwriting}{ps_134-8_028-v0008}e relation (4) generalizes a standard acyclicity criterion in
the homology theory of discrete groups ...)

Thus we get finally that in case both Rings $\underline{O}_A$ and $\underline{O}_B$ are constant,
that formation of $Lf_!$  commutes to restriction of operator rings
(provided the Rings to which we are restricting are constant too  -
say they are just the absolute $\underline{Z}_A$ and $\underline{Z}_B$).  Presumably a little more
care should show the similar result for locally constant Rings\hand{.} \var{\deleted{-- but
it does'nt seem one can get \inserted{\deleted{rid of such an}} beyond this assumption.}}{}{deletion}{ps_134-8_028-v0009}

Reverting to the case (1) of a general map between ringed objects
in (Cat),  our unability,  for a given $\underline{O}_A$-Module F on A,  to define an
operation of $\underline{O}_B$ upon $f^{ab}_{o!}$  (while there is an operation of the ring
$\hand{\Gamma}(B,\underline{O}_B)$ upon it),  is tied up with this difficulty,  that formation
of  $f^{ab}_{o!}$ ,  and a fortiori of $f_{!\text{\var{\deleted{\unsure{*}}}{}{deletion}{ps_134-8_028-v0010}}}$ for arbitrary rings $\underline{O}_A$ and $\underline{O}_B$ , does
not commute to ``localization'' (as $f_*$ and $Rf_*$ does),  namely to base
changes of the type

\margin{28}

% Visual collation: both supplied image halves inspected; see CHECK_DECISIONS_134_7_8.tsv.

%% END SOURCE ps_134-8_028
\ifnum\EditionMode=1
\psapparatusstart
\psentry{ps_134-8_028-v0001}{1}{typewriter deletion}{\xout{takin}}{}{The typewriter strikes out this material; it leaves the corrected running text.}
\psentry{ps_134-8_028-v0002}{2}{deletion}{\deleted{an}}{}{The author cancels this material by hand; it leaves the corrected running text.}
\psentry{ps_134-8_028-v0003}{3}{deletion}{\deleted{\illegible{1}}}{}{The author cancels this material by hand; it leaves the corrected running text.}
\psentry{ps_134-8_028-v0004}{4}{typewriter deletion}{\xout{beyond t}}{}{The typewriter strikes out this material; it leaves the corrected running text.}
\psentry{ps_134-8_028-v0005}{5}{typewriter deletion}{\xout{then the}}{}{The typewriter strikes out this material; it leaves the corrected running text.}
\psentry{ps_134-8_028-v0006}{6}{writing over}{\over{\unsure{ct}}{ct}}{\hand{ct}}{The superposed handwritten reading replaces the underlying reading.}
\psentry{ps_134-8_028-v0007}{7}{typewriter deletion}{\xout{b=f(a)}}{}{The typewriter strikes out this material; it leaves the corrected running text.}
\psentry{ps_134-8_028-v0008}{8}{typewriter overwriting}{\over[typed]{h}{h}}{h}{The superposed typed reading replaces the underlying reading.}
\psentry{ps_134-8_028-v0009}{9}{deletion}{\deleted{-- but
it does'nt seem one can get \inserted{\deleted{rid of such an}} beyond this assumption.}}{}{The author cancels this material by hand; it leaves the corrected running text.}
\psentry{ps_134-8_028-v0010}{10}{deletion}{\deleted{\unsure{*}}}{}{The author cancels this material by hand; it leaves the corrected running text.}
\psapparatusend
\fi
\pssource{ps_134-8_029}{134-8/29}{https://grothendieck.umontpellier.fr/web/viewer.html?file=../134-8.pdf\#page=30}
\psnumber{ps_134-8_029-v0001}{1}
\psnumber{ps_134-8_029-v0002}{2}
\psnumber{ps_134-8_029-v0003}{3}
\psnumber{ps_134-8_029-v0004}{4}
\psnumber{ps_134-8_029-v0005}{5}
\psnumber{ps_134-8_029-v0006}{6}
\psnumber{ps_134-8_029-v0007}{7}
\psnumber{ps_134-8_029-v0008}{8}
\psnumber{ps_134-8_029-v0009}{9}
\psnumber{ps_134-8_029-v0010}{10}
\psnumber{ps_134-8_029-v0011}{11}
\psnumber{ps_134-8_029-v0012}{12}
%% BEGIN SOURCE ps_134-8_029
% Pursuing Stacks (cotes 134-2 to 134-8): archive cote 134-8, site page 29 (PDF page 30), typescript page 583.
% typescript. Draft compared: ps7.tex:1344-1410.
% AI-transcribed by Claude Opus 5.5 (claude-opus-5-5, adaptive thinking, effort high), Message Batches API, batch ps4 (msgbatch_012Aoc1NuskATPsyBeKSYK7J), 2026-11-12; reviewed against both supplied scan halves; see CHECK_DECISIONS.tsv.
\scanpage{134-8/29}
\mspage{583}

(5) \qquad $B_{\hand{/}b} \hand{\longrightarrow} B$ \quad ,

where b is a\var{\xout{n}}{}{typewriter deletion}{ps_134-8_029-v0001} given object in B, and $B_{\hand{/}b}$ designates as usual the cate-%
gory of all "objects over b" in B , i.e. of all arrows \var{\inserted[typed]{in B}}{in B}{typed insertion}{ps_134-8_029-v0002} with target b.
This gives rise to the cartesian square in (Cat)

\diagram{(6) \qquad \begin{tikzcd} A_{\hand{/}b} \arrow[color=blue,r] \arrow[color=blue,d] \& A \arrow[color=blue,d] \\ B_{\hand{/}b} \arrow[color=blue,r] \& B \end{tikzcd} \quad ,}

where $A_{\hand{/}b}$ is the category of all pairs
\[ (a,u) \text{ with } u : f_{o}(a) \hand{\longrightarrow} b \quad , \]
which may be identified equally with the category $A_{\hand{/}f_{o}^{*}(b)}$ . The
commutation property we have in mind is a tautology for $f_{*}$ , and it
f\hand{o}llows for $Rf_{*}$ , because the inverse image by $A_{\hand{/}b} \hand{\to} A$ of an in\hand{j}ecti-%
ve Module on A is an injective Module on $A_{\hand{/}b}$ . \var{\deleted{\unsure{(}}}{}{deletion}{ps_134-8_029-v0003}T\var{\over[typed]{h}{h}}{h}{typewriter overwriting}{ps_134-8_029-v0004}e latter fact is
true, more generally, for any \var{\xout{map}}{}{typewriter deletion}{ps_134-8_029-v0005} \var{\inserted[typed]{"localization map" , of the type}}{"localization map" , of the type}{typed insertion}{ps_134-8_029-v0006} $A_{\hand{/}X} \hand{\to} A$, with X in $A^{\wedge}$, i.e.
any map which is fibering with discrete fibers.\var{\deleted{\unsure{)}}}{}{deletion}{ps_134-8_029-v0007} Thus, the commutation
property for $Rf_{*}$ is valid \var{\over[typed]{f}{m}}{m}{typewriter overwriting}{ps_134-8_029-v0008}ore generally for any base change of the
type
\[ B_{\hand{/}Y} \hand{\longrightarrow} Y \quad , \]
with Y in $B^{\wedge}$ (not necessarily in \var{\over{B}{B}}{\hand{B}}{writing over}{ps_134-8_029-v0009}). More generally still, it can be
shown to hold for any map
\[ B' \hand{\longrightarrow} B \]
which is fibering (not necessarily with discrete fibers).

On the other side of the mirror, when taking $f_{!}$ and it's left
derived functor, already the former \var{\inserted[typed]{definitely}}{definitely}{typed insertion}{ps_134-8_029-v0010} does \uline{not} commute to base change
of the type (5) i.e. to"localization on the base" - something a little
hard to get accustomed to ! T$\overset{\illegible{1}}{\underset{\illegible{1}}{h}}$e base changes which will do here, are
those of the dual type

(5') \qquad ${}_{b\hand{\backslash}}B \hand{\longrightarrow} B$ \quad ,

we may call them maps of "colocalization\hand{"} on the base\var{\deleted{"}}{}{deletion}{ps_134-8_029-v0011} B. \var{\xout{Here}}{}{typewriter deletion}{ps_134-8_029-v0012} It gives
rise to a cartesian diagram in (Cat) dual to (6)

\diagram{(6') \qquad \begin{tikzcd} {}_{b\hand{\backslash}}A \arrow[color=blue,r] \arrow[color=blue,d] \& A \arrow[color=blue,d] \\ {}_{b\hand{\backslash}}B \arrow[color=blue,r] \& B \end{tikzcd} \quad ,}

% Hand-drawn arrows are marked blue; typed objects remain plain.
\margin{29}

% Visual collation: both supplied image halves inspected; see CHECK_DECISIONS_134_7_8.tsv.

%% END SOURCE ps_134-8_029
\ifnum\EditionMode=1
\psapparatusstart
\psentry{ps_134-8_029-v0001}{1}{typewriter deletion}{\xout{n}}{}{The typewriter strikes out this material; it leaves the corrected running text.}
\psentry{ps_134-8_029-v0002}{2}{typed insertion}{\inserted[typed]{in B}}{in B}{The typed insertion is incorporated at its indicated position; insertion brackets leave the running text.}
\psentry{ps_134-8_029-v0003}{3}{deletion}{\deleted{\unsure{(}}}{}{The author cancels this material by hand; it leaves the corrected running text.}
\psentry{ps_134-8_029-v0004}{4}{typewriter overwriting}{\over[typed]{h}{h}}{h}{The superposed typed reading replaces the underlying reading.}
\psentry{ps_134-8_029-v0005}{5}{typewriter deletion}{\xout{map}}{}{The typewriter strikes out this material; it leaves the corrected running text.}
\psentry{ps_134-8_029-v0006}{6}{typed insertion}{\inserted[typed]{"localization map" , of the type}}{"localization map" , of the type}{The typed insertion is incorporated at its indicated position; insertion brackets leave the running text.}
\psentry{ps_134-8_029-v0007}{7}{deletion}{\deleted{\unsure{)}}}{}{The author cancels this material by hand; it leaves the corrected running text.}
\psentry{ps_134-8_029-v0008}{8}{typewriter overwriting}{\over[typed]{f}{m}}{m}{The superposed typed reading replaces the underlying reading.}
\psentry{ps_134-8_029-v0009}{9}{writing over}{\over{B}{B}}{\hand{B}}{The superposed handwritten reading replaces the underlying reading.}
\psentry{ps_134-8_029-v0010}{10}{typed insertion}{\inserted[typed]{definitely}}{definitely}{The typed insertion is incorporated at its indicated position; insertion brackets leave the running text.}
\psentry{ps_134-8_029-v0011}{11}{deletion}{\deleted{"}}{}{The author cancels this material by hand; it leaves the corrected running text.}
\psentry{ps_134-8_029-v0012}{12}{typewriter deletion}{\xout{Here}}{}{The typewriter strikes out this material; it leaves the corrected running text.}
\psapparatusend
\fi
\pssource{ps_134-8_030}{134-8/30}{https://grothendieck.umontpellier.fr/web/viewer.html?file=../134-8.pdf\#page=31}
\psnumber{ps_134-8_030-v0001}{1}
\psnumber{ps_134-8_030-v0002}{2}
\psnumber{ps_134-8_030-v0003}{3}
\psnumber{ps_134-8_030-v0004}{4}
\psnumber{ps_134-8_030-v0005}{5}
\psnumber{ps_134-8_030-v0006}{6}
\psnumber{ps_134-8_030-v0007}{7}
\psnumber{ps_134-8_030-v0008}{8}
\psnumber{ps_134-8_030-v0009}{9}
\psnumber{ps_134-8_030-v0010}{10}
\psnumber{ps_134-8_030-v0011}{11}
\psnumber{ps_134-8_030-v0012}{12}
\psnumber{ps_134-8_030-v0013}{13}
\psnumber{ps_134-8_030-v0014}{14}
\psnumber{ps_134-8_030-v0015}{15}
%% BEGIN SOURCE ps_134-8_030
% Pursuing Stacks (cotes 134-2 to 134-8): archive cote 134-8, site page 30 (PDF page 31), typescript page 584.
% typescript. Draft compared: ps7.tex:1392-1455.
% AI-transcribed by Claude Opus 5.5 (claude-opus-5-5, adaptive thinking, effort high), Message Batches API, batch ps4 (msgbatch_012Aoc1NuskATPsyBeKSYK7J), 2026-11-12; reviewed against both supplied scan halves; see CHECK_DECISIONS.tsv.
\scanpage{134-8/30}
\mspage{584}

where now ${}_{b\hand{\backslash}}A$  is the category of pairs
\[ (a,u) \ \text{ with } \ u: b \hand{\to} f_{o}(a) \quad . \]
\var{\over{\unsure{We}}{We}}{\hand{We}}{writing over}{ps_134-8_030-v0001}'ll have to assume npw that both sheaves of rings $\underline{O}_A$ , $\underline{O}_B$  are con-%
stant,  and correspond to the same ring k.  \var{\xout{The commutation of f}}{}{typewriter deletion}{ps_134-8_030-v0002} Thus
Modules on A and B are just contravariant functors from these catego-%
ries to the category  $Ab_k$  of \var{\xout{all}}{}{typewriter deletion}{ps_134-8_030-v0003} k-modules,  and accordingly $f_!$
(\var{\xout{dual}}{}{typewriter deletion}{ps_134-8_030-v0004} \var{\inserted[typed]{left adjoint \xout{to}}}{left adjoint }{typed insertion}{ps_134-8_030-v0005} to the composition functor $f^*$ )  may be computed by a well known
formula,  involving direct limits on the categories ${}_{b\hand{\backslash}}A$ :
\[ (7) \qquad f_!(F)(b) \hand{\simeq} \underset{b\hand{\backslash} A}{\underset{\hand{\longrightarrow}}{\lim}} F(a) \qquad , \]
where the limit in the second member is relative to the \var{\xout{funct}}{}{typewriter deletion}{ps_134-8_030-v0006} composi-%
tion
\[ {}_{b\hand{\backslash}}A \hand{\longrightarrow} A \overset{\ F\ }{\hand{\longrightarrow}} Ab_k \quad . \]
(NB This formula is dual to the formula for $f_*$ ,  the right adjoint of
$f^*$ ,  closer to intuition  - to mine at any rate -  because the $\underset{\hand{\longleftarrow}}{\lim}$ in
\[ f_*(F)(b) \hand{\simeq} \underset{A_{\hand{/}b}}{\underset{\hand{\longleftarrow}}{\lim}} F(a) \]
may be ``visualized'' as \var{\deleted{\xout{\illegible{1}sections\illegible{1}}}}{}{deletion}{ps_134-8_030-v0007}  the \var{\xout{set of}}{}{typewriter deletion}{ps_134-8_030-v0008} set of ``sections'' of
F over $A_{/b}$ \var{\over{\unsure{,}}{.}}{\hand{.}}{writing over}{ps_134-8_030-v0009}) From this formula (7) follows at once commutation of
$f_!$  with colocalization\var{\xout{t}}{}{typewriter deletion}{ps_134-8_030-v0010}.  To get the corresponding result for $Lf_!$ ,
we have only to use the fact that the relevant inverse image functor
(corresponding to  ${}_{b\hand{\backslash}}A \hand{\to} A$ ) transforms projective\var{\deleted{\xout{\illegible{1}}}}{}{deletion}{ps_134-8_030-v0011} modules into
% NOTE: from here to the foot of the page the text is typed on a strip pasted over the page; the older typed text underneath shows through faintly and is not transcribed
projective modules.  As the map  ${}_{b\hand{\backslash}}B \hand{\longrightarrow} B$  is a \uline{cofibering functor}
\uline{with discrete fibers},  so is the map
\[ h: A' = {}_{b\hand{\backslash}}A \hand{\longrightarrow} A \]
deduced by base change.  Now for any \var{\inserted{such}}{\hand{such}}{insertion}{ps_134-8_030-v0012} functor  $h: A' \hand{\to} A$ between small
categories,  the inverse image functor
\[ h^* : A^{\overset{\hand{\illegible{1}}}{\wedge}} \hand{\longrightarrow} A'^{\wedge} \ , \quad \text{or} \quad h^*_k : A^{\wedge}_k \hand{\to} A'^{\wedge}_k \]
carries \var{\inserted[typed]{indeed}}{indeed}{typed insertion}{ps_134-8_030-v0013} projectives into projectives.  This statement is dual formally
to the corresponding statement for injectives,  valid when we make on
h the dual assumption of being \uline{fibering with discrete fibers}  -  in the
latter case the \hand{(}well known\hand{)} proof comes out formally from the fact
that the left adjoint functor $h_!$ or $h^k_!$ carries monomorphisms into
monomorphisms \var{\deleted{i.e. is exact}}{}{deletion}{ps_134-8_030-v0014}  -  a fact which we used in section 135
\margin{\tikz{\draw[thin] (0,0)--(1.1,0);}}in the form  $Lh^k_! = h^k_!$ (in case $k=\underline{Z}$).  In the present case the proof is
essentially the dual one  -  \var{\over{as}{as}}{\hand{as}}{writing over}{ps_134-8_030-v0015} a matter of fact, as there are enough

\margin{30}

% Visual collation: both supplied image halves inspected; see CHECK_DECISIONS_134_7_8.tsv.

%% END SOURCE ps_134-8_030
\ifnum\EditionMode=1
\psapparatusstart
\psentry{ps_134-8_030-v0001}{1}{writing over}{\over{\unsure{We}}{We}}{\hand{We}}{The superposed handwritten reading replaces the underlying reading.}
\psentry{ps_134-8_030-v0002}{2}{typewriter deletion}{\xout{The commutation of f}}{}{The typewriter strikes out this material; it leaves the corrected running text.}
\psentry{ps_134-8_030-v0003}{3}{typewriter deletion}{\xout{all}}{}{The typewriter strikes out this material; it leaves the corrected running text.}
\psentry{ps_134-8_030-v0004}{4}{typewriter deletion}{\xout{dual}}{}{The typewriter strikes out this material; it leaves the corrected running text.}
\psentry{ps_134-8_030-v0005}{5}{typed insertion}{\inserted[typed]{left adjoint \xout{to}}}{left adjoint }{The typed insertion is incorporated at its indicated position; insertion brackets leave the running text.}
\psentry{ps_134-8_030-v0006}{6}{typewriter deletion}{\xout{funct}}{}{The typewriter strikes out this material; it leaves the corrected running text.}
\psentry{ps_134-8_030-v0007}{7}{deletion}{\deleted{\xout{\illegible{1}sections\illegible{1}}}}{}{The author cancels this material by hand; it leaves the corrected running text.}
\psentry{ps_134-8_030-v0008}{8}{typewriter deletion}{\xout{set of}}{}{The typewriter strikes out this material; it leaves the corrected running text.}
\psentry{ps_134-8_030-v0009}{9}{writing over}{\over{\unsure{,}}{.}}{\hand{.}}{The superposed handwritten reading replaces the underlying reading.}
\psentry{ps_134-8_030-v0010}{10}{typewriter deletion}{\xout{t}}{}{The typewriter strikes out this material; it leaves the corrected running text.}
\psentry{ps_134-8_030-v0011}{11}{deletion}{\deleted{\xout{\illegible{1}}}}{}{The author cancels this material by hand; it leaves the corrected running text.}
\psentry{ps_134-8_030-v0012}{12}{insertion}{\inserted{such}}{\hand{such}}{The handwritten insertion is incorporated at its indicated position; insertion brackets leave the running text.}
\psentry{ps_134-8_030-v0013}{13}{typed insertion}{\inserted[typed]{indeed}}{indeed}{The typed insertion is incorporated at its indicated position; insertion brackets leave the running text.}
\psentry{ps_134-8_030-v0014}{14}{deletion}{\deleted{i.e. is exact}}{}{The author cancels this material by hand; it leaves the corrected running text.}
\psentry{ps_134-8_030-v0015}{15}{writing over}{\over{as}{as}}{\hand{as}}{The superposed handwritten reading replaces the underlying reading.}
\psapparatusend
\fi
\pssource{ps_134-8_031}{134-8/31}{https://grothendieck.umontpellier.fr/web/viewer.html?file=../134-8.pdf\#page=32}
\psnumber{ps_134-8_031-v0001}{1}
\psnumber{ps_134-8_031-v0002}{2}
\psnumber{ps_134-8_031-v0003}{3}
\psnumber{ps_134-8_031-v0004}{4}
\psnumber{ps_134-8_031-v0005}{5}
\psnumber{ps_134-8_031-v0006}{6}
\psnumber{ps_134-8_031-v0007}{7}
%% BEGIN SOURCE ps_134-8_031
% Pursuing Stacks (cotes 134-2 to 134-8): archive cote 134-8, site page 31 (PDF page 32), typescript page 585.
% typescript. Draft compared: ps7.tex:1441-1512.
% AI-transcribed by Claude Opus 5.5 (claude-opus-5-5, adaptive thinking, effort high), Message Batches API, batch ps4 (msgbatch_012Aoc1NuskATPsyBeKSYK7J), 2026-11-12; reviewed against both supplied scan halves; see CHECK_DECISIONS.tsv.
\scanpage{134-8/31}
\mspage{585}
projectives,  the statement about \var{\xout{projectives}}{}{typewriter deletion}{ps_134-8_031-v0001} $h^{*}$ or $h^{*}_{k}$ taking projecti-%
ves into projectives is equivalent with the right adjoint $h_{*}$ or $h^{k}_{*}$
(for sheaves of s ets,  resp. sheaves of k-modules) transforming epimor-%
phisms into epimorphisms,  which in the case of k-modules can be
written equally under the equivalent form
\[ (8) \qquad R^{i}h_{*} = 0 \ \text{ for } \ i \hand{>} 0 \ . \]
Now th\var{\over[typed]{e}{i}}{i}{typewriter overwriting}{ps_134-8_031-v0002}s exactness property for $h_{*}$ ,  in the case when h is cofibering
with discrete fibers, is I guess well known (it is well known to me
at any rate),  and comes from the specific computation of $h_{*}(F)$ for
F in $A'^{\wedge}$,  valid whenever h is \uline{cofibering} (with arbitary fibers)
\[ (9) \qquad h_{*}(F)(x) \hand{\simeq} \hand{\Gamma}(A'_{x}, F\hand{|}A'_{x}) \qquad \text{for x in A ,} \]
where $A'_{x}$ is the fiber of A' over x (a category not to be cofused with
$A'_{\hand{/}x}$ ,  the two being closely related however ...).  In case the fibers
of h are discrete,  the right hand side of (9) may be written as a
product,  hence the formula
\[ (10) \qquad h_{*}(F)(x) \hand{\simeq} \hand{\prod}_{x' \text{ in } A'_{x}} F(x') \ , \]
(which may be viewed as the dual of the formula (11) p. 564).  As in
the category of sets (and hence also in  $Ab_{k}$ ) a product of epimorphis-%
ms is again an epimorphism,  the result we want follows indeed.

\uline{Remark\hand{s}}1)The results just given,  as well as their proof,  concerning
inverse images of injectives or projectives,  are valid not only in
the case of a \var{\xout{constant sheaf}}{}{typewriter deletion}{ps_134-8_031-v0003} common constant sheaf of rings on A and
A',  but more generally for any sheaf of rings $\underline{O}_{A}$ on A,  when taking
on A' the ``induced'' sheaf of rings
\[ \hand{(11)} \qquad \underline{O}_{A'} = h^{*}(\underline{O}_{A}) \quad . \]
H\var{\over{e}{o}}{\hand{o}}{writing over}{ps_134-8_031-v0004}wever,  it does'nt seem that the result about commutation of $Lf_{!}$
to colocalization is valid under the corresponding assumption
$\underline{O}_{A} = f^{*}(\underline{O}_{B})$ ,  without assuming moreover $\underline{O}_{B}$ \var{\xout{itself}}{}{typewriter deletion}{ps_134-8_031-v0005} to be constant
(hence $\underline{O}_{A}$ too),  because already for the functor $f_{!}$ itself for M$\overset{\hand{\scriptscriptstyle\smile}}{\mathrm{o}}$dules\hand{,}
it does'nt seem that commutation will hold.

2)  I should correct a silly mistake I made at the very beginning
of this section,  when rashly stating that the functor $Lf_{!}$  may be
defined for \uline{any} map \var{\inserted[typed]{(1)}}{(1)}{typed insertion}{ps_134-8_031-v0006} between ringed objects in (Cat).  I was thinking
of the fact that for any ri\var{\over{\unsure{gg}}{ng}}{\hand{ng}}{writing over}{ps_134-8_031-v0007}ed object $(A,\underline{O}_{A})$  in (Cat),  there are
indeed enough projectives in the category of \hand{$\underline{O}_{A}$}-Modules  --
\margin{31}

% Visual collation: both supplied image halves inspected; see CHECK_DECISIONS_134_7_8.tsv.

%% END SOURCE ps_134-8_031
\ifnum\EditionMode=1
\psapparatusstart
\psentry{ps_134-8_031-v0001}{1}{typewriter deletion}{\xout{projectives}}{}{The typewriter strikes out this material; it leaves the corrected running text.}
\psentry{ps_134-8_031-v0002}{2}{typewriter overwriting}{\over[typed]{e}{i}}{i}{The superposed typed reading replaces the underlying reading.}
\psentry{ps_134-8_031-v0003}{3}{typewriter deletion}{\xout{constant sheaf}}{}{The typewriter strikes out this material; it leaves the corrected running text.}
\psentry{ps_134-8_031-v0004}{4}{writing over}{\over{e}{o}}{\hand{o}}{The superposed handwritten reading replaces the underlying reading.}
\psentry{ps_134-8_031-v0005}{5}{typewriter deletion}{\xout{itself}}{}{The typewriter strikes out this material; it leaves the corrected running text.}
\psentry{ps_134-8_031-v0006}{6}{typed insertion}{\inserted[typed]{(1)}}{(1)}{The typed insertion is incorporated at its indicated position; insertion brackets leave the running text.}
\psentry{ps_134-8_031-v0007}{7}{writing over}{\over{\unsure{gg}}{ng}}{\hand{ng}}{The superposed handwritten reading replaces the underlying reading.}
\psapparatusend
\fi
\pssource{ps_134-8_032}{134-8/32}{https://grothendieck.umontpellier.fr/web/viewer.html?file=../134-8.pdf\#page=33}
\psnumber{ps_134-8_032-v0001}{1}
\psnumber{ps_134-8_032-v0002}{2}
\psnumber{ps_134-8_032-v0003}{3}
\psnumber{ps_134-8_032-v0004}{4}
\psnumber{ps_134-8_032-v0005}{5}
\psnumber{ps_134-8_032-v0006}{6}
\psnumber{ps_134-8_032-v0007}{7}
\psnumber{ps_134-8_032-v0008}{8}
\psnumber{ps_134-8_032-v0009}{9}
\psnumber{ps_134-8_032-v0010}{10}
\psnumber{ps_134-8_032-v0011}{11}
\psnumber{ps_134-8_032-v0012}{12}
\psnumber{ps_134-8_032-v0013}{13}
\psnumber{ps_134-8_032-v0014}{14}
\psnumber{ps_134-8_032-v0015}{15}
\psnumber{ps_134-8_032-v0016}{16}
\psnumber{ps_134-8_032-v0017}{17}
\psnumber{ps_134-8_032-v0018}{18}
\psnumber{ps_134-8_032-v0019}{19}
\psnumber{ps_134-8_032-v0020}{20}
\psnumber{ps_134-8_032-v0021}{21}
\psnumber{ps_134-8_032-v0022}{22}
\psnumber{ps_134-8_032-v0023}{23}
\psnumber{ps_134-8_032-v0024}{24}
\psnumber{ps_134-8_032-v0025}{25}
%% BEGIN SOURCE ps_134-8_032
% Pursuing Stacks (cotes 134-2 to 134-8): archive cote 134-8, site page 32 (PDF page 33), typescript page 586.
% typescript. Draft compared: ps7.tex:1496-1554.
% AI-transcribed by Claude Opus 5.5 (claude-opus-5-5, adaptive thinking, effort high), Message Batches API, batch ps4 (msgbatch_012Aoc1NuskATPsyBeKSYK7J), 2026-11-12; reviewed against both supplied scan halves; see CHECK_DECISIONS.tsv.
\scanpage{134-8/32}
\mspage{586}

thus the Modules
(12) \qquad $\underline{O}_A^{\;(a)}$ \quad , for a in A \quad ,

are clearly projective, and the\var{\over[typed]{i}{r}}{r}{typewriter overwriting}{ps_134-8_032-v0001}\var{\over[typed]{r}{e}}{e}{typewriter overwriting}{ps_134-8_032-v0002} are ``sufficiently many''. Thus,
any \var{\inserted[typed]{additive}}{additive}{typed insertion}{ps_134-8_032-v0003} functor from Mod($\underline{O}_A$) to an abelian category admits a total left
derived functor. However, it is not always true, for a map (1)
of small ringed categories, that the corresponding inverse image
functor for Modules
\[ G \hand{\longmapsto} f^{*}(G)= f_o^{\;*}(G) \otimes_{\var{\over{\unsure{f}}{f}}{\hand{f}}{writing over}{ps_134-8_032-v0004}_o^{*}(\underline{O}_B)} \underline{O}_A \]
admits \var{\xout{always}}{}{typewriter deletion}{ps_134-8_032-v0005} a left adjoint $f_{!}$ , or what amounts to the same\var{\xout{\illegible{1}}}{}{typewriter deletion}{ps_134-8_032-v0006}, that
this functor (which is right exact) is left exact and commutes to
\var{\xout{arbitrary}}{}{typewriter deletion}{ps_134-8_032-v0007} \var{\inserted[typed]{small}}{small}{typed insertion}{ps_134-8_032-v0008} products. It is'nt even true necessarily, when we assume
A and B to be the final category ! Left exactness of $f^{\hand{*}}$ just means
flatness of $\underline{O}_A$ over $\underline{O}_B$ , i.e. that for any a in A, $\underline{O}_A(a)$ is flat as
a\var{\xout{s}}{}{typewriter deletion}{ps_134-8_032-v0009} module over $\underline{O}_B(b)$, where b=f(a). As for commutation to small
products, it amounts (\var{\xout{when}}{}{typewriter deletion}{ps_134-8_032-v0010} \var{\inserted[typed]{together with}}{together with}{typed insertion}{ps_134-8_032-v0011} the first condition)\var{\xout{is satisfied)}}{}{typewriter deletion}{ps_134-8_032-v0012} to the
still more ex\var{\over{a}{a}}{\hand{a}}{writing over}{ps_134-8_032-v0013}cting condition that $\underline{O}_A(a)$ should be a \uline{projective module
of finite type over} $\underline{O}_B(b)$, for any a in A. This condition is so close
to the condition $\underline{O}_{\var{\over{A}{A}}{\hand{A}}{writing over}{ps_134-8_032-v0014}} = f_o^{\;*}(\underline{\var{\over[typed]{A}{O}}{O}{typewriter overwriting}{ps_134-8_032-v0015}}_B)$ \quad (already considered in (11) above),
that for a bird's eye view as we are aiming at here, we may as well
assume this slightly stronger condition\hand{!} Anyhow, as noticed before,
to get \var{\xout{compatibility of}}{}{typewriter deletion}{ps_134-8_032-v0016} commutation of $f_{!}$ with restriction of scalars
and with colocalization, even this assumption is'nt enough apparently\hand{,}
and we'll have to assume moreover that $\underline{O}_B$ (hence also $\underline{O}_A$) is constant,
or for the least, locally constant.

This brings us back to the case when a fixed ring k is given,
and when we are working with categories of presheaves of k-modules
- a situation studied at length in chapter V. In case the target
category B is the final category, hence L$f_{!}$ is just (absolute) total
homology of A, this then may be computed nicely, using an ``integra-%
tor'' for \hand{A} , namely a projective resolution of $k_{A^o}$ in $(A^o)^{\wedge}_{k}$ . This
is more or less where we ended up \var{\inserted{with}}{\hand{with}}{insertion}{ps_134-8_032-v0017} by the end of chapter V, when \var{\over{d}{d}}{\hand{d}}{writing over}{ps_134-8_032-v0018}eve-%
\var{\xout{were} lopping}{lopping}{line continuation spacing}{ps_134-8_032-v0019} a nicely autodual homology-cohomology set-up, \var{\xout{where}}{}{typewriter deletion}{ps_134-8_032-v0020}
replacing the category of k-modules $\mathrm{Ab}_k$ by a more or less arbitrary
abelian category. I was about to \var{\deleted{\xout{carry t\unsure{rh|hr}}}}{}{deletion}{ps_134-8_032-v0021} go on and carry t\var{\over{\unsure{h}}{h}}{\hand{h}}{writing over}{ps_134-8_032-v0022}\var{\over[typed]{h}{r}}{r}{typewriter overwriting}{ps_134-8_032-v0023}ough
a similar treatment in the ``r\var{\over{\unsure{ee}}{el}}{\hand{el}}{writing over}{ps_134-8_032-v0024}ative'' case, namely for an arbitrary
map f in (Cat) \var{\deleted{\unsure{,}}}{}{deletion}{ps_134-8_032-v0025}\hand{(}but then I got caught unsuspectingly by that unen-%
ding digression on schematic homotopy types, finally making up a whole
chapter by itself\hand{)}. Maybe it is still worthwhile to come back to this\hand{,}

\margin{32}

% Visual collation: both supplied image halves inspected; see CHECK_DECISIONS_134_7_8.tsv.

%% END SOURCE ps_134-8_032
\ifnum\EditionMode=1
\psapparatusstart
\psentry{ps_134-8_032-v0001}{1}{typewriter overwriting}{\over[typed]{i}{r}}{r}{The superposed typed reading replaces the underlying reading.}
\psentry{ps_134-8_032-v0002}{2}{typewriter overwriting}{\over[typed]{r}{e}}{e}{The superposed typed reading replaces the underlying reading.}
\psentry{ps_134-8_032-v0003}{3}{typed insertion}{\inserted[typed]{additive}}{additive}{The typed insertion is incorporated at its indicated position; insertion brackets leave the running text.}
\psentry{ps_134-8_032-v0004}{4}{writing over}{$\over{\unsure{f}}{f}$}{$\hand{f}$}{The superposed handwritten reading replaces the underlying reading.}
\psentry{ps_134-8_032-v0005}{5}{typewriter deletion}{\xout{always}}{}{The typewriter strikes out this material; it leaves the corrected running text.}
\psentry{ps_134-8_032-v0006}{6}{typewriter deletion}{\xout{\illegible{1}}}{}{The typewriter strikes out this material; it leaves the corrected running text.}
\psentry{ps_134-8_032-v0007}{7}{typewriter deletion}{\xout{arbitrary}}{}{The typewriter strikes out this material; it leaves the corrected running text.}
\psentry{ps_134-8_032-v0008}{8}{typed insertion}{\inserted[typed]{small}}{small}{The typed insertion is incorporated at its indicated position; insertion brackets leave the running text.}
\psentry{ps_134-8_032-v0009}{9}{typewriter deletion}{\xout{s}}{}{The typewriter strikes out this material; it leaves the corrected running text.}
\psentry{ps_134-8_032-v0010}{10}{typewriter deletion}{\xout{when}}{}{The typewriter strikes out this material; it leaves the corrected running text.}
\psentry{ps_134-8_032-v0011}{11}{typed insertion}{\inserted[typed]{together with}}{together with}{The typed insertion is incorporated at its indicated position; insertion brackets leave the running text.}
\psentry{ps_134-8_032-v0012}{12}{typewriter deletion}{\xout{is satisfied)}}{}{The typewriter strikes out this material; it leaves the corrected running text.}
\psentry{ps_134-8_032-v0013}{13}{writing over}{\over{a}{a}}{\hand{a}}{The superposed handwritten reading replaces the underlying reading.}
\psentry{ps_134-8_032-v0014}{14}{writing over}{$\over{A}{A}$}{$\hand{A}$}{The superposed handwritten reading replaces the underlying reading.}
\psentry{ps_134-8_032-v0015}{15}{typewriter overwriting}{$\over[typed]{A}{O}$}{$O$}{The superposed typed reading replaces the underlying reading.}
\psentry{ps_134-8_032-v0016}{16}{typewriter deletion}{\xout{compatibility of}}{}{The typewriter strikes out this material; it leaves the corrected running text.}
\psentry{ps_134-8_032-v0017}{17}{insertion}{\inserted{with}}{\hand{with}}{The handwritten insertion is incorporated at its indicated position; insertion brackets leave the running text.}
\psentry{ps_134-8_032-v0018}{18}{writing over}{\over{d}{d}}{\hand{d}}{The superposed handwritten reading replaces the underlying reading.}
\psentry{ps_134-8_032-v0019}{19}{line continuation spacing}{\xout{were} lopping}{lopping}{Omitting the cancelled were must join lopping to the preceding visible deve- without a residual word space. The source hyphen is retained;diplomatic placement remains unchanged.}
\psentry{ps_134-8_032-v0020}{20}{typewriter deletion}{\xout{where}}{}{The typewriter strikes out this material; it leaves the corrected running text.}
\psentry{ps_134-8_032-v0021}{21}{deletion}{\deleted{\xout{carry t\unsure{rh|hr}}}}{}{The author cancels this material by hand; it leaves the corrected running text.}
\psentry{ps_134-8_032-v0022}{22}{writing over}{\over{\unsure{h}}{h}}{\hand{h}}{The superposed handwritten reading replaces the underlying reading.}
\psentry{ps_134-8_032-v0023}{23}{typewriter overwriting}{\over[typed]{h}{r}}{r}{The superposed typed reading replaces the underlying reading.}
\psentry{ps_134-8_032-v0024}{24}{writing over}{\over{\unsure{ee}}{el}}{\hand{el}}{The superposed handwritten reading replaces the underlying reading.}
\psentry{ps_134-8_032-v0025}{25}{deletion}{\deleted{\unsure{,}}}{}{The author cancels this material by hand; it leaves the corrected running text.}
\psapparatusend
\fi
\pssource{ps_134-8_033}{134-8/33}{https://grothendieck.umontpellier.fr/web/viewer.html?file=../134-8.pdf\#page=34}
\psnumber{ps_134-8_033-v0001}{1}
\psnumber{ps_134-8_033-v0002}{2}
\psnumber{ps_134-8_033-v0003}{3}
\psnumber{ps_134-8_033-v0004}{4}
\psnumber{ps_134-8_033-v0005}{5}
\psnumber{ps_134-8_033-v0006}{6}
\psnumber{ps_134-8_033-v0007}{7}
\psnumber{ps_134-8_033-v0008}{8}
\psnumber{ps_134-8_033-v0009}{9}
\psnumber{ps_134-8_033-v0010}{10}
\psnumber{ps_134-8_033-v0011}{11}
\psnumber{ps_134-8_033-v0012}{12}
\psnumber{ps_134-8_033-v0013}{13}
%% BEGIN SOURCE ps_134-8_033
% Pursuing Stacks (cotes 134-2 to 134-8): archive cote 134-8, site page 33 (PDF page 34), typescript page 587.
% typescript. Draft compared: ps7.tex:1543-1608.
% AI-transcribed by Claude Opus 5.5 (claude-opus-5-5, adaptive thinking, effort high), Message Batches API, batch ps4 (msgbatch_012Aoc1NuskATPsyBeKSYK7J), 2026-11-12; reviewed against both supplied scan halves; see CHECK_DECISIONS.tsv.
\scanpage{134-8/33}
\mspage{587}
without necessarily grinding through a complete formulaire for the
five \var{\inserted[typed]{main}}{main}{typed insertion}{ps_134-8_033-v0001} operations we got so far (namely $Lf_{!}$,$f^{*}$,$Rf_{*}$, $\overset{L}{\otimes}$ , $R\underline{Hom}$ ). Not
later than two pages ago or so, we were faced again with two visibly
dual st\var{\over{\unsure{at}}{at}}{\hand{at}}{writing over}{ps_134-8_033-v0002}em\var{\over{\unsure{ents}}{ents}}{\hand{ents}}{writing over}{ps_134-8_033-v0003}, one about $\var{\over[typed]{R}{L}}{L}{typewriter overwriting}{ps_134-8_033-v0004}f_{!}$ , the other about $Rf_{*}$ - and feeling
silly not to be able to merely deduce one from the other !

29.10.

139. Last night I still pondered a little about the $\var{\over[typed]{R}{L}}{L}{typewriter overwriting}{ps_134-8_033-v0005}f_{!}$ operation,
and did some more reading in the notes of chapter V of about three
months ago, which had been getting a little distant in my mind. In
those notes a great deal of emphasis goes with the notions of integra-%
tors and cointegrators - as a matter of fact, \var{\xout{it was}}{}{typewriter deletion}{ps_134-8_033-v0006} that whole
chapter sprung from an attempt to understand the meaning \var{\xout{and role}}{}{typewriter deletion}{ps_134-8_033-v0007} of
certain ``standard complexes'' associated with standard test categories
such as \hand{$\Delta$}, which then led us to the notion\var{\xout{s}}{}{typewriter deletion}{ps_134-8_033-v0008} of an integrator (via the
intermediate one of an ``abelianizator''). This in turn brought us to
\margin{\textcircled{$*_k$}}the $*_{k}$ - formalism, expressing most conveniently the relationship
between categories such as $A^{\wedge}_{k}$ and $(A^{o})^{\wedge}_{k}$ (k any commutative ring),
and it's various avatars. I was so pleased with this formalism and it'%
s ``computational'' flavor, that in my enthousiasm I subsumed under it
the dual treatment of homology and cohomology for an\var{\xout{category (with}}{}{typewriter deletion}{ps_134-8_033-v0009}
object A of (Cat) (with coefficients in a complex of \hand{$\mathcal{M}$}-valued
presheaves,\hand{$\mathcal{M}$} being an abelian category satisfying some mild conditi-%
ons), in section 108, as a particular case of the total derived
functors $F.\overset{L}{*}_{k}L^{\hand{\prime}}.$ and $R\mathrm{Hom}_{k}(L.,K^{\cdot})$ ( with values in the derived
category of \hand{$\mathcal{M}$} ). T$\overset{\illegible{1}}{\underset{\illegible{1}}{h}}$e main point here was using \uline{projective resolutions}
of the argument \hand{$L.,L^{\hand{\prime}}.$} in $B^{\wedge}_{k}$ or $A^{\wedge}_{k}$ (where $B=A^{o}$), which in \var{\xout{most cases}}{}{typewriter deletion}{ps_134-8_033-v0010}
the most important case was just the constant sheaf of rings $k_{B}$ or $k_{A}$
- rather than resolve the argument $F.$ (projectively) or $K^{\cdot}$ (injective-%
ly), namely the coefficients for homology or cohomology\hand{.}\var{\deleted{respectively.}}{}{deletion}{ps_134-8_033-v0011}
It was the enthousiasm of the adept of a game he just discovered -
I was going to try it out for the next step, namely relative homology
and cohomology $Lf_{!}$ and $Rf_{*}$ , with f any map in (Cat)\var{\deleted{,}}{}{deletion}{ps_134-8_033-v0012} - but then I
got caught by the more fascinating schematization game. Coming back
now upon the rather routine matter of looking up a comprehensive
mutually dual treatment of $Lf_{!}$ and $Rf_{*}$ , it does'nt seem that the
formalism of integrators and cointegrators is going to be of much help.
To be more specific, in order to compute (or simply define) $Lf_{!}(F.)$
\var{\xout{or $Rf_{*}$}}{}{typewriter deletion}{ps_134-8_033-v0013}

\margin{33}

% Visual collation: both supplied image halves inspected; see CHECK_DECISIONS_134_7_8.tsv.

%% END SOURCE ps_134-8_033
\ifnum\EditionMode=1
\psapparatusstart
\psentry{ps_134-8_033-v0001}{1}{typed insertion}{\inserted[typed]{main}}{main}{The typed insertion is incorporated at its indicated position; insertion brackets leave the running text.}
\psentry{ps_134-8_033-v0002}{2}{writing over}{\over{\unsure{at}}{at}}{\hand{at}}{The superposed handwritten reading replaces the underlying reading.}
\psentry{ps_134-8_033-v0003}{3}{writing over}{\over{\unsure{ents}}{ents}}{\hand{ents}}{The superposed handwritten reading replaces the underlying reading.}
\psentry{ps_134-8_033-v0004}{4}{typewriter overwriting}{$\over[typed]{R}{L}$}{$L$}{The superposed typed reading replaces the underlying reading.}
\psentry{ps_134-8_033-v0005}{5}{typewriter overwriting}{$\over[typed]{R}{L}$}{$L$}{The superposed typed reading replaces the underlying reading.}
\psentry{ps_134-8_033-v0006}{6}{typewriter deletion}{\xout{it was}}{}{The typewriter strikes out this material; it leaves the corrected running text.}
\psentry{ps_134-8_033-v0007}{7}{typewriter deletion}{\xout{and role}}{}{The typewriter strikes out this material; it leaves the corrected running text.}
\psentry{ps_134-8_033-v0008}{8}{typewriter deletion}{\xout{s}}{}{The typewriter strikes out this material; it leaves the corrected running text.}
\psentry{ps_134-8_033-v0009}{9}{typewriter deletion}{\xout{category (with}}{}{The typewriter strikes out this material; it leaves the corrected running text.}
\psentry{ps_134-8_033-v0010}{10}{typewriter deletion}{\xout{most cases}}{}{The typewriter strikes out this material; it leaves the corrected running text.}
\psentry{ps_134-8_033-v0011}{11}{deletion}{\deleted{respectively.}}{}{The author cancels this material by hand; it leaves the corrected running text.}
\psentry{ps_134-8_033-v0012}{12}{deletion}{\deleted{,}}{}{The author cancels this material by hand; it leaves the corrected running text.}
\psentry{ps_134-8_033-v0013}{13}{typewriter deletion}{\xout{or $Rf_{*}$}}{}{The typewriter strikes out this material; it leaves the corrected running text.}
\psapparatusend
\fi
\pssource{ps_134-8_034}{134-8/34}{https://grothendieck.umontpellier.fr/web/viewer.html?file=../134-8.pdf\#page=35}
\psnumber{ps_134-8_034-v0001}{1}
\psnumber{ps_134-8_034-v0002}{2}
\psnumber{ps_134-8_034-v0003}{3}
\psnumber{ps_134-8_034-v0004}{4}
\psnumber{ps_134-8_034-v0005}{5}
\psnumber{ps_134-8_034-v0006}{6}
\psnumber{ps_134-8_034-v0007}{7}
\psnumber{ps_134-8_034-v0008}{8}
\psnumber{ps_134-8_034-v0009}{9}
\psnumber{ps_134-8_034-v0010}{10}
\psnumber{ps_134-8_034-v0011}{11}
\psnumber{ps_134-8_034-v0012}{12}
\psnumber{ps_134-8_034-v0013}{13}
\psnumber{ps_134-8_034-v0014}{14}
\psnumber{ps_134-8_034-v0015}{15}
\psnumber{ps_134-8_034-v0016}{16}
\psnumber{ps_134-8_034-v0017}{17}
\psnumber{ps_134-8_034-v0018}{18}
\psnumber{ps_134-8_034-v0019}{19}
\psnumber{ps_134-8_034-v0020}{20}
\psnumber{ps_134-8_034-v0021}{21}
\psnumber{ps_134-8_034-v0022}{22}
\psnumber{ps_134-8_034-v0023}{23}
%% BEGIN SOURCE ps_134-8_034
% Pursuing Stacks (cotes 134-2 to 134-8): archive cote 134-8, site page 34 (PDF page 35), typescript page 588.
% typescript. Draft compared: ps7.tex:1594-1658.
% AI-transcribed by Claude Opus 5.5 (claude-opus-5-5, adaptive thinking, effort high), Message Batches API, batch ps4 (msgbatch_012Aoc1NuskATPsyBeKSYK7J), 2026-11-12; reviewed against both supplied scan halves; see CHECK_DECISIONS.tsv.
\scanpage{134-8/34}
\mspage{588}

or $Rf_{*}(K^{\bullet})$,  for a general map \var{\xout{f: A B in (Cat) and}}{}{typewriter deletion}{ps_134-8_034-v0001}
\[ f:\ A \hand{\longrightarrow} B \]
in (Cat),  and  \var{\xout{\unsure{L}}}{}{typewriter deletion}{ps_134-8_034-v0002} $F.\hand{\in} D^{-}(A^{\wedge}_{k})$,  $K^{\bullet}\hand{\in} D^{+}(A^{\wedge}_{k})$,  I do not see any means of
\margin{\textcircled{L.}}bypassing projective resolutions of $\var{\over{\unsure{F|L}}{L}}{\hand{L}}{writing over}{ps_134-8_034-v0003}.$ , injective on\var{\over{\unsure{e}}{e}}{\hand{e}}{writing over}{ps_134-8_034-v0004}s of $K^{\bullet}$.  Taking
the more familiar case of $Rf_{*}(K^{\bullet})$,  one natural idea of course would
be to take a projective resolution \var{\xout{of}}{}{typewriter deletion}{ps_134-8_034-v0005} $L^{A}.$ of $k_{A}$ (i.e. a coin\var{\over{e}{t}}{\hand{t}}{writing over}{ps_134-8_034-v0006}\var{\over{t}{e}}{\hand{e}}{writing over}{ps_134-8_034-v0007}grator for
A),  and write  tentatively
\[ (1) \qquad Rf_{*}(K^{\bullet}) \overset{?}{\hand{\simeq}} f_{*}(\underline{Hom}^{\bullet\bullet}_{k}(L^{A}.,K^{\bullet}))\ . \]
The (misleading) reflex inducing us to write down this formula, is that
this formula looks as if it were to boil down to the similar (correct)
formula for the maps $A_{\hand{/}b}\hand{\longrightarrow} e$ ,  when taking the localizations on B
\[ B_{\hand{/}b}\hand{\longrightarrow} B\ . \]
For this intuition to be correct,  \var{\xout{would}}{}{typewriter deletion}{ps_134-8_034-v0008} it should be true that the
restriction \var{\inserted[typed]{(or ``localization'')}}{(or ``localization'')}{typed insertion}{ps_134-8_034-v0009} of  $L^{A}.$  to $A_{\hand{/}b}$  (which is of course a resolution of the
constant sheaf k on $A_{\hand{/}b}$ ) is indeed a cointegrator on $A_{\hand{/}b}$ ,  namely
that it's components are still projective.  This however,  we suspect,
will hold true only under very special assumptions  -  as in general,
it is inverse image\var{\xout{s}}{}{typewriter deletion}{ps_134-8_034-v0010} by \uline{colocalization} (not by localization) which
takes projectives into projectives.  Thus I do'nt expect a relation (1)
to hold,  except under most exacting conditions on f and $K^{\bullet}$,  which I
did'nt try to pin down.  To take an example,  assume A has a final
object hence $k_{A}$ is projective and we may take $L^{A}.=k_{A}$ ,  then (1)
reads ( when $K^{\bullet}=K$ is reduced to degree zero)
\[ Rf_{*}(K)\hand{\simeq f_{*}}(K) \text{ i.e. } R^{i}f_{*}(K)=0 \text{ for } i \hand{>} 0, \]
which need not hold true even if \var{\xout{K is}}{}{typewriter deletion}{ps_134-8_034-v0011} K is constant ( = $k_{A}$ say).  For
instance,  we may start with an arbitrary object $A_{0}$ of (Cat) and add
a final object \var{\inserted[typed]{$e_1$}}{$e_1$}{typed insertion}{ps_134-8_034-v0012} to get A (intuitively,  it is the \hand{``}cone\hand{''} over $A_{0}$),  which
is mapped into \var{\inserted{the cone over e,}}{\hand{the cone over e,}}{insertion}{ps_134-8_034-v0013} $B = \hand{\Delta}_{1} = (0 \hand{\to} 1)$\hand{,} in the obvious way.  Then for any
sheaf \var{\over{K}{K}}{\hand{K}}{writing over}{ps_134-8_034-v0014} on A,  with restriction $K_{0}$ to $A_{0}$ ,  we get
\[ R^{i}f_{\text{\var{\deleted{\illegible{1}}}{}{deletion}{ps_134-8_034-v0015}}*}(K)_{0} \text{ (fiber at 0) } = H^{i}(A_{0},K_{0}) \qquad , \]
which needs not be 0 for $i \hand{>} 0$.

There is a big blunder at the end of section 101,  where the
formulae (12),(12'),(13) (p. 377)  (supposed to be ``essentially trivi\var{\deleted{al}}{}{deletion}{ps_134-8_034-v0016}\var{\inserted{-}}{\hand{-}}{insertion}{ps_134-8_034-v0017}%
al''),  are false for essentially the same kihd of reason.  \var{\over{\unsure{Th}}{Th}}{\hand{Th}}{writing over}{ps_134-8_034-v0018}e\var{\deleted{se}}{}{deletion}{ps_134-8_034-v0019} \var{\deleted{latter}}{}{deletion}{ps_134-8_034-v0020}
\var{\deleted{two}}{}{deletion}{ps_134-8_034-v0021} formulae r\var{\over{u}{a}}{\hand{a}}{writing over}{ps_134-8_034-v0022}n into the typewriter as a matter of course,  as \var{\deleted{they}}{}{deletion}{ps_134-8_034-v0023}

\margin{34}

% Visual collation: both supplied image halves inspected; see CHECK_DECISIONS_134_7_8.tsv.

%% END SOURCE ps_134-8_034
\ifnum\EditionMode=1
\psapparatusstart
\psentry{ps_134-8_034-v0001}{1}{typewriter deletion}{\xout{f: A B in (Cat) and}}{}{The typewriter strikes out this material; it leaves the corrected running text.}
\psentry{ps_134-8_034-v0002}{2}{typewriter deletion}{\xout{\unsure{L}}}{}{The typewriter strikes out this material; it leaves the corrected running text.}
\psentry{ps_134-8_034-v0003}{3}{writing over}{$\over{\unsure{F|L}}{L}$}{$\hand{L}$}{The superposed handwritten reading replaces the underlying reading.}
\psentry{ps_134-8_034-v0004}{4}{writing over}{\over{\unsure{e}}{e}}{\hand{e}}{The superposed handwritten reading replaces the underlying reading.}
\psentry{ps_134-8_034-v0005}{5}{typewriter deletion}{\xout{of}}{}{The typewriter strikes out this material; it leaves the corrected running text.}
\psentry{ps_134-8_034-v0006}{6}{writing over}{\over{e}{t}}{\hand{t}}{The superposed handwritten reading replaces the underlying reading.}
\psentry{ps_134-8_034-v0007}{7}{writing over}{\over{t}{e}}{\hand{e}}{The superposed handwritten reading replaces the underlying reading.}
\psentry{ps_134-8_034-v0008}{8}{typewriter deletion}{\xout{would}}{}{The typewriter strikes out this material; it leaves the corrected running text.}
\psentry{ps_134-8_034-v0009}{9}{typed insertion}{\inserted[typed]{(or ``localization'')}}{(or ``localization'')}{The typed insertion is incorporated at its indicated position; insertion brackets leave the running text.}
\psentry{ps_134-8_034-v0010}{10}{typewriter deletion}{\xout{s}}{}{The typewriter strikes out this material; it leaves the corrected running text.}
\psentry{ps_134-8_034-v0011}{11}{typewriter deletion}{\xout{K is}}{}{The typewriter strikes out this material; it leaves the corrected running text.}
\psentry{ps_134-8_034-v0012}{12}{typed insertion}{\inserted[typed]{$e_1$}}{$e_1$}{The typed insertion is incorporated at its indicated position; insertion brackets leave the running text.}
\psentry{ps_134-8_034-v0013}{13}{insertion}{\inserted{the cone over e,}}{\hand{the cone over e,}}{The handwritten insertion is incorporated at its indicated position; insertion brackets leave the running text.}
\psentry{ps_134-8_034-v0014}{14}{writing over}{\over{K}{K}}{\hand{K}}{The superposed handwritten reading replaces the underlying reading.}
\psentry{ps_134-8_034-v0015}{15}{deletion}{\deleted{\illegible{1}}}{}{The author cancels this material by hand; it leaves the corrected running text.}
\psentry{ps_134-8_034-v0016}{16}{deletion}{\deleted{al}}{}{The author cancels this material by hand; it leaves the corrected running text.}
\psentry{ps_134-8_034-v0017}{17}{insertion}{\inserted{-}}{\hand{-}}{The handwritten insertion is incorporated at its indicated position; insertion brackets leave the running text.}
\psentry{ps_134-8_034-v0018}{18}{writing over}{\over{\unsure{Th}}{Th}}{\hand{Th}}{The superposed handwritten reading replaces the underlying reading.}
\psentry{ps_134-8_034-v0019}{19}{deletion}{\deleted{se}}{}{The author cancels this material by hand; it leaves the corrected running text.}
\psentry{ps_134-8_034-v0020}{20}{deletion}{\deleted{latter}}{}{The author cancels this material by hand; it leaves the corrected running text.}
\psentry{ps_134-8_034-v0021}{21}{deletion}{\deleted{two}}{}{The author cancels this material by hand; it leaves the corrected running text.}
\psentry{ps_134-8_034-v0022}{22}{writing over}{\over{u}{a}}{\hand{a}}{The superposed handwritten reading replaces the underlying reading.}
\psentry{ps_134-8_034-v0023}{23}{deletion}{\deleted{they}}{}{The author cancels this material by hand; it leaves the corrected running text.}
\psapparatusend
\fi
\pssource{ps_134-8_035}{134-8/35}{https://grothendieck.umontpellier.fr/web/viewer.html?file=../134-8.pdf\#page=36}
\psnumber{ps_134-8_035-v0001}{1}
\psnumber{ps_134-8_035-v0002}{2}
\psnumber{ps_134-8_035-v0003}{3}
\psnumber{ps_134-8_035-v0004}{4}
\psnumber{ps_134-8_035-v0005}{5}
\psnumber{ps_134-8_035-v0006}{6}
\psnumber{ps_134-8_035-v0007}{7}
\psnumber{ps_134-8_035-v0008}{8}
\psnumber{ps_134-8_035-v0009}{9}
\psnumber{ps_134-8_035-v0010}{10}
\psnumber{ps_134-8_035-v0011}{11}
\psnumber{ps_134-8_035-v0012}{12}
\psnumber{ps_134-8_035-v0013}{13}
\psnumber{ps_134-8_035-v0014}{14}
\psnumber{ps_134-8_035-v0015}{15}
%% BEGIN SOURCE ps_134-8_035
% Pursuing Stacks (cotes 134-2 to 134-8): archive cote 134-8, site page 35 (PDF page 36), typescript page 589.
% typescript. Draft compared: ps7.tex:1641-1691.
% AI-transcribed by Claude Opus 5.5 (claude-opus-5-5, adaptive thinking, effort high), Message Batches API, batch ps4 (msgbatch_012Aoc1NuskATPsyBeKSYK7J), 2026-11-12; reviewed against both supplied scan halves; see CHECK_DECISIONS.tsv.
\scanpage{134-8/35}
\mspage{589}

they looked just the same as familiar \var{\xout{formu}}{}{typewriter deletion}{ps_134-8_035-v0001} one's from the standard
duality formulaire (with Lf$_{!}$, f$^*$ replaced by Rf$_{!}$, f$^{!}$)\hand{.} The first
one reads, in case as ab\hand{\textsubscript{\illegible{1}}}\var{\over{\unsure{xo}}{o}}{\hand{o}}{writing over}{ps_134-8_035-v0002}ve when A has a final object \hand{$\varepsilon$}
\[ \hand{\text{(2)}} \qquad Rf_*(f^*(K)) \overset{\hand{?}}{\hand{\simeq}} \underline{Hom}_k(f_{!}(k_A),K) \quad , \]
\var{\deleted{now in case formation of f$_{!}$ commutes \xout{to} to arbitrary base change}}{}{deletion}{ps_134-8_035-v0003}
\var{\deleted{(which is true when f is a fibering functor), the right hand side}}{}{deletion}{ps_134-8_035-v0004}
\var{\deleted{is isomorphic to \xout{Rx} f$_*$($\underline{\mathrm{Hom}}_k$(k$_A$,f$^*$(K)), and (2) may be rewritten}}{}{deletion}{ps_134-8_035-v0005}
\var{\deleted{\[ Rf_*(f^*(K)) \quad f_*(f^*(K)) \quad \text{i.e.} \quad R^if_*(f^*(K))= 0 \text{ for } i \ 0 \]}}{}{deletion}{ps_134-8_035-v0006}
\var{\deleted{which should be valid for any sheaf K on B, provided f is fibering and}}{}{deletion}{ps_134-8_035-v0007}
\var{\deleted{A has a final object.}}{}{deletion}{ps_134-8_035-v0008}
But we have
\[ f_{!}(k_A) = f_{!}(k^{(\hand{\varepsilon})}) = k^{(b)} \quad , \quad \text{where} \quad b = f_{\hand{!}}\hand{(\varepsilon) = f(\varepsilon)}\hand{,} \]
and if f takes final object into final object, we thus get
f$_{!}$(k$_A$) = k$_B$ , and the right hand side of (2) is just K , and hence
(2) implies \var{\deleted{\xout{R$^i$f$_*$(f$^*$(}}}{}{deletion}{ps_134-8_035-v0009}
\[ R^if_*(f^*(K)) = 0 \quad \text{for} \quad i \hand{>} 0 , \]
which however needs not be true even for K = k$_B$ , as we saw with the
previous example.

To come back to the general Lf$_{!}$, Rf$_*$ formalism, it seems it
can't be helped we'll have to do the usual silly thing and just resol-%
ve the argument involved, projectively in one case, injectively in
the other. Even in the ``absolute case'' when B = e , we could'nt help
either taking such resolutions, when it comes to working with \hand{the}
inter\var{\over{h}{n}}{\hand{n}}{writing over}{ps_134-8_035-v0010}al operations \var{\deleted{as}}{}{deletion}{ps_134-8_035-v0011} R$\underline{\mathrm{Hom}}^{\hand{\illegible{1}}}_k$(L.,K$^\bullet$) or F. $\var{\over{@}{\otimes}}{\hand{\otimes}}{writing over}{ps_134-8_035-v0012}^{\var{\over{L}{L}}{\hand{L}}{writing over}{ps_134-8_035-v0013}}_k$ L! , as we saw already
in section 108 (which should have tuned down a little my commitedness
to integrators, but it did'nt ! ).

Once this got clear, in order to turn the homology algebra
mill, all we still need is a handy criterion for existence of \var{\inserted[typed]{``\over{\unsure{en}}{en}ough''}}{``\hand{en}ough''}{typed insertion}{ps_134-8_035-v0014} projec-%
tives or injectives in categories of the type
\[ A^{\wedge}_{\hand{\mathcal{M}}} = \underline{Hom}(A^o, \hand{\mathcal{M}}) \quad , \]
with \hand{$\mathcal{M}$} an abelian category. In section 109 we got such a criterion
(prop.4 p.433) - the simplest one can imagine: it is \var{\xout{enough}}{}{typewriter deletion}{ps_134-8_035-v0015} suffi-%
cient that (besides the stability under direct or inverse limits,
needed anyhow in order for the functor f$_{!}$ resp. f$_*$ from A$^{\wedge}_{\hand{\mathcal{M}}}$ to B$^{\wedge}_{\hand{\mathcal{M}}}$

\margin{35}

% Visual collation: both supplied image halves inspected; see CHECK_DECISIONS_134_7_8.tsv.

%% END SOURCE ps_134-8_035
\ifnum\EditionMode=1
\psapparatusstart
\psentry{ps_134-8_035-v0001}{1}{typewriter deletion}{\xout{formu}}{}{The typewriter strikes out this material; it leaves the corrected running text.}
\psentry{ps_134-8_035-v0002}{2}{writing over}{\over{\unsure{xo}}{o}}{\hand{o}}{The superposed handwritten reading replaces the underlying reading.}
\psentry{ps_134-8_035-v0003}{3}{deletion}{\deleted{now in case formation of f$_{!}$ commutes \xout{to} to arbitrary base change}}{}{The author cancels this material by hand; it leaves the corrected running text.}
\psentry{ps_134-8_035-v0004}{4}{deletion}{\deleted{(which is true when f is a fibering functor), the right hand side}}{}{The author cancels this material by hand; it leaves the corrected running text.}
\psentry{ps_134-8_035-v0005}{5}{deletion}{\deleted{is isomorphic to \xout{Rx} f$_*$($\underline{\mathrm{Hom}}_k$(k$_A$,f$^*$(K)), and (2) may be rewritten}}{}{The author cancels this material by hand; it leaves the corrected running text.}
\psentry{ps_134-8_035-v0006}{6}{deletion}{\deleted{\[ Rf_*(f^*(K)) \quad f_*(f^*(K)) \quad \text{i.e.} \quad R^if_*(f^*(K))= 0 \text{ for } i \ 0 \]}}{}{The author cancels this material by hand; it leaves the corrected running text.}
\psentry{ps_134-8_035-v0007}{7}{deletion}{\deleted{which should be valid for any sheaf K on B, provided f is fibering and}}{}{The author cancels this material by hand; it leaves the corrected running text.}
\psentry{ps_134-8_035-v0008}{8}{deletion}{\deleted{A has a final object.}}{}{The author cancels this material by hand; it leaves the corrected running text.}
\psentry{ps_134-8_035-v0009}{9}{deletion}{\deleted{\xout{R$^i$f$_*$(f$^*$(}}}{}{The author cancels this material by hand; it leaves the corrected running text.}
\psentry{ps_134-8_035-v0010}{10}{writing over}{\over{h}{n}}{\hand{n}}{The superposed handwritten reading replaces the underlying reading.}
\psentry{ps_134-8_035-v0011}{11}{deletion}{\deleted{as}}{}{The author cancels this material by hand; it leaves the corrected running text.}
\psentry{ps_134-8_035-v0012}{12}{writing over}{$\over{@}{\otimes}$}{$\hand{\otimes}$}{The superposed handwritten reading replaces the underlying reading.}
\psentry{ps_134-8_035-v0013}{13}{writing over}{$\over{L}{L}$}{$\hand{L}$}{The superposed handwritten reading replaces the underlying reading.}
\psentry{ps_134-8_035-v0014}{14}{typed insertion}{\inserted[typed]{``\over{\unsure{en}}{en}ough''}}{``\hand{en}ough''}{The typed insertion is incorporated at its indicated position; insertion brackets leave the running text.}
\psentry{ps_134-8_035-v0015}{15}{typewriter deletion}{\xout{enough}}{}{The typewriter strikes out this material; it leaves the corrected running text.}
\psapparatusend
\fi
\pssource{ps_134-8_036}{134-8/36}{https://grothendieck.umontpellier.fr/web/viewer.html?file=../134-8.pdf\#page=37}
\psnumber{ps_134-8_036-v0001}{1}
\psnumber{ps_134-8_036-v0002}{2}
\psnumber{ps_134-8_036-v0003}{3}
\psnumber{ps_134-8_036-v0004}{4}
\psnumber{ps_134-8_036-v0005}{5}
\psnumber{ps_134-8_036-v0006}{6}
\psnumber{ps_134-8_036-v0007}{7}
\psnumber{ps_134-8_036-v0008}{8}
\psnumber{ps_134-8_036-v0009}{9}
\psnumber{ps_134-8_036-v0010}{10}
\psnumber{ps_134-8_036-v0011}{11}
\psnumber{ps_134-8_036-v0012}{12}
\psnumber{ps_134-8_036-v0013}{13}
\psnumber{ps_134-8_036-v0014}{14}
\psnumber{ps_134-8_036-v0015}{15}
\psnumber{ps_134-8_036-v0016}{16}
%% BEGIN SOURCE ps_134-8_036
% Pursuing Stacks (cotes 134-2 to 134-8): archive cote 134-8, site page 36 (PDF page 37), typescript page 590.
% typescript. Draft compared: ps7.tex:1679-1733.
% AI-transcribed by Claude Opus 5.5 (claude-opus-5-5, adaptive thinking, effort high), Message Batches API, batch ps4 (msgbatch_012Aoc1NuskATPsyBeKSYK7J), 2026-11-12; reviewed against both supplied scan halves; see CHECK_DECISIONS.tsv.
\scanpage{134-8/36}
\mspage{590}

to exist)  that in \hand{$\mathcal{M}$} there should be enough projectives resp. injecti-%
ves !  This gives what is needed \var{\inserted[typed]{surely}}{surely}{typed insertion}{ps_134-8_036-v0001} in order to grind through a mutually
dual treatment of relative homology Lf$_{!}$  and cohomology Rf$_{*}$ .  I have
the feeling that the little work ahead,for defining the basic operati-%
ons and working out the relevant formulaire \var{\inserted{(including ``cap products'')}}{\hand{(including ``cap products'')}}{insertion}{ps_134-8_036-v0002}, is a matter of mere
routine,  and I do'nt really expect any surprise may come up.  There-%
fore I do'nt feel like grinding through this,  and rather will feel
free to use whenever needed the most evi\var{\over[typed]{a}{d}}{d}{typewriter overwriting}{ps_134-8_036-v0003}ent formulae,  as the adjunc-%
tion formulae between the three functors
\[ \mathrm{Lf}_{!}\ ,\ f^{*}\ ,\ \mathrm{Rf}_{*}\ \ , \]
transitivity isomorphisms for a composition of maps,  possibly also
various projection formulae  -  trying however to be careful with
base change questions,  notably for Lf$_{!}$ ,  and not repeat the same
blunders !

In retrospect,  the main role \var{\inserted[typed]{for me}}{for me}{typed insertion}{ps_134-8_036-v0004} of the re\var{\over{f}{f}}{\hand{f}}{writing over}{ps_134-8_036-v0005}\var{\over{\unsure{f|l}}{l}}{\hand{l}}{writing over}{ps_134-8_036-v0006}ections of chapter V
on abelianization has been to become a little more familiar with the
\uline{homology} formalism \var{\inserted[typed]{in (Cat)}}{in (Cat)}{typed insertion}{ps_134-8_036-v0007}, namely essentially with the Lf$_{!}$ operation, which
has been more or less a white spot in my former experience,  centered
rather upon cohomology.  \var{\xout{This}}{}{typewriter deletion}{ps_134-8_036-v0008} An interesting,  still somewhat routine
byproduct of these ponderings has been the careful formulation of the
duality relationship \var{\xout{\illegible{1}}}{}{typewriter deletion}{ps_134-8_036-v0009} around the pairings between sheaves and
co-sheaves, namely the operations $*_{k}$ and $\hand{\circledast}_{k}$ .  Granting this,  the
game with integrators and cointegrators \var{\xout{is j\unsure{u}st the most}}{}{typewriter deletion}{ps_134-8_036-v0010} \var{\inserted[typed]{boil\over{\unsure{s}}{s} down to the}}{boil\hand{s} down to the}{typed insertion}{ps_134-8_036-v0011} standard
reflex,  of taking projective resolutions of the "unit" sheaf or
cosheaf,  k$_{A}$ or k$_{A^{o}}$  - which are the most obvious objects at hand of
all,  \var{\xout{among our}}{}{typewriter deletion}{ps_134-8_036-v0012} in our coefficient categories.

T$\overset{\illegible{1}}{h}$e  one idea which seems to me of wider scope and significance,
in this whole reflection on abelianization,  is the "further step in
linearization",  whereby the models in (Cat) are replaced by their
k-additive envelope,  endowed moreover with their natural \uline{diagonal}
\uline{map} \var{\inserted{(section 109)}}{\hand{(section 109)}}{insertion}{ps_134-8_036-v0013} .  The psychological effect of this discovery has been an immedia-%
te one  -  \var{\deleted{as}}{}{deletion}{ps_134-8_036-v0014} it triggered at once the reflection on schematization
of chapter VI.  This reflection,  apparently,  took me into a rather
different direction from \var{those\hand{$\vert$}"k-coalgebra}{those "k-coalgebra}{word separation}{ps_134-8_036-v0015} structures in (Cat)" ,
which had reawakened and made more acute the feeling that homotopy
types should make a sense "over any ground ring".  \var{\over{\unsure{c|C}}{C}}{\hand{C}}{writing over}{ps_134-8_036-v0016}oming back now to
the homology and cohomology formalism within (Cat),  it remains \hand{\unsure{a}}

\margin{36}

% Visual collation: both supplied image halves inspected; see CHECK_DECISIONS_134_7_8.tsv.

%% END SOURCE ps_134-8_036
\ifnum\EditionMode=1
\psapparatusstart
\psentry{ps_134-8_036-v0001}{1}{typed insertion}{\inserted[typed]{surely}}{surely}{The typed insertion is incorporated at its indicated position; insertion brackets leave the running text.}
\psentry{ps_134-8_036-v0002}{2}{insertion}{\inserted{(including ``cap products'')}}{\hand{(including ``cap products'')}}{The handwritten insertion is incorporated at its indicated position; insertion brackets leave the running text.}
\psentry{ps_134-8_036-v0003}{3}{typewriter overwriting}{\over[typed]{a}{d}}{d}{The superposed typed reading replaces the underlying reading.}
\psentry{ps_134-8_036-v0004}{4}{typed insertion}{\inserted[typed]{for me}}{for me}{The typed insertion is incorporated at its indicated position; insertion brackets leave the running text.}
\psentry{ps_134-8_036-v0005}{5}{writing over}{\over{f}{f}}{\hand{f}}{The superposed handwritten reading replaces the underlying reading.}
\psentry{ps_134-8_036-v0006}{6}{writing over}{\over{\unsure{f|l}}{l}}{\hand{l}}{The superposed handwritten reading replaces the underlying reading.}
\psentry{ps_134-8_036-v0007}{7}{typed insertion}{\inserted[typed]{in (Cat)}}{in (Cat)}{The typed insertion is incorporated at its indicated position; insertion brackets leave the running text.}
\psentry{ps_134-8_036-v0008}{8}{typewriter deletion}{\xout{This}}{}{The typewriter strikes out this material; it leaves the corrected running text.}
\psentry{ps_134-8_036-v0009}{9}{typewriter deletion}{\xout{\illegible{1}}}{}{The typewriter strikes out this material; it leaves the corrected running text.}
\psentry{ps_134-8_036-v0010}{10}{typewriter deletion}{\xout{is j\unsure{u}st the most}}{}{The typewriter strikes out this material; it leaves the corrected running text.}
\psentry{ps_134-8_036-v0011}{11}{typed insertion}{\inserted[typed]{boil\over{\unsure{s}}{s} down to the}}{boil\hand{s} down to the}{The typed insertion is incorporated at its indicated position; insertion brackets leave the running text.}
\psentry{ps_134-8_036-v0012}{12}{typewriter deletion}{\xout{among our}}{}{The typewriter strikes out this material; it leaves the corrected running text.}
\psentry{ps_134-8_036-v0013}{13}{insertion}{\inserted{(section 109)}}{\hand{(section 109)}}{The handwritten insertion is incorporated at its indicated position; insertion brackets leave the running text.}
\psentry{ps_134-8_036-v0014}{14}{deletion}{\deleted{as}}{}{The author cancels this material by hand; it leaves the corrected running text.}
\psentry{ps_134-8_036-v0015}{15}{word separation}{those\hand{$\vert$}"k-coalgebra}{those "k-coalgebra}{The pen vertical stroke is a separator cue;replace it by a word space in reading mode.}
\psentry{ps_134-8_036-v0016}{16}{writing over}{\over{\unsure{c|C}}{C}}{\hand{C}}{The superposed handwritten reading replaces the underlying reading.}
\psapparatusend
\fi
\pssource{ps_134-8_037}{134-8/37}{https://grothendieck.umontpellier.fr/web/viewer.html?file=../134-8.pdf\#page=38}
\psnumber{ps_134-8_037-v0001}{1}
\psnumber{ps_134-8_037-v0002}{2}
\psnumber{ps_134-8_037-v0003}{3}
\psnumber{ps_134-8_037-v0004}{4}
\psnumber{ps_134-8_037-v0005}{5}
\psnumber{ps_134-8_037-v0006}{6}
\psnumber{ps_134-8_037-v0007}{7}
\psnumber{ps_134-8_037-v0008}{8}
\psnumber{ps_134-8_037-v0009}{9}
\psnumber{ps_134-8_037-v0010}{10}
\psnumber{ps_134-8_037-v0011}{11}
\psnumber{ps_134-8_037-v0012}{12}
\psnumber{ps_134-8_037-v0013}{13}
\psnumber{ps_134-8_037-v0014}{14}
\psnumber{ps_134-8_037-v0015}{15}
\psnumber{ps_134-8_037-v0016}{16}
\psnumber{ps_134-8_037-v0017}{17}
\psnumber{ps_134-8_037-v0018}{18}
\psnumber{ps_134-8_037-v0019}{19}
\psnumber{ps_134-8_037-v0020}{20}
\psnumber{ps_134-8_037-v0021}{21}
\psnumber{ps_134-8_037-v0022}{22}
\psnumber{ps_134-8_037-v0023}{23}
\psnumber{ps_134-8_037-v0024}{24}
\psnumber{ps_134-8_037-v0025}{25}
\psnumber{ps_134-8_037-v0026}{26}
\psnumber{ps_134-8_037-v0027}{27}
\psnumber{ps_134-8_037-v0028}{28}
\psnumber{ps_134-8_037-v0029}{29}
\psnumber{ps_134-8_037-v0030}{30}
\psnumber{ps_134-8_037-v0031}{31}
%% BEGIN SOURCE ps_134-8_037
% Pursuing Stacks (cotes 134-2 to 134-8): archive cote 134-8, site page 37 (PDF page 38), typescript page 591.
% typescript. Draft compared: ps7.tex:1721-1780.
% AI-transcribed by Claude Opus 5.5 (claude-opus-5-5, adaptive thinking, effort high), Message Batches API, batch ps4 (msgbatch_012Aoc1NuskATPsyBeKSYK7J), 2026-11-12; reviewed against both supplied scan halves; see CHECK_DECISIONS.tsv.
\scanpage{134-8/37}
\mspage{591}

a very striking fact indeed for me, that as far as I can see at pre-%
sent, all basic operations in the \var{\inserted[typed]{(commutative)}}{(commutative)}{typed insertion}{ps_134-8_037-v0001} homology and cohomology formalism
in (Cat), and their basic properties and interrelations, should make
a sense for these linearized objects, which therefore may be viewed
as more perfect \var{\xout{models}}{}{typewriter deletion}{ps_134-8_037-v0002} \var{\inserted[typed]{carriers}}{carriers}{typed insertion}{ps_134-8_037-v0003} still than small categories for embodying the
relevant formalism. To what extent this feeling is indeed justified,
cannot of course be decided beforehand - only experience can tell.
For instance, does the ``six operations duality formalism'' contemplated
lately for (Cat) (sections 136,137), which has still to be worked
out, carry over to this wider, linearized \var{\xout{context}}{}{typewriter deletion}{ps_134-8_037-v0004} set-up ? This
will become clearer when the relevant notions in (Cat) are understood,
so that \unsure{if|it} will become a meaningful question whether \var{\xout{the}}{}{typewriter deletion}{ps_134-8_037-v0005} for a map
in (Cat), the property of being proper, or smooth, or an immersion,
may be read off in terms of properties of the corresponding map of
coalgebra structyres\var{\deleted{\unsure{,}}}{}{deletion}{ps_134-8_037-v0006} in (Cat). And what about subtler types of cohomo-%
logy operations, such as the Steenrod operations (which, I remember,
\var{\over{\unsure{can}}{may}}{\hand{may}}{writing over}{ps_134-8_037-v0007} be defined in the context of cohomology with coefficients in
general sheaves of $\underline{F}_p$-modules) ?

When concerned with homology and cohomology formalism in (Cat),
invo\unsure{l|f}\var{\over{av}{\unsure{v}}}{\hand{\unsure{v}}}{writing over}{ps_134-8_037-v0008}i\var{\over{n}{n}}{\hand{n}}{writing over}{ps_134-8_037-v0009}g general sheaves of coefficients, not merely locally con-%
stant ones, we are leaving, strictly speaking, the waters of a
re\var{\over{ff}{fl}}{\hand{fl}}{writing over}{ps_134-8_037-v0010}ection on``homotopy models''. The objects of (Cat) now are no lon-%
ger viewed as mere models \var{\xout{as}}{}{typewriter deletion}{ps_134-8_037-v0011} for homotopy types, but rather\unsure{,} \var{\xout{as}}{}{typewriter deletion}{ps_134-8_037-v0012} each
one as defining a topos, w\var{\over[typed]{i}{i}}{i}{typewriter overwriting}{ps_134-8_037-v0013}th the manifold richess it carries; \var{\inserted[typed]{a richess}}{a richess}{typed insertion}{ps_134-8_037-v0014} simi-%
lar to the \var{\xout{richess}}{}{typewriter deletion}{ps_134-8_037-v0015} \var{\inserted[typed]{one}}{one}{typed insertion}{ps_134-8_037-v0016} of a topological space, almost all of which is
being stripped off when looking at the mere homotopy type - including
\var{\inserted{\deleted{\unsure{finiteness}}}}{\hand{}}{insertion}{ps_134-8_037-v0017}
even such basic properties as dimension, compactness, smoothness, \var{\deleted{and}}{}{deletion}{ps_134-8_037-v0018}
\var{\inserted{cardinality and}}{\hand{cardinality and}}{insertion}{ps_134-8_037-v0019} the like. When passing from an object in (Cat) to the corresponding
\var{\xout{k-al}}{}{typewriter deletion}{ps_134-8_037-v0020} coalgebra structure in (Cat), much of this richess is preserved
- maybe everything, indeed, which \var{\xout{can}}{}{typewriter deletion}{ps_134-8_037-v0021} can be expressed in terms of
\uline{commutative} sheaves of coefficients. As for the realm of non commuta-%
tive cohomology formalism (which is supp\var{\over{s}{o}}{\hand{o}}{writing over}{ps_134-8_037-v0022}sed to be the main theme of
these notes, with overall title ``Pursuing Stacks'' !), it does'nt \var{\xout{seem}}{}{typewriter deletion}{ps_134-8_037-v0023}
seem, not at first sight at any rate, that much of this \var{\over{\unsure{can}}{could}}{\hand{could}}{writing over}{ps_134-8_037-v0024} be read
off the enveloping coalgebra \var{\inserted[typed]{$P=$}}{$P=$}{typed insertion}{ps_134-8_037-v0025} Add$_k$(A) of a given object A in (Cat) -
except\var{\xout{\unsure{,}}}{}{typewriter deletion}{ps_134-8_037-v0026} \var{\inserted{of course}}{\hand{of course}}{insertion}{ps_134-8_037-v0027} \var{\xout{maybe,}}{}{typewriter deletion}{ps_134-8_037-v0028} in the case \var{\deleted{of course}}{}{deletion}{ps_134-8_037-v0029} when A can be recovered in terms
of P, maybe as the category of \var{\inserted{``exponential''}}{\hand{``exponential''}}{insertion}{ps_134-8_037-v0030} pairs (x,u), where x is in P and
\[ u : \hand{\delta}(x) \hand{\simeq} x \var{\over{\unsure{@}}{\otimes}}{\hand{\otimes}}{writing over}{ps_134-8_037-v0031}_k x \]
an isomorphism.

\margin{37}

% Visual collation: both supplied image halves inspected; see CHECK_DECISIONS_134_7_8.tsv.

%% END SOURCE ps_134-8_037
\ifnum\EditionMode=1
\psapparatusstart
\psentry{ps_134-8_037-v0001}{1}{typed insertion}{\inserted[typed]{(commutative)}}{(commutative)}{The typed insertion is incorporated at its indicated position; insertion brackets leave the running text.}
\psentry{ps_134-8_037-v0002}{2}{typewriter deletion}{\xout{models}}{}{The typewriter strikes out this material; it leaves the corrected running text.}
\psentry{ps_134-8_037-v0003}{3}{typed insertion}{\inserted[typed]{carriers}}{carriers}{The typed insertion is incorporated at its indicated position; insertion brackets leave the running text.}
\psentry{ps_134-8_037-v0004}{4}{typewriter deletion}{\xout{context}}{}{The typewriter strikes out this material; it leaves the corrected running text.}
\psentry{ps_134-8_037-v0005}{5}{typewriter deletion}{\xout{the}}{}{The typewriter strikes out this material; it leaves the corrected running text.}
\psentry{ps_134-8_037-v0006}{6}{deletion}{\deleted{\unsure{,}}}{}{The author cancels this material by hand; it leaves the corrected running text.}
\psentry{ps_134-8_037-v0007}{7}{writing over}{\over{\unsure{can}}{may}}{\hand{may}}{The superposed handwritten reading replaces the underlying reading.}
\psentry{ps_134-8_037-v0008}{8}{writing over}{\over{av}{\unsure{v}}}{\hand{\unsure{v}}}{The superposed handwritten reading replaces the underlying reading.}
\psentry{ps_134-8_037-v0009}{9}{writing over}{\over{n}{n}}{\hand{n}}{The superposed handwritten reading replaces the underlying reading.}
\psentry{ps_134-8_037-v0010}{10}{writing over}{\over{ff}{fl}}{\hand{fl}}{The superposed handwritten reading replaces the underlying reading.}
\psentry{ps_134-8_037-v0011}{11}{typewriter deletion}{\xout{as}}{}{The typewriter strikes out this material; it leaves the corrected running text.}
\psentry{ps_134-8_037-v0012}{12}{typewriter deletion}{\xout{as}}{}{The typewriter strikes out this material; it leaves the corrected running text.}
\psentry{ps_134-8_037-v0013}{13}{typewriter overwriting}{\over[typed]{i}{i}}{i}{The superposed typed reading replaces the underlying reading.}
\psentry{ps_134-8_037-v0014}{14}{typed insertion}{\inserted[typed]{a richess}}{a richess}{The typed insertion is incorporated at its indicated position; insertion brackets leave the running text.}
\psentry{ps_134-8_037-v0015}{15}{typewriter deletion}{\xout{richess}}{}{The typewriter strikes out this material; it leaves the corrected running text.}
\psentry{ps_134-8_037-v0016}{16}{typed insertion}{\inserted[typed]{one}}{one}{The typed insertion is incorporated at its indicated position; insertion brackets leave the running text.}
\psentry{ps_134-8_037-v0017}{17}{insertion}{\inserted{\deleted{\unsure{finiteness}}}}{\hand{}}{The handwritten insertion is incorporated at its indicated position; insertion brackets leave the running text.}
\psentry{ps_134-8_037-v0018}{18}{deletion}{\deleted{and}}{}{The author cancels this material by hand; it leaves the corrected running text.}
\psentry{ps_134-8_037-v0019}{19}{insertion}{\inserted{cardinality and}}{\hand{cardinality and}}{The handwritten insertion is incorporated at its indicated position; insertion brackets leave the running text.}
\psentry{ps_134-8_037-v0020}{20}{typewriter deletion}{\xout{k-al}}{}{The typewriter strikes out this material; it leaves the corrected running text.}
\psentry{ps_134-8_037-v0021}{21}{typewriter deletion}{\xout{can}}{}{The typewriter strikes out this material; it leaves the corrected running text.}
\psentry{ps_134-8_037-v0022}{22}{writing over}{\over{s}{o}}{\hand{o}}{The superposed handwritten reading replaces the underlying reading.}
\psentry{ps_134-8_037-v0023}{23}{typewriter deletion}{\xout{seem}}{}{The typewriter strikes out this material; it leaves the corrected running text.}
\psentry{ps_134-8_037-v0024}{24}{writing over}{\over{\unsure{can}}{could}}{\hand{could}}{The superposed handwritten reading replaces the underlying reading.}
\psentry{ps_134-8_037-v0025}{25}{typed insertion}{\inserted[typed]{$P=$}}{$P=$}{The typed insertion is incorporated at its indicated position; insertion brackets leave the running text.}
\psentry{ps_134-8_037-v0026}{26}{typewriter deletion}{\xout{\unsure{,}}}{}{The typewriter strikes out this material; it leaves the corrected running text.}
\psentry{ps_134-8_037-v0027}{27}{insertion}{\inserted{of course}}{\hand{of course}}{The handwritten insertion is incorporated at its indicated position; insertion brackets leave the running text.}
\psentry{ps_134-8_037-v0028}{28}{typewriter deletion}{\xout{maybe,}}{}{The typewriter strikes out this material; it leaves the corrected running text.}
\psentry{ps_134-8_037-v0029}{29}{deletion}{\deleted{of course}}{}{The author cancels this material by hand; it leaves the corrected running text.}
\psentry{ps_134-8_037-v0030}{30}{insertion}{\inserted{``exponential''}}{\hand{``exponential''}}{The handwritten insertion is incorporated at its indicated position; insertion brackets leave the running text.}
\psentry{ps_134-8_037-v0031}{31}{writing over}{$\over{\unsure{@}}{\otimes}$}{$\hand{\otimes}$}{The superposed handwritten reading replaces the underlying reading.}
\psapparatusend
\fi
\pssource{ps_134-8_038}{134-8/38}{https://grothendieck.umontpellier.fr/web/viewer.html?file=../134-8.pdf\#page=39}
\psnumber{ps_134-8_038-v0001}{1}
\psnumber{ps_134-8_038-v0002}{2}
\psnumber{ps_134-8_038-v0003}{3}
\psnumber{ps_134-8_038-v0004}{4}
\psnumber{ps_134-8_038-v0005}{5}
\psnumber{ps_134-8_038-v0006}{6}
\psnumber{ps_134-8_038-v0007}{7}
\psnumber{ps_134-8_038-v0008}{8}
\psnumber{ps_134-8_038-v0009}{9}
\psnumber{ps_134-8_038-v0010}{10}
\psnumber{ps_134-8_038-v0011}{11}
\psnumber{ps_134-8_038-v0012}{12}
\psnumber{ps_134-8_038-v0013}{13}
\psnumber{ps_134-8_038-v0014}{14}
\psnumber{ps_134-8_038-v0015}{15}
\psnumber{ps_134-8_038-v0016}{16}
\psnumber{ps_134-8_038-v0017}{17}
\psnumber{ps_134-8_038-v0018}{18}
\psnumber{ps_134-8_038-v0019}{19}
\psnumber{ps_134-8_038-v0020}{20}
\psnumber{ps_134-8_038-v0021}{21}
\psnumber{ps_134-8_038-v0022}{22}
\psnumber{ps_134-8_038-v0023}{23}
\psnumber{ps_134-8_038-v0024}{24}
\psnumber{ps_134-8_038-v0025}{25}
\psnumber{ps_134-8_038-v0026}{26}
\psnumber{ps_134-8_038-v0027}{27}
\psnumber{ps_134-8_038-v0028}{28}
\psnumber{ps_134-8_038-v0029}{29}
\psnumber{ps_134-8_038-v0030}{30}
\psnumber{ps_134-8_038-v0031}{31}
\psnumber{ps_134-8_038-v0032}{32}
\psnumber{ps_134-8_038-v0033}{33}
%% BEGIN SOURCE ps_134-8_038
% Pursuing Stacks (cotes 134-2 to 134-8): archive cote 134-8, site page 38 (PDF page 39), typescript page 592.
% typescript. Draft compared: ps7.tex:1762-1823.
% AI-transcribed by Claude Opus 5.5 (claude-opus-5-5, adaptive thinking, effort high), Message Batches API, batch ps4 (msgbatch_012Aoc1NuskATPsyBeKSYK7J), 2026-11-12; reviewed against both supplied scan halves; see CHECK_DECISIONS.tsv.
\scanpage{134-8/38}
\mspage{592}

4.11.

140.\quad Life keeps pushing open the do\var{\over[typed]{r}{o}}{o}{typewriter overwriting}{ps_134-8_038-v0001}rs of that well-tempered hothouse
of my mathematical reflections, as a fresh wind and \var{\inserted[typed]{often an impetuous}}{often an impetuous}{typed insertion}{ps_134-8_038-v0002} \var{\xout{sometimes an stormy}}{}{typewriter deletion}{ps_134-8_038-v0003}
one, sweeping off the \var{\inserted[typed]{serene}}{serene}{typed insertion}{ps_134-8_038-v0004} quietness of \var{\xout{the world of}}{}{typewriter deletion}{ps_134-8_038-v0005} abstraction, \var{\inserted[typed]{- a breath}}{- a breath}{typed insertion}{ps_134-8_038-v0006} rich with
the manifold fragrance of the world we live in. This is the world of
conflict, weaving around each birth and each death and around the
lover\var{\over{s}{s}}{\hand{s}}{writing over}{ps_134-8_038-v0007}\hand{'}\var{\xout{\unsure{a}}}{}{typewriter deletion}{ps_134-8_038-v0008} play alike - it is the world we all have been born i n to
without our choosing. I used to see it as a stage - the stage set for
our acting. \var{O$\overset{\hand{\scriptscriptstyle\smile}}{\underset{\hand{\scriptscriptstyle\frown\,\frown}}{u}}$r}{Our}{nonlexical marks}{ps_134-8_038-v0009} freedom (rarely used indeed !) includes \var{\xout{changing}}{}{typewriter deletion}{ps_134-8_038-v0010}
choosing the role we are playing, possibly changing \var{\over[typed]{i}{r}}{r}{typewriter overwriting}{ps_134-8_038-v0011}\var{\over[typed]{t}{o}}{o}{typewriter overwriting}{ps_134-8_038-v0012}les - but not
choosing or changing the stage. It does'nt seem the \var{\xout{stage}}{}{typewriter deletion}{ps_134-8_038-v0013}\var{\inserted[typed]{stage}}{stage}{typed insertion}{ps_134-8_038-v0014} ever changed
during the history of mankind - only the décors kept changing. More
and more \var{\inserted[typed]{however}}{however}{typed insertion}{ps_134-8_038-v0015} over the last years, I h\var{\over{\unsure{a}}{a}}{\hand{a}}{writing over}{ps_134-8_038-v0016}\var{\over{\unsure{v}}{v}}{\hand{v}}{writing over}{ps_134-8_038-v0017}e been feeling th\var{\over[typed]{e}{i}}{i}{typewriter overwriting}{ps_134-8_038-v0018}\var{\over{s}{s}}{\hand{s}}{writing over}{ps_134-8_038-v0019} world I am
li\var{\over[typed]{g}{v}}{v}{typewriter overwriting}{ps_134-8_038-v0020}ing in, the world of conflict, somehow as a\var{\xout{n}}{}{typewriter deletion}{ps_134-8_038-v0021} \uline{meal} - a meal of
unexhaustible richess. Maybe the ultimate fruit and meaning of all my
acting on that stage, is that parts of that meal, of that richess,
be actually eaten, digested, assimilated - that they become part of
the flesh and bones of my own being. Maybe the ultimate purpose of
conflict, so deeply rooted in every human being, is to be the raw ma-%
terial, to be eaten and digested and changed into understanding about
conflict. Not a collective``understanding''(I doubt there is such a
thing !),written down in \var{text\deleted{\unsure{s}}\hand{\textsubscript{$\smile$}}books}{textbooks}{word joining}{ps_134-8_038-v0022} or sacred books or whatever, nor
even something expressed or expressible in wor\var{\over{d}{d}}{\hand{d}}{writing over}{ps_134-8_038-v0023}s necessarily - but
the kind of immediate knowledge rather the walker ha\var{\over[typed]{d}{s}}{s}{typewriter overwriting}{ps_134-8_038-v0024} about walking,
the swimmer about swimming, or the suckling \var{\over[typed]{a}{a}}{a}{typewriter overwriting}{ps_134-8_038-v0025}bout milk and mother's
breast\var{\over{s}{s}}{\hand{s}}{writing over}{ps_134-8_038-v0026}\var{\xout{\illegible{1}}}{}{typewriter deletion}{ps_134-8_038-v0027}. My business is to be a learner, not a teacher - na\var{\over[typed]{\unsure{e|o}}{m}}{m}{typewriter overwriting}{ps_134-8_038-v0028}ely
to allow this process to take place in my being, letting the world
of conflict, of suffering \var{\inserted[typed]{and of joy}}{and of joy}{typed insertion}{ps_134-8_038-v0029}, of violence and of tenderness, enter
and be digested \var{\inserted[typed]{and}}{and}{typed insertion}{ps_134-8_038-v0030} \var{\xout{to}}{}{typewriter deletion}{ps_134-8_038-v0031} become knowledge about myself.

I am not out though to write a \var{\deleted{diary}}{}{deletion}{ps_134-8_038-v0032} \var{\inserted{``journal intime''}}{\hand{``journal intime''}}{insertion}{ps_134-8_038-v0033} or meditation notes, so I
guess I better get back to the thread of mathematical reflection where
I left it, rather than write allusively about the events of these last
days, telling me about life and\unsure{:} about myself through one of my child-%
ren.

12.11.

Maybe at times I like to give the impression, to myself and hence
to others, that I am the easy learner of things of life, wholly rela-%

\margin{38}

% Visual collation: both supplied image halves inspected; see CHECK_DECISIONS_134_7_8.tsv.

%% END SOURCE ps_134-8_038
\ifnum\EditionMode=1
\psapparatusstart
\psentry{ps_134-8_038-v0001}{1}{typewriter overwriting}{\over[typed]{r}{o}}{o}{The superposed typed reading replaces the underlying reading.}
\psentry{ps_134-8_038-v0002}{2}{typed insertion}{\inserted[typed]{often an impetuous}}{often an impetuous}{The typed insertion is incorporated at its indicated position; insertion brackets leave the running text.}
\psentry{ps_134-8_038-v0003}{3}{typewriter deletion}{\xout{sometimes an stormy}}{}{The typewriter strikes out this material; it leaves the corrected running text.}
\psentry{ps_134-8_038-v0004}{4}{typed insertion}{\inserted[typed]{serene}}{serene}{The typed insertion is incorporated at its indicated position; insertion brackets leave the running text.}
\psentry{ps_134-8_038-v0005}{5}{typewriter deletion}{\xout{the world of}}{}{The typewriter strikes out this material; it leaves the corrected running text.}
\psentry{ps_134-8_038-v0006}{6}{typed insertion}{\inserted[typed]{- a breath}}{- a breath}{The typed insertion is incorporated at its indicated position; insertion brackets leave the running text.}
\psentry{ps_134-8_038-v0007}{7}{writing over}{\over{s}{s}}{\hand{s}}{The superposed handwritten reading replaces the underlying reading.}
\psentry{ps_134-8_038-v0008}{8}{typewriter deletion}{\xout{\unsure{a}}}{}{The typewriter strikes out this material; it leaves the corrected running text.}
\psentry{ps_134-8_038-v0009}{9}{nonlexical marks}{O$\overset{\hand{\scriptscriptstyle\smile}}{\underset{\hand{\scriptscriptstyle\frown\,\frown}}{u}}$r}{Our}{Omit the graphical cup and paired lower curves around u from the reading text while preserving them in the diplomatic source.}
\psentry{ps_134-8_038-v0010}{10}{typewriter deletion}{\xout{changing}}{}{The typewriter strikes out this material; it leaves the corrected running text.}
\psentry{ps_134-8_038-v0011}{11}{typewriter overwriting}{\over[typed]{i}{r}}{r}{The superposed typed reading replaces the underlying reading.}
\psentry{ps_134-8_038-v0012}{12}{typewriter overwriting}{\over[typed]{t}{o}}{o}{The superposed typed reading replaces the underlying reading.}
\psentry{ps_134-8_038-v0013}{13}{typewriter deletion}{\xout{stage}}{}{The typewriter strikes out this material; it leaves the corrected running text.}
\psentry{ps_134-8_038-v0014}{14}{typed insertion}{\inserted[typed]{stage}}{stage}{The typed insertion is incorporated at its indicated position; insertion brackets leave the running text.}
\psentry{ps_134-8_038-v0015}{15}{typed insertion}{\inserted[typed]{however}}{however}{The typed insertion is incorporated at its indicated position; insertion brackets leave the running text.}
\psentry{ps_134-8_038-v0016}{16}{writing over}{\over{\unsure{a}}{a}}{\hand{a}}{The superposed handwritten reading replaces the underlying reading.}
\psentry{ps_134-8_038-v0017}{17}{writing over}{\over{\unsure{v}}{v}}{\hand{v}}{The superposed handwritten reading replaces the underlying reading.}
\psentry{ps_134-8_038-v0018}{18}{typewriter overwriting}{\over[typed]{e}{i}}{i}{The superposed typed reading replaces the underlying reading.}
\psentry{ps_134-8_038-v0019}{19}{writing over}{\over{s}{s}}{\hand{s}}{The superposed handwritten reading replaces the underlying reading.}
\psentry{ps_134-8_038-v0020}{20}{typewriter overwriting}{\over[typed]{g}{v}}{v}{The superposed typed reading replaces the underlying reading.}
\psentry{ps_134-8_038-v0021}{21}{typewriter deletion}{\xout{n}}{}{The typewriter strikes out this material; it leaves the corrected running text.}
\psentry{ps_134-8_038-v0022}{22}{word joining}{text\deleted{\unsure{s}}\hand{\textsubscript{$\smile$}}books}{textbooks}{Apply the pen close-up cue beneath the cancelled extra letter between text and books.}
\psentry{ps_134-8_038-v0023}{23}{writing over}{\over{d}{d}}{\hand{d}}{The superposed handwritten reading replaces the underlying reading.}
\psentry{ps_134-8_038-v0024}{24}{typewriter overwriting}{\over[typed]{d}{s}}{s}{The superposed typed reading replaces the underlying reading.}
\psentry{ps_134-8_038-v0025}{25}{typewriter overwriting}{\over[typed]{a}{a}}{a}{The superposed typed reading replaces the underlying reading.}
\psentry{ps_134-8_038-v0026}{26}{writing over}{\over{s}{s}}{\hand{s}}{The superposed handwritten reading replaces the underlying reading.}
\psentry{ps_134-8_038-v0027}{27}{typewriter deletion}{\xout{\illegible{1}}}{}{The typewriter strikes out this material; it leaves the corrected running text.}
\psentry{ps_134-8_038-v0028}{28}{typewriter overwriting}{\over[typed]{\unsure{e|o}}{m}}{m}{The superposed typed reading replaces the underlying reading.}
\psentry{ps_134-8_038-v0029}{29}{typed insertion}{\inserted[typed]{and of joy}}{and of joy}{The typed insertion is incorporated at its indicated position; insertion brackets leave the running text.}
\psentry{ps_134-8_038-v0030}{30}{typed insertion}{\inserted[typed]{and}}{and}{The typed insertion is incorporated at its indicated position; insertion brackets leave the running text.}
\psentry{ps_134-8_038-v0031}{31}{typewriter deletion}{\xout{to}}{}{The typewriter strikes out this material; it leaves the corrected running text.}
\psentry{ps_134-8_038-v0032}{32}{deletion}{\deleted{diary}}{}{The author cancels this material by hand; it leaves the corrected running text.}
\psentry{ps_134-8_038-v0033}{33}{insertion}{\inserted{``journal intime''}}{\hand{``journal intime''}}{The handwritten insertion is incorporated at its indicated position; insertion brackets leave the running text.}
\psapparatusend
\fi
\pssource{ps_134-8_039}{134-8/39}{https://grothendieck.umontpellier.fr/web/viewer.html?file=../134-8.pdf\#page=40}
\psnumber{ps_134-8_039-v0001}{1}
\psnumber{ps_134-8_039-v0002}{2}
\psnumber{ps_134-8_039-v0003}{3}
\psnumber{ps_134-8_039-v0004}{4}
\psnumber{ps_134-8_039-v0005}{5}
\psnumber{ps_134-8_039-v0006}{6}
\psnumber{ps_134-8_039-v0007}{7}
\psnumber{ps_134-8_039-v0008}{8}
\psnumber{ps_134-8_039-v0009}{9}
\psnumber{ps_134-8_039-v0010}{10}
\psnumber{ps_134-8_039-v0011}{11}
\psnumber{ps_134-8_039-v0012}{12}
\psnumber{ps_134-8_039-v0013}{13}
\psnumber{ps_134-8_039-v0014}{14}
\psnumber{ps_134-8_039-v0015}{15}
\psnumber{ps_134-8_039-v0016}{16}
\psnumber{ps_134-8_039-v0017}{17}
\psnumber{ps_134-8_039-v0018}{18}
\psnumber{ps_134-8_039-v0019}{19}
\psnumber{ps_134-8_039-v0020}{20}
\psnumber{ps_134-8_039-v0021}{21}
\psnumber{ps_134-8_039-v0022}{22}
\psnumber{ps_134-8_039-v0023}{23}
\psnumber{ps_134-8_039-v0024}{24}
\psnumber{ps_134-8_039-v0025}{25}
\psnumber{ps_134-8_039-v0026}{26}
\psnumber{ps_134-8_039-v0027}{27}
\psnumber{ps_134-8_039-v0028}{28}
%% BEGIN SOURCE ps_134-8_039
% Pursuing Stacks (cotes 134-2 to 134-8): archive cote 134-8, site page 39 (PDF page 40), typescript page 593.
% typescript. Draft compared: ps7.tex:1806-1865.
% AI-transcribed by Claude Opus 5.5 (claude-opus-5-5, adaptive thinking, effort high), Message Batches API, batch ps4 (msgbatch_012Aoc1NuskATPsyBeKSYK7J), 2026-11-12; reviewed against both supplied scan halves; see CHECK_DECISIONS.tsv.
\scanpage{134-8/39}
\mspage{593}

xed, ``cool'' and all that - just keen \hand{for} learning, \hand{for} eating the
meal and wel\var{\over[typed]{w}{c}}{c}{typewriter overwriting}{ps_134-8_039-v0001}ome smilingly whatever comes with it's message, frustrati-%
on and sorrow and destructiveness and the softer dishes alike. This of
course is just humbug, an image d'Epinal which at whiles I'll kid \var{\xout{mysel}}{}{typewriter deletion}{ps_134-8_039-v0002}
myself into believing I am like. Truth is that I am a hard learner,
maybe as hard and reluctant as anyone. At any rate, the inbuilt mecha-%
nisms causing rejection of the dishes unpalatable to \var{\xout{our}}{}{typewriter deletion}{ps_134-8_039-v0003} \var{\inserted[typed]{my}}{my}{typed insertion}{ps_134-8_039-v0004} wholly conditi-%
oned, wholly ego-controlled taste, are as much present in me now as
they have ever been in my life, \var{\inserted{and}}{\hand{and}}{insertion}{ps_134-8_039-v0005} as much as in anyone else I know. This
interferes a lot with the learning - it causes a tremendous amount of
friction and energy dispersion (wholly unlike what happens when I am
learning mathematics say, namely discovering things about any kind of
substance \var{\xout{where}}{}{typewriter deletion}{ps_134-8_039-v0006} \var{\inserted[typed]{\deleted{\unsure{of}} which}}{ which}{typed insertion}{ps_134-8_039-v0007} my own person and ego is not part of ...). If \var{\over[typed]{h}{t}}{t}{typewriter overwriting}{ps_134-8_039-v0008}\var{\over[typed]{\unsure{t|h}}{h}}{h}{typewriter overwriting}{ps_134-8_039-v0009}ere
has been anything new appearing in my life, it is surely not the end
of this process of dispersion, or the end of inertia, closely related
to dispersion. It is something, rather, which causes learning to take
place all the same - be it the hard way, as it often happens, very
much like the troubled digestion of one who took \var{\xout{reluctantly}}{}{typewriter deletion}{ps_134-8_039-v0010} a substa\var{\over{\unsure{n}}{n}}{\hand{n}}{writing over}{ps_134-8_039-v0011}-%
tial meal \var{\over{\unsure{r}}{r}}{\hand{r}}{writing over}{ps_134-8_039-v0012}\var{\over{\unsure{e}}{e}}{\hand{e}}{writing over}{ps_134-8_039-v0013}luctantly or in a state of nervousness, of c\var{\over{r}{r}}{\hand{r}}{writing over}{ps_134-8_039-v0014}\var{\over{r}{i}}{\hand{i}}{writing over}{ps_134-8_039-v0015}spation.
Once \var{\xout{the}}{}{typewriter deletion}{ps_134-8_039-v0016} one is through with the digestion, though, the food \var{\xout{we}}{}{typewriter deletion}{ps_134-8_039-v0017} \var{\inserted[typed]{one}}{one}{typed insertion}{ps_134-8_039-v0018} ate
is transformed into flesh and muscles, blood and bones and the like,
just as good and g\var{\xout{\illegible{1}}}{}{typewriter deletion}{ps_134-8_039-v0019}enuine as if the meal had been taken rela\var{\over{x}{x}}{\hand{x}}{writing over}{ps_134-8_039-v0020}edly and
with eag\var{\xout{\illegible{1}}}{}{typewriter deletion}{ps_134-8_039-v0021}er \var{appetite,\hand{|}as}{appetite, as}{word separation}{ps_134-8_039-v0022} it deserved. What really counts for the process
of assimilation to be able to take place, is that in a certain sense
the food, palatable or not, be \uline{accepted} - not v\var{\over{o}{o}}{\hand{o}}{writing over}{ps_134-8_039-v0023}mited, or just
\var{\xout{keeping it}}{}{typewriter deletion}{ps_134-8_039-v0024} \var{\inserted[typed]{kept.}}{kept.}{typed insertion}{ps_134-8_039-v0025} in the bowels like a foreign body, sometimes for decennarie\hand{s.}
The remarkable fact I came to know through experience, is that even
after having \hand{been} kept thus inertly for a lifetime, a process of digesti-%
on and assimilation may still come into being and transform the \var{\inserted[typed]{obstrusive}}{obstrusive}{typed insertion}{ps_134-8_039-v0026} stuff
into li\var{\over{g}{v}}{\hand{v}}{writing over}{ps_134-8_039-v0027}ing substance.

During the last week I have been sick for a few days - a grippe
I might say, but surely a case of troubled digestion too. It seems
though I am through now - till the next case at any rate ! I daresay
life has been generous with me for these last three months, while I
hav\hand{'}nt even taken the trouble to stop with the mathe\var{\over[typed]{i}{m}}{m}{typewriter overwriting}{ps_134-8_039-v0028}atical nonsense
for any more than a week or two. This week too I still did some mathe-%
matical scratchwork, still along the lines of abelianization, which
keeps showing a lot richer than suspected.

\margin{39}

% Visual collation: both supplied image halves inspected; see CHECK_DECISIONS_134_7_8.tsv.

%% END SOURCE ps_134-8_039
\ifnum\EditionMode=1
\psapparatusstart
\psentry{ps_134-8_039-v0001}{1}{typewriter overwriting}{\over[typed]{w}{c}}{c}{The superposed typed reading replaces the underlying reading.}
\psentry{ps_134-8_039-v0002}{2}{typewriter deletion}{\xout{mysel}}{}{The typewriter strikes out this material; it leaves the corrected running text.}
\psentry{ps_134-8_039-v0003}{3}{typewriter deletion}{\xout{our}}{}{The typewriter strikes out this material; it leaves the corrected running text.}
\psentry{ps_134-8_039-v0004}{4}{typed insertion}{\inserted[typed]{my}}{my}{The typed insertion is incorporated at its indicated position; insertion brackets leave the running text.}
\psentry{ps_134-8_039-v0005}{5}{insertion}{\inserted{and}}{\hand{and}}{The handwritten insertion is incorporated at its indicated position; insertion brackets leave the running text.}
\psentry{ps_134-8_039-v0006}{6}{typewriter deletion}{\xout{where}}{}{The typewriter strikes out this material; it leaves the corrected running text.}
\psentry{ps_134-8_039-v0007}{7}{typed insertion}{\inserted[typed]{\deleted{\unsure{of}} which}}{ which}{The typed insertion is incorporated at its indicated position; insertion brackets leave the running text.}
\psentry{ps_134-8_039-v0008}{8}{typewriter overwriting}{\over[typed]{h}{t}}{t}{The superposed typed reading replaces the underlying reading.}
\psentry{ps_134-8_039-v0009}{9}{typewriter overwriting}{\over[typed]{\unsure{t|h}}{h}}{h}{The superposed typed reading replaces the underlying reading.}
\psentry{ps_134-8_039-v0010}{10}{typewriter deletion}{\xout{reluctantly}}{}{The typewriter strikes out this material; it leaves the corrected running text.}
\psentry{ps_134-8_039-v0011}{11}{writing over}{\over{\unsure{n}}{n}}{\hand{n}}{The superposed handwritten reading replaces the underlying reading.}
\psentry{ps_134-8_039-v0012}{12}{writing over}{\over{\unsure{r}}{r}}{\hand{r}}{The superposed handwritten reading replaces the underlying reading.}
\psentry{ps_134-8_039-v0013}{13}{writing over}{\over{\unsure{e}}{e}}{\hand{e}}{The superposed handwritten reading replaces the underlying reading.}
\psentry{ps_134-8_039-v0014}{14}{writing over}{\over{r}{r}}{\hand{r}}{The superposed handwritten reading replaces the underlying reading.}
\psentry{ps_134-8_039-v0015}{15}{writing over}{\over{r}{i}}{\hand{i}}{The superposed handwritten reading replaces the underlying reading.}
\psentry{ps_134-8_039-v0016}{16}{typewriter deletion}{\xout{the}}{}{The typewriter strikes out this material; it leaves the corrected running text.}
\psentry{ps_134-8_039-v0017}{17}{typewriter deletion}{\xout{we}}{}{The typewriter strikes out this material; it leaves the corrected running text.}
\psentry{ps_134-8_039-v0018}{18}{typed insertion}{\inserted[typed]{one}}{one}{The typed insertion is incorporated at its indicated position; insertion brackets leave the running text.}
\psentry{ps_134-8_039-v0019}{19}{typewriter deletion}{\xout{\illegible{1}}}{}{The typewriter strikes out this material; it leaves the corrected running text.}
\psentry{ps_134-8_039-v0020}{20}{writing over}{\over{x}{x}}{\hand{x}}{The superposed handwritten reading replaces the underlying reading.}
\psentry{ps_134-8_039-v0021}{21}{typewriter deletion}{\xout{\illegible{1}}}{}{The typewriter strikes out this material; it leaves the corrected running text.}
\psentry{ps_134-8_039-v0022}{22}{word separation}{appetite,\hand{|}as}{appetite, as}{Apply the pen word-separation cue after the comma.}
\psentry{ps_134-8_039-v0023}{23}{writing over}{\over{o}{o}}{\hand{o}}{The superposed handwritten reading replaces the underlying reading.}
\psentry{ps_134-8_039-v0024}{24}{typewriter deletion}{\xout{keeping it}}{}{The typewriter strikes out this material; it leaves the corrected running text.}
\psentry{ps_134-8_039-v0025}{25}{typed insertion}{\inserted[typed]{kept.}}{kept.}{The typed insertion is incorporated at its indicated position; insertion brackets leave the running text.}
\psentry{ps_134-8_039-v0026}{26}{typed insertion}{\inserted[typed]{obstrusive}}{obstrusive}{The typed insertion is incorporated at its indicated position; insertion brackets leave the running text.}
\psentry{ps_134-8_039-v0027}{27}{writing over}{\over{g}{v}}{\hand{v}}{The superposed handwritten reading replaces the underlying reading.}
\psentry{ps_134-8_039-v0028}{28}{typewriter overwriting}{\over[typed]{i}{m}}{m}{The superposed typed reading replaces the underlying reading.}
\psapparatusend
\fi

\end{document}
